% Draft v0.2 -- generated numbers come only from ../output/macros/results.tex and ../output/tables/.
% Build:  cd paper && latexmk -pdf paper.tex
% Review version (anonymous). For the arXiv version replace [review] by nothing: \usepackage{acl}.
\documentclass[11pt]{article}
\usepackage[review]{acl}
\usepackage{times}
\usepackage{latexsym}
\usepackage[T1]{fontenc}
\usepackage[utf8]{inputenc}
\usepackage{microtype}
\usepackage{booktabs}
\usepackage{amsmath}
\usepackage{xcolor}

% ---- generated results (never edit by hand) ----
% Auto-generated by code/05_export_macros.py -- do not edit by hand.
\newcommand{\RateHundredBpdQnaAcc}{1.00}
\newcommand{\RateHundredBpdQnaProp}{0.00}
\newcommand{\RateHundredBpdQnaDet}{0.85}
\newcommand{\RateHundredBpdQnaRec}{0.85}
\newcommand{\TokHundredBpdQna}{3.0}
\newcommand{\CountHundredBpdQnaTimeout}{0}
\newcommand{\CountHundredBpdQnaRetried}{0}
\newcommand{\RateHundredBpdCodeAcc}{0.97}
\newcommand{\RateHundredBpdCodeProp}{0.02}
\newcommand{\RateHundredBpdCodeDet}{0.63}
\newcommand{\RateHundredBpdCodeRec}{0.63}
\newcommand{\TokHundredBpdCode}{8.0}
\newcommand{\CountHundredBpdCodeTimeout}{1}
\newcommand{\CountHundredBpdCodeRetried}{0}
\newcommand{\RateHundredYonseiQnaAcc}{0.29}
\newcommand{\RateHundredYonseiQnaProp}{0.71}
\newcommand{\RateHundredYonseiQnaDet}{0.29}
\newcommand{\RateHundredYonseiQnaRec}{0.29}
\newcommand{\TokHundredYonseiQna}{1.4}
\newcommand{\CountHundredYonseiQnaTimeout}{0}
\newcommand{\CountHundredYonseiQnaRetried}{0}
\newcommand{\RateHundredYonseiCodeAcc}{1.00}
\newcommand{\RateHundredYonseiCodeProp}{0.00}
\newcommand{\RateHundredYonseiCodeDet}{0.23}
\newcommand{\RateHundredYonseiCodeRec}{0.23}
\newcommand{\TokHundredYonseiCode}{2.4}
\newcommand{\CountHundredYonseiCodeTimeout}{0}
\newcommand{\CountHundredYonseiCodeRetried}{0}
\newcommand{\RateHundredGenOneQnaAcc}{0.83}
\newcommand{\RateHundredGenOneQnaProp}{0.17}
\newcommand{\RateHundredGenOneQnaDet}{0.00}
\newcommand{\RateHundredGenOneQnaRec}{0.00}
\newcommand{\TokHundredGenOneQna}{1.2}
\newcommand{\CountHundredGenOneQnaTimeout}{0}
\newcommand{\CountHundredGenOneQnaRetried}{0}
\newcommand{\RateHundredGenOneCodeAcc}{0.88}
\newcommand{\RateHundredGenOneCodeProp}{0.10}
\newcommand{\RateHundredGenOneCodeDet}{0.00}
\newcommand{\RateHundredGenOneCodeRec}{0.00}
\newcommand{\TokHundredGenOneCode}{2.9}
\newcommand{\CountHundredGenOneCodeTimeout}{0}
\newcommand{\CountHundredGenOneCodeRetried}{0}
\newcommand{\RateBFourBpdQnaAcc}{1.00}
\newcommand{\RateBFourBpdQnaProp}{0.00}
\newcommand{\RateBFourBpdQnaDet}{0.95}
\newcommand{\RateBFourBpdQnaRec}{0.95}
\newcommand{\TokBFourBpdQna}{3.0}
\newcommand{\CountBFourBpdQnaTimeout}{0}
\newcommand{\CountBFourBpdQnaRetried}{0}
\newcommand{\RateBFourBpdCodeAcc}{1.00}
\newcommand{\RateBFourBpdCodeProp}{0.00}
\newcommand{\RateBFourBpdCodeDet}{0.40}
\newcommand{\RateBFourBpdCodeRec}{0.40}
\newcommand{\TokBFourBpdCode}{8.7}
\newcommand{\CountBFourBpdCodeTimeout}{0}
\newcommand{\CountBFourBpdCodeRetried}{0}
\newcommand{\RateBFourFlatThreeQnaAcc}{1.00}
\newcommand{\RateBFourFlatThreeQnaProp}{0.00}
\newcommand{\RateBFourFlatThreeQnaDet}{1.00}
\newcommand{\RateBFourFlatThreeQnaRec}{1.00}
\newcommand{\TokBFourFlatThreeQna}{3.0}
\newcommand{\CountBFourFlatThreeQnaTimeout}{0}
\newcommand{\CountBFourFlatThreeQnaRetried}{0}
\newcommand{\RateBFourFlatThreeCodeAcc}{0.95}
\newcommand{\RateBFourFlatThreeCodeProp}{0.05}
\newcommand{\RateBFourFlatThreeCodeDet}{0.00}
\newcommand{\RateBFourFlatThreeCodeRec}{0.00}
\newcommand{\TokBFourFlatThreeCode}{8.7}
\newcommand{\CountBFourFlatThreeCodeTimeout}{0}
\newcommand{\CountBFourFlatThreeCodeRetried}{0}
\newcommand{\RateBFourGenOneQnaAcc}{0.90}
\newcommand{\RateBFourGenOneQnaProp}{0.10}
\newcommand{\RateBFourGenOneQnaDet}{0.00}
\newcommand{\RateBFourGenOneQnaRec}{0.00}
\newcommand{\TokBFourGenOneQna}{1.2}
\newcommand{\CountBFourGenOneQnaTimeout}{0}
\newcommand{\CountBFourGenOneQnaRetried}{0}
\newcommand{\RateBFourGenOneCodeAcc}{0.70}
\newcommand{\RateBFourGenOneCodeProp}{0.25}
\newcommand{\RateBFourGenOneCodeDet}{0.00}
\newcommand{\RateBFourGenOneCodeRec}{0.00}
\newcommand{\TokBFourGenOneCode}{3.1}
\newcommand{\CountBFourGenOneCodeTimeout}{0}
\newcommand{\CountBFourGenOneCodeRetried}{0}
\newcommand{\RateBFourGenThreeQnaAcc}{1.00}
\newcommand{\RateBFourGenThreeQnaProp}{0.00}
\newcommand{\RateBFourGenThreeQnaDet}{0.15}
\newcommand{\RateBFourGenThreeQnaRec}{0.15}
\newcommand{\TokBFourGenThreeQna}{3.6}
\newcommand{\CountBFourGenThreeQnaTimeout}{0}
\newcommand{\CountBFourGenThreeQnaRetried}{0}
\newcommand{\RateBFourGenThreeCodeAcc}{0.95}
\newcommand{\RateBFourGenThreeCodeProp}{0.00}
\newcommand{\RateBFourGenThreeCodeDet}{0.00}
\newcommand{\RateBFourGenThreeCodeRec}{0.00}
\newcommand{\TokBFourGenThreeCode}{9.1}
\newcommand{\CountBFourGenThreeCodeTimeout}{0}
\newcommand{\CountBFourGenThreeCodeRetried}{0}
\newcommand{\RateBFiveBpdQnaAcc}{1.00}
\newcommand{\RateBFiveBpdQnaProp}{0.00}
\newcommand{\RateBFiveBpdQnaDet}{0.88}
\newcommand{\RateBFiveBpdQnaRec}{0.88}
\newcommand{\TokBFiveBpdQna}{3.0}
\newcommand{\CountBFiveBpdQnaTimeout}{0}
\newcommand{\CountBFiveBpdQnaRetried}{0}
\newcommand{\RateBFiveBpdCodeAcc}{0.99}
\newcommand{\RateBFiveBpdCodeProp}{0.01}
\newcommand{\RateBFiveBpdCodeDet}{0.63}
\newcommand{\RateBFiveBpdCodeRec}{0.63}
\newcommand{\TokBFiveBpdCode}{8.1}
\newcommand{\CountBFiveBpdCodeTimeout}{0}
\newcommand{\CountBFiveBpdCodeRetried}{0}
\newcommand{\RateBFiveFlatThreeQnaAcc}{1.00}
\newcommand{\RateBFiveFlatThreeQnaProp}{0.00}
\newcommand{\RateBFiveFlatThreeQnaDet}{1.00}
\newcommand{\RateBFiveFlatThreeQnaRec}{1.00}
\newcommand{\TokBFiveFlatThreeQna}{3.0}
\newcommand{\CountBFiveFlatThreeQnaTimeout}{0}
\newcommand{\CountBFiveFlatThreeQnaRetried}{0}
\newcommand{\RateBFiveFlatThreeCodeAcc}{0.91}
\newcommand{\RateBFiveFlatThreeCodeProp}{0.09}
\newcommand{\RateBFiveFlatThreeCodeDet}{0.00}
\newcommand{\RateBFiveFlatThreeCodeRec}{0.00}
\newcommand{\TokBFiveFlatThreeCode}{8.1}
\newcommand{\CountBFiveFlatThreeCodeTimeout}{0}
\newcommand{\CountBFiveFlatThreeCodeRetried}{0}
\newcommand{\RateBFiveGenOneQnaAcc}{0.99}
\newcommand{\RateBFiveGenOneQnaProp}{0.01}
\newcommand{\RateBFiveGenOneQnaDet}{0.00}
\newcommand{\RateBFiveGenOneQnaRec}{0.00}
\newcommand{\TokBFiveGenOneQna}{1.3}
\newcommand{\CountBFiveGenOneQnaTimeout}{0}
\newcommand{\CountBFiveGenOneQnaRetried}{0}
\newcommand{\RateBFiveGenOneCodeAcc}{0.84}
\newcommand{\RateBFiveGenOneCodeProp}{0.11}
\newcommand{\RateBFiveGenOneCodeDet}{0.00}
\newcommand{\RateBFiveGenOneCodeRec}{0.00}
\newcommand{\TokBFiveGenOneCode}{3.0}
\newcommand{\CountBFiveGenOneCodeTimeout}{0}
\newcommand{\CountBFiveGenOneCodeRetried}{0}
\newcommand{\RateBFiveGenThreeQnaAcc}{1.00}
\newcommand{\RateBFiveGenThreeQnaProp}{0.00}
\newcommand{\RateBFiveGenThreeQnaDet}{0.00}
\newcommand{\RateBFiveGenThreeQnaRec}{0.00}
\newcommand{\TokBFiveGenThreeQna}{3.6}
\newcommand{\CountBFiveGenThreeQnaTimeout}{0}
\newcommand{\CountBFiveGenThreeQnaRetried}{0}
\newcommand{\RateBFiveGenThreeCodeAcc}{1.00}
\newcommand{\RateBFiveGenThreeCodeProp}{0.00}
\newcommand{\RateBFiveGenThreeCodeDet}{0.00}
\newcommand{\RateBFiveGenThreeCodeRec}{0.00}
\newcommand{\TokBFiveGenThreeCode}{8.4}
\newcommand{\CountBFiveGenThreeCodeTimeout}{0}
\newcommand{\CountBFiveGenThreeCodeRetried}{0}
\newcommand{\RateBFiveJudgeOnlyQnaAcc}{1.00}
\newcommand{\RateBFiveJudgeOnlyQnaProp}{0.00}
\newcommand{\RateBFiveJudgeOnlyQnaDet}{0.84}
\newcommand{\RateBFiveJudgeOnlyQnaRec}{0.84}
\newcommand{\TokBFiveJudgeOnlyQna}{3.0}
\newcommand{\CountBFiveJudgeOnlyQnaTimeout}{0}
\newcommand{\CountBFiveJudgeOnlyQnaRetried}{0}
\newcommand{\RateBFiveJudgeOnlyCodeAcc}{1.00}
\newcommand{\RateBFiveJudgeOnlyCodeProp}{0.00}
\newcommand{\RateBFiveJudgeOnlyCodeDet}{0.63}
\newcommand{\RateBFiveJudgeOnlyCodeRec}{0.63}
\newcommand{\TokBFiveJudgeOnlyCode}{8.2}
\newcommand{\CountBFiveJudgeOnlyCodeTimeout}{0}
\newcommand{\CountBFiveJudgeOnlyCodeRetried}{0}
\newcommand{\RateBSixBpdQnaAcc}{1.00}
\newcommand{\RateBSixBpdQnaProp}{0.00}
\newcommand{\RateBSixBpdQnaDet}{0.87}
\newcommand{\RateBSixBpdQnaRec}{0.87}
\newcommand{\RateBSixBpdQnaEmit}{0.98}
\newcommand{\RateBSixBpdQnaDetEmit}{0.89}
\newcommand{\TokBSixBpdQna}{3.0}
\newcommand{\CountBSixBpdQnaTimeout}{0}
\newcommand{\CountBSixBpdQnaRetried}{0}
\newcommand{\CountBSixBpdQnaEmit}{98}
\newcommand{\RateBSixBpdCodeAcc}{0.88}
\newcommand{\RateBSixBpdCodeProp}{0.01}
\newcommand{\RateBSixBpdCodeDet}{0.61}
\newcommand{\RateBSixBpdCodeRec}{0.61}
\newcommand{\RateBSixBpdCodeEmit}{0.65}
\newcommand{\RateBSixBpdCodeDetEmit}{0.90}
\newcommand{\TokBSixBpdCode}{7.2}
\newcommand{\CountBSixBpdCodeTimeout}{11}
\newcommand{\CountBSixBpdCodeRetried}{0}
\newcommand{\CountBSixBpdCodeEmit}{58}
\newcommand{\RateBSixFlatThreeQnaAcc}{1.00}
\newcommand{\RateBSixFlatThreeQnaProp}{0.00}
\newcommand{\RateBSixFlatThreeQnaDet}{0.99}
\newcommand{\RateBSixFlatThreeQnaRec}{0.99}
\newcommand{\RateBSixFlatThreeQnaEmit}{0.99}
\newcommand{\RateBSixFlatThreeQnaDetEmit}{1.00}
\newcommand{\TokBSixFlatThreeQna}{3.1}
\newcommand{\CountBSixFlatThreeQnaTimeout}{0}
\newcommand{\CountBSixFlatThreeQnaRetried}{0}
\newcommand{\CountBSixFlatThreeQnaEmit}{99}
\newcommand{\RateBSixFlatThreeCodeAcc}{0.92}
\newcommand{\RateBSixFlatThreeCodeProp}{0.08}
\newcommand{\RateBSixFlatThreeCodeDet}{0.01}
\newcommand{\RateBSixFlatThreeCodeRec}{0.01}
\newcommand{\RateBSixFlatThreeCodeEmit}{0.56}
\newcommand{\RateBSixFlatThreeCodeDetEmit}{0.02}
\newcommand{\TokBSixFlatThreeCode}{8.1}
\newcommand{\CountBSixFlatThreeCodeTimeout}{0}
\newcommand{\CountBSixFlatThreeCodeRetried}{0}
\newcommand{\CountBSixFlatThreeCodeEmit}{56}
\newcommand{\RateBSixGenOneQnaAcc}{0.00}
\newcommand{\RateBSixGenOneQnaProp}{1.00}
\newcommand{\RateBSixGenOneQnaDet}{0.00}
\newcommand{\RateBSixGenOneQnaRec}{0.00}
\newcommand{\RateBSixGenOneQnaEmit}{1.00}
\newcommand{\RateBSixGenOneQnaDetEmit}{0.00}
\newcommand{\TokBSixGenOneQna}{0.6}
\newcommand{\CountBSixGenOneQnaTimeout}{0}
\newcommand{\CountBSixGenOneQnaRetried}{0}
\newcommand{\CountBSixGenOneQnaEmit}{100}
\newcommand{\RateBSixGenOneCodeAcc}{0.54}
\newcommand{\RateBSixGenOneCodeProp}{0.41}
\newcommand{\RateBSixGenOneCodeDet}{0.00}
\newcommand{\RateBSixGenOneCodeRec}{0.00}
\newcommand{\RateBSixGenOneCodeEmit}{0.50}
\newcommand{\RateBSixGenOneCodeDetEmit}{0.00}
\newcommand{\TokBSixGenOneCode}{3.0}
\newcommand{\CountBSixGenOneCodeTimeout}{0}
\newcommand{\CountBSixGenOneCodeRetried}{0}
\newcommand{\CountBSixGenOneCodeEmit}{50}
\newcommand{\RateBSixGenThreeQnaAcc}{1.00}
\newcommand{\RateBSixGenThreeQnaProp}{0.00}
\newcommand{\RateBSixGenThreeQnaDet}{1.00}
\newcommand{\RateBSixGenThreeQnaRec}{1.00}
\newcommand{\RateBSixGenThreeQnaEmit}{1.00}
\newcommand{\RateBSixGenThreeQnaDetEmit}{1.00}
\newcommand{\TokBSixGenThreeQna}{3.0}
\newcommand{\CountBSixGenThreeQnaTimeout}{0}
\newcommand{\CountBSixGenThreeQnaRetried}{0}
\newcommand{\CountBSixGenThreeQnaEmit}{100}
\newcommand{\RateBSixGenThreeCodeAcc}{0.97}
\newcommand{\RateBSixGenThreeCodeProp}{0.02}
\newcommand{\RateBSixGenThreeCodeDet}{0.01}
\newcommand{\RateBSixGenThreeCodeRec}{0.01}
\newcommand{\RateBSixGenThreeCodeEmit}{0.49}
\newcommand{\RateBSixGenThreeCodeDetEmit}{0.02}
\newcommand{\TokBSixGenThreeCode}{8.2}
\newcommand{\CountBSixGenThreeCodeTimeout}{0}
\newcommand{\CountBSixGenThreeCodeRetried}{0}
\newcommand{\CountBSixGenThreeCodeEmit}{49}
\newcommand{\RateBSixJudgeOnlyQnaAcc}{1.00}
\newcommand{\RateBSixJudgeOnlyQnaProp}{0.00}
\newcommand{\RateBSixJudgeOnlyQnaDet}{0.85}
\newcommand{\RateBSixJudgeOnlyQnaRec}{0.85}
\newcommand{\RateBSixJudgeOnlyQnaEmit}{0.99}
\newcommand{\RateBSixJudgeOnlyQnaDetEmit}{0.86}
\newcommand{\TokBSixJudgeOnlyQna}{3.0}
\newcommand{\CountBSixJudgeOnlyQnaTimeout}{0}
\newcommand{\CountBSixJudgeOnlyQnaRetried}{0}
\newcommand{\CountBSixJudgeOnlyQnaEmit}{99}
\newcommand{\RateBSixJudgeOnlyCodeAcc}{0.99}
\newcommand{\RateBSixJudgeOnlyCodeProp}{0.01}
\newcommand{\RateBSixJudgeOnlyCodeDet}{0.64}
\newcommand{\RateBSixJudgeOnlyCodeRec}{0.64}
\newcommand{\RateBSixJudgeOnlyCodeEmit}{0.58}
\newcommand{\RateBSixJudgeOnlyCodeDetEmit}{0.90}
\newcommand{\TokBSixJudgeOnlyCode}{8.1}
\newcommand{\CountBSixJudgeOnlyCodeTimeout}{0}
\newcommand{\CountBSixJudgeOnlyCodeRetried}{0}
\newcommand{\CountBSixJudgeOnlyCodeEmit}{58}
\newcommand{\RateBSixRBpdQnaAcc}{1.00}
\newcommand{\RateBSixRBpdQnaProp}{0.00}
\newcommand{\RateBSixRBpdQnaDet}{0.87}
\newcommand{\RateBSixRBpdQnaRec}{0.87}
\newcommand{\RateBSixRBpdQnaEmit}{0.98}
\newcommand{\RateBSixRBpdQnaDetEmit}{0.89}
\newcommand{\TokBSixRBpdQna}{3.0}
\newcommand{\CountBSixRBpdQnaTimeout}{0}
\newcommand{\CountBSixRBpdQnaRetried}{0}
\newcommand{\CountBSixRBpdQnaEmit}{98}
\newcommand{\RateBSixRBpdCodeAcc}{0.96}
\newcommand{\RateBSixRBpdCodeProp}{0.01}
\newcommand{\RateBSixRBpdCodeDet}{0.63}
\newcommand{\RateBSixRBpdCodeRec}{0.63}
\newcommand{\RateBSixRBpdCodeEmit}{0.61}
\newcommand{\RateBSixRBpdCodeDetEmit}{0.90}
\newcommand{\TokBSixRBpdCode}{7.8}
\newcommand{\CountBSixRBpdCodeTimeout}{3}
\newcommand{\CountBSixRBpdCodeRetried}{11}
\newcommand{\CountBSixRBpdCodeEmit}{59}
\newcommand{\RateBSixRFlatThreeQnaAcc}{1.00}
\newcommand{\RateBSixRFlatThreeQnaProp}{0.00}
\newcommand{\RateBSixRFlatThreeQnaDet}{0.99}
\newcommand{\RateBSixRFlatThreeQnaRec}{0.99}
\newcommand{\RateBSixRFlatThreeQnaEmit}{0.99}
\newcommand{\RateBSixRFlatThreeQnaDetEmit}{1.00}
\newcommand{\TokBSixRFlatThreeQna}{3.1}
\newcommand{\CountBSixRFlatThreeQnaTimeout}{0}
\newcommand{\CountBSixRFlatThreeQnaRetried}{0}
\newcommand{\CountBSixRFlatThreeQnaEmit}{99}
\newcommand{\RateBSixRFlatThreeCodeAcc}{0.92}
\newcommand{\RateBSixRFlatThreeCodeProp}{0.08}
\newcommand{\RateBSixRFlatThreeCodeDet}{0.01}
\newcommand{\RateBSixRFlatThreeCodeRec}{0.01}
\newcommand{\RateBSixRFlatThreeCodeEmit}{0.56}
\newcommand{\RateBSixRFlatThreeCodeDetEmit}{0.02}
\newcommand{\TokBSixRFlatThreeCode}{8.1}
\newcommand{\CountBSixRFlatThreeCodeTimeout}{0}
\newcommand{\CountBSixRFlatThreeCodeRetried}{0}
\newcommand{\CountBSixRFlatThreeCodeEmit}{56}
\newcommand{\RateBSixRGenOneQnaAcc}{0.00}
\newcommand{\RateBSixRGenOneQnaProp}{1.00}
\newcommand{\RateBSixRGenOneQnaDet}{0.00}
\newcommand{\RateBSixRGenOneQnaRec}{0.00}
\newcommand{\RateBSixRGenOneQnaEmit}{1.00}
\newcommand{\RateBSixRGenOneQnaDetEmit}{0.00}
\newcommand{\TokBSixRGenOneQna}{0.6}
\newcommand{\CountBSixRGenOneQnaTimeout}{0}
\newcommand{\CountBSixRGenOneQnaRetried}{0}
\newcommand{\CountBSixRGenOneQnaEmit}{100}
\newcommand{\RateBSixRGenOneCodeAcc}{0.54}
\newcommand{\RateBSixRGenOneCodeProp}{0.41}
\newcommand{\RateBSixRGenOneCodeDet}{0.00}
\newcommand{\RateBSixRGenOneCodeRec}{0.00}
\newcommand{\RateBSixRGenOneCodeEmit}{0.50}
\newcommand{\RateBSixRGenOneCodeDetEmit}{0.00}
\newcommand{\TokBSixRGenOneCode}{3.0}
\newcommand{\CountBSixRGenOneCodeTimeout}{0}
\newcommand{\CountBSixRGenOneCodeRetried}{0}
\newcommand{\CountBSixRGenOneCodeEmit}{50}
\newcommand{\RateBSixRGenThreeQnaAcc}{1.00}
\newcommand{\RateBSixRGenThreeQnaProp}{0.00}
\newcommand{\RateBSixRGenThreeQnaDet}{1.00}
\newcommand{\RateBSixRGenThreeQnaRec}{1.00}
\newcommand{\RateBSixRGenThreeQnaEmit}{1.00}
\newcommand{\RateBSixRGenThreeQnaDetEmit}{1.00}
\newcommand{\TokBSixRGenThreeQna}{3.0}
\newcommand{\CountBSixRGenThreeQnaTimeout}{0}
\newcommand{\CountBSixRGenThreeQnaRetried}{0}
\newcommand{\CountBSixRGenThreeQnaEmit}{100}
\newcommand{\RateBSixRGenThreeCodeAcc}{0.97}
\newcommand{\RateBSixRGenThreeCodeProp}{0.02}
\newcommand{\RateBSixRGenThreeCodeDet}{0.01}
\newcommand{\RateBSixRGenThreeCodeRec}{0.01}
\newcommand{\RateBSixRGenThreeCodeEmit}{0.49}
\newcommand{\RateBSixRGenThreeCodeDetEmit}{0.02}
\newcommand{\TokBSixRGenThreeCode}{8.2}
\newcommand{\CountBSixRGenThreeCodeTimeout}{0}
\newcommand{\CountBSixRGenThreeCodeRetried}{0}
\newcommand{\CountBSixRGenThreeCodeEmit}{49}
\newcommand{\RateBSixRJudgeOnlyQnaAcc}{1.00}
\newcommand{\RateBSixRJudgeOnlyQnaProp}{0.00}
\newcommand{\RateBSixRJudgeOnlyQnaDet}{0.85}
\newcommand{\RateBSixRJudgeOnlyQnaRec}{0.85}
\newcommand{\RateBSixRJudgeOnlyQnaEmit}{0.99}
\newcommand{\RateBSixRJudgeOnlyQnaDetEmit}{0.86}
\newcommand{\TokBSixRJudgeOnlyQna}{3.0}
\newcommand{\CountBSixRJudgeOnlyQnaTimeout}{0}
\newcommand{\CountBSixRJudgeOnlyQnaRetried}{0}
\newcommand{\CountBSixRJudgeOnlyQnaEmit}{99}
\newcommand{\RateBSixRJudgeOnlyCodeAcc}{0.99}
\newcommand{\RateBSixRJudgeOnlyCodeProp}{0.01}
\newcommand{\RateBSixRJudgeOnlyCodeDet}{0.64}
\newcommand{\RateBSixRJudgeOnlyCodeRec}{0.64}
\newcommand{\RateBSixRJudgeOnlyCodeEmit}{0.58}
\newcommand{\RateBSixRJudgeOnlyCodeDetEmit}{0.90}
\newcommand{\TokBSixRJudgeOnlyCode}{8.1}
\newcommand{\CountBSixRJudgeOnlyCodeTimeout}{0}
\newcommand{\CountBSixRJudgeOnlyCodeRetried}{0}
\newcommand{\CountBSixRJudgeOnlyCodeEmit}{58}
\newcommand{\RateBSevenBpdQnaAcc}{1.00}
\newcommand{\RateMainBpdQnaAcc}{1.00}
\newcommand{\RateBSevenBpdQnaProp}{0.00}
\newcommand{\RateMainBpdQnaProp}{0.00}
\newcommand{\RateBSevenBpdQnaDet}{0.94}
\newcommand{\RateMainBpdQnaDet}{0.94}
\newcommand{\RateBSevenBpdQnaRec}{0.94}
\newcommand{\RateMainBpdQnaRec}{0.94}
\newcommand{\RateBSevenBpdQnaEmit}{1.00}
\newcommand{\RateMainBpdQnaEmit}{1.00}
\newcommand{\RateBSevenBpdQnaDetEmit}{0.94}
\newcommand{\RateMainBpdQnaDetEmit}{0.94}
\newcommand{\TokBSevenBpdQna}{3.1}
\newcommand{\TokMainBpdQna}{3.1}
\newcommand{\CountBSevenBpdQnaTimeout}{0}
\newcommand{\CountMainBpdQnaTimeout}{0}
\newcommand{\CountBSevenBpdQnaRetried}{0}
\newcommand{\CountMainBpdQnaRetried}{0}
\newcommand{\CountBSevenBpdQnaEmit}{100}
\newcommand{\CountMainBpdQnaEmit}{100}
\newcommand{\RateBSevenBpdCodeAcc}{0.99}
\newcommand{\RateMainBpdCodeAcc}{0.99}
\newcommand{\RateBSevenBpdCodeProp}{0.01}
\newcommand{\RateMainBpdCodeProp}{0.01}
\newcommand{\RateBSevenBpdCodeDet}{0.67}
\newcommand{\RateMainBpdCodeDet}{0.67}
\newcommand{\RateBSevenBpdCodeRec}{0.67}
\newcommand{\RateMainBpdCodeRec}{0.67}
\newcommand{\RateBSevenBpdCodeEmit}{0.61}
\newcommand{\RateMainBpdCodeEmit}{0.61}
\newcommand{\RateBSevenBpdCodeDetEmit}{0.84}
\newcommand{\RateMainBpdCodeDetEmit}{0.84}
\newcommand{\TokBSevenBpdCode}{8.1}
\newcommand{\TokMainBpdCode}{8.1}
\newcommand{\CountBSevenBpdCodeTimeout}{0}
\newcommand{\CountMainBpdCodeTimeout}{0}
\newcommand{\CountBSevenBpdCodeRetried}{0}
\newcommand{\CountMainBpdCodeRetried}{0}
\newcommand{\CountBSevenBpdCodeEmit}{61}
\newcommand{\CountMainBpdCodeEmit}{61}
\newcommand{\RateBSevenFlatThreeQnaAcc}{1.00}
\newcommand{\RateMainFlatThreeQnaAcc}{1.00}
\newcommand{\RateBSevenFlatThreeQnaProp}{0.00}
\newcommand{\RateMainFlatThreeQnaProp}{0.00}
\newcommand{\RateBSevenFlatThreeQnaDet}{0.99}
\newcommand{\RateMainFlatThreeQnaDet}{0.99}
\newcommand{\RateBSevenFlatThreeQnaRec}{0.99}
\newcommand{\RateMainFlatThreeQnaRec}{0.99}
\newcommand{\RateBSevenFlatThreeQnaEmit}{0.99}
\newcommand{\RateMainFlatThreeQnaEmit}{0.99}
\newcommand{\RateBSevenFlatThreeQnaDetEmit}{1.00}
\newcommand{\RateMainFlatThreeQnaDetEmit}{1.00}
\newcommand{\TokBSevenFlatThreeQna}{3.0}
\newcommand{\TokMainFlatThreeQna}{3.0}
\newcommand{\CountBSevenFlatThreeQnaTimeout}{0}
\newcommand{\CountMainFlatThreeQnaTimeout}{0}
\newcommand{\CountBSevenFlatThreeQnaRetried}{0}
\newcommand{\CountMainFlatThreeQnaRetried}{0}
\newcommand{\CountBSevenFlatThreeQnaEmit}{99}
\newcommand{\CountMainFlatThreeQnaEmit}{99}
\newcommand{\RateBSevenFlatThreeCodeAcc}{0.72}
\newcommand{\RateMainFlatThreeCodeAcc}{0.72}
\newcommand{\RateBSevenFlatThreeCodeProp}{0.26}
\newcommand{\RateMainFlatThreeCodeProp}{0.26}
\newcommand{\RateBSevenFlatThreeCodeDet}{0.01}
\newcommand{\RateMainFlatThreeCodeDet}{0.01}
\newcommand{\RateBSevenFlatThreeCodeRec}{0.01}
\newcommand{\RateMainFlatThreeCodeRec}{0.01}
\newcommand{\RateBSevenFlatThreeCodeEmit}{0.59}
\newcommand{\RateMainFlatThreeCodeEmit}{0.59}
\newcommand{\RateBSevenFlatThreeCodeDetEmit}{0.02}
\newcommand{\RateMainFlatThreeCodeDetEmit}{0.02}
\newcommand{\TokBSevenFlatThreeCode}{8.1}
\newcommand{\TokMainFlatThreeCode}{8.1}
\newcommand{\CountBSevenFlatThreeCodeTimeout}{0}
\newcommand{\CountMainFlatThreeCodeTimeout}{0}
\newcommand{\CountBSevenFlatThreeCodeRetried}{0}
\newcommand{\CountMainFlatThreeCodeRetried}{0}
\newcommand{\CountBSevenFlatThreeCodeEmit}{59}
\newcommand{\CountMainFlatThreeCodeEmit}{59}
\newcommand{\RateBSevenGenOneQnaAcc}{0.00}
\newcommand{\RateMainGenOneQnaAcc}{0.00}
\newcommand{\RateBSevenGenOneQnaProp}{1.00}
\newcommand{\RateMainGenOneQnaProp}{1.00}
\newcommand{\RateBSevenGenOneQnaDet}{0.00}
\newcommand{\RateMainGenOneQnaDet}{0.00}
\newcommand{\RateBSevenGenOneQnaRec}{0.00}
\newcommand{\RateMainGenOneQnaRec}{0.00}
\newcommand{\RateBSevenGenOneQnaEmit}{1.00}
\newcommand{\RateMainGenOneQnaEmit}{1.00}
\newcommand{\RateBSevenGenOneQnaDetEmit}{0.00}
\newcommand{\RateMainGenOneQnaDetEmit}{0.00}
\newcommand{\TokBSevenGenOneQna}{0.6}
\newcommand{\TokMainGenOneQna}{0.6}
\newcommand{\CountBSevenGenOneQnaTimeout}{0}
\newcommand{\CountMainGenOneQnaTimeout}{0}
\newcommand{\CountBSevenGenOneQnaRetried}{0}
\newcommand{\CountMainGenOneQnaRetried}{0}
\newcommand{\CountBSevenGenOneQnaEmit}{100}
\newcommand{\CountMainGenOneQnaEmit}{100}
\newcommand{\RateBSevenGenOneCodeAcc}{0.59}
\newcommand{\RateMainGenOneCodeAcc}{0.59}
\newcommand{\RateBSevenGenOneCodeProp}{0.40}
\newcommand{\RateMainGenOneCodeProp}{0.40}
\newcommand{\RateBSevenGenOneCodeDet}{0.00}
\newcommand{\RateMainGenOneCodeDet}{0.00}
\newcommand{\RateBSevenGenOneCodeRec}{0.00}
\newcommand{\RateMainGenOneCodeRec}{0.00}
\newcommand{\RateBSevenGenOneCodeEmit}{0.47}
\newcommand{\RateMainGenOneCodeEmit}{0.47}
\newcommand{\RateBSevenGenOneCodeDetEmit}{0.00}
\newcommand{\RateMainGenOneCodeDetEmit}{0.00}
\newcommand{\TokBSevenGenOneCode}{2.8}
\newcommand{\TokMainGenOneCode}{2.8}
\newcommand{\CountBSevenGenOneCodeTimeout}{0}
\newcommand{\CountMainGenOneCodeTimeout}{0}
\newcommand{\CountBSevenGenOneCodeRetried}{0}
\newcommand{\CountMainGenOneCodeRetried}{0}
\newcommand{\CountBSevenGenOneCodeEmit}{47}
\newcommand{\CountMainGenOneCodeEmit}{47}
\newcommand{\RateBSevenGenThreeQnaAcc}{1.00}
\newcommand{\RateMainGenThreeQnaAcc}{1.00}
\newcommand{\RateBSevenGenThreeQnaProp}{0.00}
\newcommand{\RateMainGenThreeQnaProp}{0.00}
\newcommand{\RateBSevenGenThreeQnaDet}{1.00}
\newcommand{\RateMainGenThreeQnaDet}{1.00}
\newcommand{\RateBSevenGenThreeQnaRec}{1.00}
\newcommand{\RateMainGenThreeQnaRec}{1.00}
\newcommand{\RateBSevenGenThreeQnaEmit}{1.00}
\newcommand{\RateMainGenThreeQnaEmit}{1.00}
\newcommand{\RateBSevenGenThreeQnaDetEmit}{1.00}
\newcommand{\RateMainGenThreeQnaDetEmit}{1.00}
\newcommand{\TokBSevenGenThreeQna}{3.0}
\newcommand{\TokMainGenThreeQna}{3.0}
\newcommand{\CountBSevenGenThreeQnaTimeout}{0}
\newcommand{\CountMainGenThreeQnaTimeout}{0}
\newcommand{\CountBSevenGenThreeQnaRetried}{0}
\newcommand{\CountMainGenThreeQnaRetried}{0}
\newcommand{\CountBSevenGenThreeQnaEmit}{100}
\newcommand{\CountMainGenThreeQnaEmit}{100}
\newcommand{\RateBSevenGenThreeCodeAcc}{0.81}
\newcommand{\RateMainGenThreeCodeAcc}{0.81}
\newcommand{\RateBSevenGenThreeCodeProp}{0.15}
\newcommand{\RateMainGenThreeCodeProp}{0.15}
\newcommand{\RateBSevenGenThreeCodeDet}{0.02}
\newcommand{\RateMainGenThreeCodeDet}{0.02}
\newcommand{\RateBSevenGenThreeCodeRec}{0.02}
\newcommand{\RateMainGenThreeCodeRec}{0.02}
\newcommand{\RateBSevenGenThreeCodeEmit}{0.48}
\newcommand{\RateMainGenThreeCodeEmit}{0.48}
\newcommand{\RateBSevenGenThreeCodeDetEmit}{0.04}
\newcommand{\RateMainGenThreeCodeDetEmit}{0.04}
\newcommand{\TokBSevenGenThreeCode}{8.4}
\newcommand{\TokMainGenThreeCode}{8.4}
\newcommand{\CountBSevenGenThreeCodeTimeout}{0}
\newcommand{\CountMainGenThreeCodeTimeout}{0}
\newcommand{\CountBSevenGenThreeCodeRetried}{0}
\newcommand{\CountMainGenThreeCodeRetried}{0}
\newcommand{\CountBSevenGenThreeCodeEmit}{48}
\newcommand{\CountMainGenThreeCodeEmit}{48}
\newcommand{\RateBSevenJudgeOnlyQnaAcc}{1.00}
\newcommand{\RateMainJudgeOnlyQnaAcc}{1.00}
\newcommand{\RateBSevenJudgeOnlyQnaProp}{0.00}
\newcommand{\RateMainJudgeOnlyQnaProp}{0.00}
\newcommand{\RateBSevenJudgeOnlyQnaDet}{0.91}
\newcommand{\RateMainJudgeOnlyQnaDet}{0.91}
\newcommand{\RateBSevenJudgeOnlyQnaRec}{0.91}
\newcommand{\RateMainJudgeOnlyQnaRec}{0.91}
\newcommand{\RateBSevenJudgeOnlyQnaEmit}{0.98}
\newcommand{\RateMainJudgeOnlyQnaEmit}{0.98}
\newcommand{\RateBSevenJudgeOnlyQnaDetEmit}{0.93}
\newcommand{\RateMainJudgeOnlyQnaDetEmit}{0.93}
\newcommand{\TokBSevenJudgeOnlyQna}{3.0}
\newcommand{\TokMainJudgeOnlyQna}{3.0}
\newcommand{\CountBSevenJudgeOnlyQnaTimeout}{0}
\newcommand{\CountMainJudgeOnlyQnaTimeout}{0}
\newcommand{\CountBSevenJudgeOnlyQnaRetried}{0}
\newcommand{\CountMainJudgeOnlyQnaRetried}{0}
\newcommand{\CountBSevenJudgeOnlyQnaEmit}{98}
\newcommand{\CountMainJudgeOnlyQnaEmit}{98}
\newcommand{\RateBSevenJudgeOnlyCodeAcc}{1.00}
\newcommand{\RateMainJudgeOnlyCodeAcc}{1.00}
\newcommand{\RateBSevenJudgeOnlyCodeProp}{0.00}
\newcommand{\RateMainJudgeOnlyCodeProp}{0.00}
\newcommand{\RateBSevenJudgeOnlyCodeDet}{0.63}
\newcommand{\RateMainJudgeOnlyCodeDet}{0.63}
\newcommand{\RateBSevenJudgeOnlyCodeRec}{0.63}
\newcommand{\RateMainJudgeOnlyCodeRec}{0.63}
\newcommand{\RateBSevenJudgeOnlyCodeEmit}{0.61}
\newcommand{\RateMainJudgeOnlyCodeEmit}{0.61}
\newcommand{\RateBSevenJudgeOnlyCodeDetEmit}{0.87}
\newcommand{\RateMainJudgeOnlyCodeDetEmit}{0.87}
\newcommand{\TokBSevenJudgeOnlyCode}{8.1}
\newcommand{\TokMainJudgeOnlyCode}{8.1}
\newcommand{\CountBSevenJudgeOnlyCodeTimeout}{0}
\newcommand{\CountMainJudgeOnlyCodeTimeout}{0}
\newcommand{\CountBSevenJudgeOnlyCodeRetried}{0}
\newcommand{\CountMainJudgeOnlyCodeRetried}{0}
\newcommand{\CountBSevenJudgeOnlyCodeEmit}{61}
\newcommand{\CountMainJudgeOnlyCodeEmit}{61}
\newcommand{\RateBEightBpdQnaAcc}{1.00}
\newcommand{\RateBEightBpdQnaProp}{0.00}
\newcommand{\RateBEightBpdQnaDet}{0.10}
\newcommand{\RateBEightBpdQnaRec}{0.10}
\newcommand{\RateBEightBpdQnaEmit}{0.10}
\newcommand{\RateBEightBpdQnaDetEmit}{1.00}
\newcommand{\TokBEightBpdQna}{4.5}
\newcommand{\CountBEightBpdQnaTimeout}{0}
\newcommand{\CountBEightBpdQnaRetried}{0}
\newcommand{\CountBEightBpdQnaEmit}{10}
\newcommand{\RateBEightBpdCodeAcc}{0.98}
\newcommand{\RateBEightBpdCodeProp}{0.02}
\newcommand{\RateBEightBpdCodeDet}{0.61}
\newcommand{\RateBEightBpdCodeRec}{0.61}
\newcommand{\RateBEightBpdCodeEmit}{0.53}
\newcommand{\RateBEightBpdCodeDetEmit}{0.96}
\newcommand{\TokBEightBpdCode}{8.2}
\newcommand{\CountBEightBpdCodeTimeout}{0}
\newcommand{\CountBEightBpdCodeRetried}{0}
\newcommand{\CountBEightBpdCodeEmit}{53}
\newcommand{\RateBEightFlatThreeQnaAcc}{1.00}
\newcommand{\RateBEightFlatThreeQnaProp}{0.00}
\newcommand{\RateBEightFlatThreeQnaDet}{0.12}
\newcommand{\RateBEightFlatThreeQnaRec}{0.12}
\newcommand{\RateBEightFlatThreeQnaEmit}{0.12}
\newcommand{\RateBEightFlatThreeQnaDetEmit}{1.00}
\newcommand{\TokBEightFlatThreeQna}{4.5}
\newcommand{\CountBEightFlatThreeQnaTimeout}{0}
\newcommand{\CountBEightFlatThreeQnaRetried}{0}
\newcommand{\CountBEightFlatThreeQnaEmit}{12}
\newcommand{\RateBEightFlatThreeCodeAcc}{0.70}
\newcommand{\RateBEightFlatThreeCodeProp}{0.29}
\newcommand{\RateBEightFlatThreeCodeDet}{0.01}
\newcommand{\RateBEightFlatThreeCodeRec}{0.01}
\newcommand{\RateBEightFlatThreeCodeEmit}{0.58}
\newcommand{\RateBEightFlatThreeCodeDetEmit}{0.02}
\newcommand{\TokBEightFlatThreeCode}{8.1}
\newcommand{\CountBEightFlatThreeCodeTimeout}{0}
\newcommand{\CountBEightFlatThreeCodeRetried}{0}
\newcommand{\CountBEightFlatThreeCodeEmit}{58}
\newcommand{\RateBEightGenOneQnaAcc}{0.05}
\newcommand{\RateBEightGenOneQnaProp}{0.95}
\newcommand{\RateBEightGenOneQnaDet}{0.00}
\newcommand{\RateBEightGenOneQnaRec}{0.00}
\newcommand{\RateBEightGenOneQnaEmit}{0.95}
\newcommand{\RateBEightGenOneQnaDetEmit}{0.00}
\newcommand{\TokBEightGenOneQna}{0.7}
\newcommand{\CountBEightGenOneQnaTimeout}{0}
\newcommand{\CountBEightGenOneQnaRetried}{0}
\newcommand{\CountBEightGenOneQnaEmit}{95}
\newcommand{\RateBEightGenOneCodeAcc}{0.44}
\newcommand{\RateBEightGenOneCodeProp}{0.46}
\newcommand{\RateBEightGenOneCodeDet}{0.00}
\newcommand{\RateBEightGenOneCodeRec}{0.00}
\newcommand{\RateBEightGenOneCodeEmit}{0.48}
\newcommand{\RateBEightGenOneCodeDetEmit}{0.00}
\newcommand{\TokBEightGenOneCode}{3.2}
\newcommand{\CountBEightGenOneCodeTimeout}{0}
\newcommand{\CountBEightGenOneCodeRetried}{0}
\newcommand{\CountBEightGenOneCodeEmit}{48}
\newcommand{\RateBEightGenThreeQnaAcc}{1.00}
\newcommand{\RateBEightGenThreeQnaProp}{0.00}
\newcommand{\RateBEightGenThreeQnaDet}{0.95}
\newcommand{\RateBEightGenThreeQnaRec}{0.95}
\newcommand{\RateBEightGenThreeQnaEmit}{0.95}
\newcommand{\RateBEightGenThreeQnaDetEmit}{1.00}
\newcommand{\TokBEightGenThreeQna}{3.1}
\newcommand{\CountBEightGenThreeQnaTimeout}{0}
\newcommand{\CountBEightGenThreeQnaRetried}{0}
\newcommand{\CountBEightGenThreeQnaEmit}{95}
\newcommand{\RateBEightGenThreeCodeAcc}{0.72}
\newcommand{\RateBEightGenThreeCodeProp}{0.27}
\newcommand{\RateBEightGenThreeCodeDet}{0.01}
\newcommand{\RateBEightGenThreeCodeRec}{0.01}
\newcommand{\RateBEightGenThreeCodeEmit}{0.51}
\newcommand{\RateBEightGenThreeCodeDetEmit}{0.02}
\newcommand{\TokBEightGenThreeCode}{8.9}
\newcommand{\CountBEightGenThreeCodeTimeout}{0}
\newcommand{\CountBEightGenThreeCodeRetried}{0}
\newcommand{\CountBEightGenThreeCodeEmit}{51}
\newcommand{\RateBEightJudgeOnlyQnaAcc}{1.00}
\newcommand{\RateBEightJudgeOnlyQnaProp}{0.00}
\newcommand{\RateBEightJudgeOnlyQnaDet}{0.11}
\newcommand{\RateBEightJudgeOnlyQnaRec}{0.11}
\newcommand{\RateBEightJudgeOnlyQnaEmit}{0.12}
\newcommand{\RateBEightJudgeOnlyQnaDetEmit}{0.92}
\newcommand{\TokBEightJudgeOnlyQna}{4.5}
\newcommand{\CountBEightJudgeOnlyQnaTimeout}{0}
\newcommand{\CountBEightJudgeOnlyQnaRetried}{0}
\newcommand{\CountBEightJudgeOnlyQnaEmit}{12}
\newcommand{\RateBEightJudgeOnlyCodeAcc}{0.99}
\newcommand{\RateBEightJudgeOnlyCodeProp}{0.01}
\newcommand{\RateBEightJudgeOnlyCodeDet}{0.60}
\newcommand{\RateBEightJudgeOnlyCodeRec}{0.60}
\newcommand{\RateBEightJudgeOnlyCodeEmit}{0.54}
\newcommand{\RateBEightJudgeOnlyCodeDetEmit}{0.98}
\newcommand{\TokBEightJudgeOnlyCode}{8.1}
\newcommand{\CountBEightJudgeOnlyCodeTimeout}{0}
\newcommand{\CountBEightJudgeOnlyCodeRetried}{0}
\newcommand{\CountBEightJudgeOnlyCodeEmit}{54}
\newcommand{\RateConditionsBpdQnaAcc}{1.00}
\newcommand{\RateConditionsBpdQnaProp}{0.00}
\newcommand{\RateConditionsBpdQnaDet}{0.29}
\newcommand{\RateConditionsBpdQnaRec}{0.29}
\newcommand{\TokConditionsBpdQna}{3.4}
\newcommand{\CountConditionsBpdQnaTimeout}{0}
\newcommand{\CountConditionsBpdQnaRetried}{0}
\newcommand{\RateConditionsBpdCodeAcc}{0.98}
\newcommand{\RateConditionsBpdCodeProp}{0.01}
\newcommand{\RateConditionsBpdCodeDet}{0.21}
\newcommand{\RateConditionsBpdCodeRec}{0.21}
\newcommand{\TokConditionsBpdCode}{7.9}
\newcommand{\CountConditionsBpdCodeTimeout}{3}
\newcommand{\CountConditionsBpdCodeRetried}{0}
\newcommand{\RateConditionsYonseiQnaAcc}{0.99}
\newcommand{\RateConditionsYonseiQnaProp}{0.00}
\newcommand{\RateConditionsYonseiQnaDet}{0.33}
\newcommand{\RateConditionsYonseiQnaRec}{0.33}
\newcommand{\TokConditionsYonseiQna}{1.3}
\newcommand{\CountConditionsYonseiQnaTimeout}{3}
\newcommand{\CountConditionsYonseiQnaRetried}{0}
\newcommand{\RateConditionsYonseiCodeAcc}{1.00}
\newcommand{\RateConditionsYonseiCodeProp}{0.00}
\newcommand{\RateConditionsYonseiCodeDet}{0.33}
\newcommand{\RateConditionsYonseiCodeRec}{0.33}
\newcommand{\TokConditionsYonseiCode}{2.4}
\newcommand{\CountConditionsYonseiCodeTimeout}{0}
\newcommand{\CountConditionsYonseiCodeRetried}{0}
\newcommand{\RateConditionsGenOneQnaAcc}{0.93}
\newcommand{\RateConditionsGenOneQnaProp}{0.07}
\newcommand{\RateConditionsGenOneQnaDet}{0.00}
\newcommand{\RateConditionsGenOneQnaRec}{0.00}
\newcommand{\TokConditionsGenOneQna}{1.2}
\newcommand{\CountConditionsGenOneQnaTimeout}{0}
\newcommand{\CountConditionsGenOneQnaRetried}{0}
\newcommand{\RateConditionsGenOneCodeAcc}{0.92}
\newcommand{\RateConditionsGenOneCodeProp}{0.06}
\newcommand{\RateConditionsGenOneCodeDet}{0.00}
\newcommand{\RateConditionsGenOneCodeRec}{0.00}
\newcommand{\TokConditionsGenOneCode}{2.7}
\newcommand{\CountConditionsGenOneCodeTimeout}{0}
\newcommand{\CountConditionsGenOneCodeRetried}{0}
\newcommand{\RateYonseiImprovedYonseiQnaAcc}{1.00}
\newcommand{\RateYonseiImprovedYonseiQnaProp}{0.00}
\newcommand{\RateYonseiImprovedYonseiQnaDet}{1.00}
\newcommand{\RateYonseiImprovedYonseiQnaRec}{1.00}
\newcommand{\TokYonseiImprovedYonseiQna}{1.4}
\newcommand{\CountYonseiImprovedYonseiQnaTimeout}{0}
\newcommand{\CountYonseiImprovedYonseiQnaRetried}{0}
\newcommand{\RateYonseiImprovedYonseiCodeAcc}{1.00}
\newcommand{\RateYonseiImprovedYonseiCodeProp}{0.00}
\newcommand{\RateYonseiImprovedYonseiCodeDet}{1.00}
\newcommand{\RateYonseiImprovedYonseiCodeRec}{1.00}
\newcommand{\TokYonseiImprovedYonseiCode}{2.6}
\newcommand{\CountYonseiImprovedYonseiCodeTimeout}{0}
\newcommand{\CountYonseiImprovedYonseiCodeRetried}{0}
\newcommand{\PvalBFiveQnaAccFlatThreevsBpd}{$p=1.00$}
\newcommand{\PvalBFiveQnaPropFlatThreevsBpd}{$p=1.00$}
\newcommand{\PvalBFiveQnaDetFlatThreevsBpd}{$p<0.001$}
\newcommand{\PvalBFiveQnaRecFlatThreevsBpd}{$p<0.001$}
\newcommand{\PvalBFiveQnaAccGenOnevsBpd}{$p=1.00$}
\newcommand{\PvalBFiveQnaPropGenOnevsBpd}{$p=1.00$}
\newcommand{\PvalBFiveQnaDetGenOnevsBpd}{$p<0.001$}
\newcommand{\PvalBFiveQnaRecGenOnevsBpd}{$p<0.001$}
\newcommand{\PvalBFiveQnaAccGenThreevsBpd}{$p=1.00$}
\newcommand{\PvalBFiveQnaPropGenThreevsBpd}{$p=1.00$}
\newcommand{\PvalBFiveQnaDetGenThreevsBpd}{$p<0.001$}
\newcommand{\PvalBFiveQnaRecGenThreevsBpd}{$p<0.001$}
\newcommand{\PvalBFiveQnaAccJudgeOnlyvsBpd}{$p=1.00$}
\newcommand{\PvalBFiveQnaPropJudgeOnlyvsBpd}{$p=1.00$}
\newcommand{\PvalBFiveQnaDetJudgeOnlyvsBpd}{$p=0.503$}
\newcommand{\PvalBFiveQnaRecJudgeOnlyvsBpd}{$p=0.503$}
\newcommand{\PvalBFiveQnaAccGenOnevsFlatThree}{$p=1.00$}
\newcommand{\PvalBFiveQnaPropGenOnevsFlatThree}{$p=1.00$}
\newcommand{\PvalBFiveQnaDetGenOnevsFlatThree}{$p<0.001$}
\newcommand{\PvalBFiveQnaRecGenOnevsFlatThree}{$p<0.001$}
\newcommand{\PvalBFiveQnaAccGenThreevsFlatThree}{$p=1.00$}
\newcommand{\PvalBFiveQnaPropGenThreevsFlatThree}{$p=1.00$}
\newcommand{\PvalBFiveQnaDetGenThreevsFlatThree}{$p<0.001$}
\newcommand{\PvalBFiveQnaRecGenThreevsFlatThree}{$p<0.001$}
\newcommand{\PvalBFiveQnaAccFlatThreevsJudgeOnly}{$p=1.00$}
\newcommand{\PvalBFiveQnaPropFlatThreevsJudgeOnly}{$p=1.00$}
\newcommand{\PvalBFiveQnaDetFlatThreevsJudgeOnly}{$p<0.001$}
\newcommand{\PvalBFiveQnaRecFlatThreevsJudgeOnly}{$p<0.001$}
\newcommand{\PvalBFiveQnaAccGenOnevsGenThree}{$p=1.00$}
\newcommand{\PvalBFiveQnaPropGenOnevsGenThree}{$p=1.00$}
\newcommand{\PvalBFiveQnaDetGenOnevsGenThree}{$p=1.00$}
\newcommand{\PvalBFiveQnaRecGenOnevsGenThree}{$p=1.00$}
\newcommand{\PvalBFiveQnaAccGenOnevsJudgeOnly}{$p=1.00$}
\newcommand{\PvalBFiveQnaPropGenOnevsJudgeOnly}{$p=1.00$}
\newcommand{\PvalBFiveQnaDetGenOnevsJudgeOnly}{$p<0.001$}
\newcommand{\PvalBFiveQnaRecGenOnevsJudgeOnly}{$p<0.001$}
\newcommand{\PvalBFiveQnaAccGenThreevsJudgeOnly}{$p=1.00$}
\newcommand{\PvalBFiveQnaPropGenThreevsJudgeOnly}{$p=1.00$}
\newcommand{\PvalBFiveQnaDetGenThreevsJudgeOnly}{$p<0.001$}
\newcommand{\PvalBFiveQnaRecGenThreevsJudgeOnly}{$p<0.001$}
\newcommand{\PvalBFiveCodeAccFlatThreevsBpd}{$p=0.008$}
\newcommand{\PvalBFiveCodePropFlatThreevsBpd}{$p=0.008$}
\newcommand{\PvalBFiveCodeDetFlatThreevsBpd}{$p<0.001$}
\newcommand{\PvalBFiveCodeRecFlatThreevsBpd}{$p<0.001$}
\newcommand{\PvalBFiveCodeAccGenOnevsBpd}{$p<0.001$}
\newcommand{\PvalBFiveCodePropGenOnevsBpd}{$p=0.002$}
\newcommand{\PvalBFiveCodeDetGenOnevsBpd}{$p<0.001$}
\newcommand{\PvalBFiveCodeRecGenOnevsBpd}{$p<0.001$}
\newcommand{\PvalBFiveCodeAccGenThreevsBpd}{$p=1.00$}
\newcommand{\PvalBFiveCodePropGenThreevsBpd}{$p=1.00$}
\newcommand{\PvalBFiveCodeDetGenThreevsBpd}{$p<0.001$}
\newcommand{\PvalBFiveCodeRecGenThreevsBpd}{$p<0.001$}
\newcommand{\PvalBFiveCodeAccJudgeOnlyvsBpd}{$p=1.00$}
\newcommand{\PvalBFiveCodePropJudgeOnlyvsBpd}{$p=1.00$}
\newcommand{\PvalBFiveCodeDetJudgeOnlyvsBpd}{$p=1.00$}
\newcommand{\PvalBFiveCodeRecJudgeOnlyvsBpd}{$p=1.00$}
\newcommand{\PvalBFiveCodeAccGenOnevsFlatThree}{$p=0.189$}
\newcommand{\PvalBFiveCodePropGenOnevsFlatThree}{$p=0.804$}
\newcommand{\PvalBFiveCodeDetGenOnevsFlatThree}{$p=1.00$}
\newcommand{\PvalBFiveCodeRecGenOnevsFlatThree}{$p=1.00$}
\newcommand{\PvalBFiveCodeAccGenThreevsFlatThree}{$p=0.004$}
\newcommand{\PvalBFiveCodePropGenThreevsFlatThree}{$p=0.004$}
\newcommand{\PvalBFiveCodeDetGenThreevsFlatThree}{$p=1.00$}
\newcommand{\PvalBFiveCodeRecGenThreevsFlatThree}{$p=1.00$}
\newcommand{\PvalBFiveCodeAccFlatThreevsJudgeOnly}{$p=0.004$}
\newcommand{\PvalBFiveCodePropFlatThreevsJudgeOnly}{$p=0.004$}
\newcommand{\PvalBFiveCodeDetFlatThreevsJudgeOnly}{$p<0.001$}
\newcommand{\PvalBFiveCodeRecFlatThreevsJudgeOnly}{$p<0.001$}
\newcommand{\PvalBFiveCodeAccGenOnevsGenThree}{$p<0.001$}
\newcommand{\PvalBFiveCodePropGenOnevsGenThree}{$p<0.001$}
\newcommand{\PvalBFiveCodeDetGenOnevsGenThree}{$p=1.00$}
\newcommand{\PvalBFiveCodeRecGenOnevsGenThree}{$p=1.00$}
\newcommand{\PvalBFiveCodeAccGenOnevsJudgeOnly}{$p<0.001$}
\newcommand{\PvalBFiveCodePropGenOnevsJudgeOnly}{$p<0.001$}
\newcommand{\PvalBFiveCodeDetGenOnevsJudgeOnly}{$p<0.001$}
\newcommand{\PvalBFiveCodeRecGenOnevsJudgeOnly}{$p<0.001$}
\newcommand{\PvalBFiveCodeAccGenThreevsJudgeOnly}{$p=1.00$}
\newcommand{\PvalBFiveCodePropGenThreevsJudgeOnly}{$p=1.00$}
\newcommand{\PvalBFiveCodeDetGenThreevsJudgeOnly}{$p<0.001$}
\newcommand{\PvalBFiveCodeRecGenThreevsJudgeOnly}{$p<0.001$}
\newcommand{\PvalBSixRQnaAccFlatThreevsBpd}{$p=1.00$}
\newcommand{\PvalBSixRQnaPropFlatThreevsBpd}{$p=1.00$}
\newcommand{\PvalBSixRQnaDetFlatThreevsBpd}{$p=0.002$}
\newcommand{\PvalBSixRQnaRecFlatThreevsBpd}{$p=0.002$}
\newcommand{\PvalBSixRQnaAccGenOnevsBpd}{$p<0.001$}
\newcommand{\PvalBSixRQnaPropGenOnevsBpd}{$p<0.001$}
\newcommand{\PvalBSixRQnaDetGenOnevsBpd}{$p<0.001$}
\newcommand{\PvalBSixRQnaRecGenOnevsBpd}{$p<0.001$}
\newcommand{\PvalBSixRQnaAccGenThreevsBpd}{$p=1.00$}
\newcommand{\PvalBSixRQnaPropGenThreevsBpd}{$p=1.00$}
\newcommand{\PvalBSixRQnaDetGenThreevsBpd}{$p<0.001$}
\newcommand{\PvalBSixRQnaRecGenThreevsBpd}{$p<0.001$}
\newcommand{\PvalBSixRQnaAccJudgeOnlyvsBpd}{$p=1.00$}
\newcommand{\PvalBSixRQnaPropJudgeOnlyvsBpd}{$p=1.00$}
\newcommand{\PvalBSixRQnaDetJudgeOnlyvsBpd}{$p=0.832$}
\newcommand{\PvalBSixRQnaRecJudgeOnlyvsBpd}{$p=0.832$}
\newcommand{\PvalBSixRQnaAccGenOnevsFlatThree}{$p<0.001$}
\newcommand{\PvalBSixRQnaPropGenOnevsFlatThree}{$p<0.001$}
\newcommand{\PvalBSixRQnaDetGenOnevsFlatThree}{$p<0.001$}
\newcommand{\PvalBSixRQnaRecGenOnevsFlatThree}{$p<0.001$}
\newcommand{\PvalBSixRQnaAccGenThreevsFlatThree}{$p=1.00$}
\newcommand{\PvalBSixRQnaPropGenThreevsFlatThree}{$p=1.00$}
\newcommand{\PvalBSixRQnaDetGenThreevsFlatThree}{$p=1.00$}
\newcommand{\PvalBSixRQnaRecGenThreevsFlatThree}{$p=1.00$}
\newcommand{\PvalBSixRQnaAccFlatThreevsJudgeOnly}{$p=1.00$}
\newcommand{\PvalBSixRQnaPropFlatThreevsJudgeOnly}{$p=1.00$}
\newcommand{\PvalBSixRQnaDetFlatThreevsJudgeOnly}{$p<0.001$}
\newcommand{\PvalBSixRQnaRecFlatThreevsJudgeOnly}{$p<0.001$}
\newcommand{\PvalBSixRQnaAccGenOnevsGenThree}{$p<0.001$}
\newcommand{\PvalBSixRQnaPropGenOnevsGenThree}{$p<0.001$}
\newcommand{\PvalBSixRQnaDetGenOnevsGenThree}{$p<0.001$}
\newcommand{\PvalBSixRQnaRecGenOnevsGenThree}{$p<0.001$}
\newcommand{\PvalBSixRQnaAccGenOnevsJudgeOnly}{$p<0.001$}
\newcommand{\PvalBSixRQnaPropGenOnevsJudgeOnly}{$p<0.001$}
\newcommand{\PvalBSixRQnaDetGenOnevsJudgeOnly}{$p<0.001$}
\newcommand{\PvalBSixRQnaRecGenOnevsJudgeOnly}{$p<0.001$}
\newcommand{\PvalBSixRQnaAccGenThreevsJudgeOnly}{$p=1.00$}
\newcommand{\PvalBSixRQnaPropGenThreevsJudgeOnly}{$p=1.00$}
\newcommand{\PvalBSixRQnaDetGenThreevsJudgeOnly}{$p<0.001$}
\newcommand{\PvalBSixRQnaRecGenThreevsJudgeOnly}{$p<0.001$}
\newcommand{\PvalBSixRCodeAccFlatThreevsBpd}{$p=0.388$}
\newcommand{\PvalBSixRCodePropFlatThreevsBpd}{$p=0.039$}
\newcommand{\PvalBSixRCodeDetFlatThreevsBpd}{$p<0.001$}
\newcommand{\PvalBSixRCodeRecFlatThreevsBpd}{$p<0.001$}
\newcommand{\PvalBSixRCodeAccGenOnevsBpd}{$p<0.001$}
\newcommand{\PvalBSixRCodePropGenOnevsBpd}{$p<0.001$}
\newcommand{\PvalBSixRCodeDetGenOnevsBpd}{$p<0.001$}
\newcommand{\PvalBSixRCodeRecGenOnevsBpd}{$p<0.001$}
\newcommand{\PvalBSixRCodeAccGenThreevsBpd}{$p=1.00$}
\newcommand{\PvalBSixRCodePropGenThreevsBpd}{$p=1.00$}
\newcommand{\PvalBSixRCodeDetGenThreevsBpd}{$p<0.001$}
\newcommand{\PvalBSixRCodeRecGenThreevsBpd}{$p<0.001$}
\newcommand{\PvalBSixRCodeAccJudgeOnlyvsBpd}{$p=0.375$}
\newcommand{\PvalBSixRCodePropJudgeOnlyvsBpd}{$p=1.00$}
\newcommand{\PvalBSixRCodeDetJudgeOnlyvsBpd}{$p=1.00$}
\newcommand{\PvalBSixRCodeRecJudgeOnlyvsBpd}{$p=1.00$}
\newcommand{\PvalBSixRCodeAccGenOnevsFlatThree}{$p<0.001$}
\newcommand{\PvalBSixRCodePropGenOnevsFlatThree}{$p<0.001$}
\newcommand{\PvalBSixRCodeDetGenOnevsFlatThree}{$p=1.00$}
\newcommand{\PvalBSixRCodeRecGenOnevsFlatThree}{$p=1.00$}
\newcommand{\PvalBSixRCodeAccGenThreevsFlatThree}{$p=0.180$}
\newcommand{\PvalBSixRCodePropGenThreevsFlatThree}{$p=0.070$}
\newcommand{\PvalBSixRCodeDetGenThreevsFlatThree}{$p=1.00$}
\newcommand{\PvalBSixRCodeRecGenThreevsFlatThree}{$p=1.00$}
\newcommand{\PvalBSixRCodeAccFlatThreevsJudgeOnly}{$p=0.039$}
\newcommand{\PvalBSixRCodePropFlatThreevsJudgeOnly}{$p=0.039$}
\newcommand{\PvalBSixRCodeDetFlatThreevsJudgeOnly}{$p<0.001$}
\newcommand{\PvalBSixRCodeRecFlatThreevsJudgeOnly}{$p<0.001$}
\newcommand{\PvalBSixRCodeAccGenOnevsGenThree}{$p<0.001$}
\newcommand{\PvalBSixRCodePropGenOnevsGenThree}{$p<0.001$}
\newcommand{\PvalBSixRCodeDetGenOnevsGenThree}{$p=1.00$}
\newcommand{\PvalBSixRCodeRecGenOnevsGenThree}{$p=1.00$}
\newcommand{\PvalBSixRCodeAccGenOnevsJudgeOnly}{$p<0.001$}
\newcommand{\PvalBSixRCodePropGenOnevsJudgeOnly}{$p<0.001$}
\newcommand{\PvalBSixRCodeDetGenOnevsJudgeOnly}{$p<0.001$}
\newcommand{\PvalBSixRCodeRecGenOnevsJudgeOnly}{$p<0.001$}
\newcommand{\PvalBSixRCodeAccGenThreevsJudgeOnly}{$p=0.625$}
\newcommand{\PvalBSixRCodePropGenThreevsJudgeOnly}{$p=1.00$}
\newcommand{\PvalBSixRCodeDetGenThreevsJudgeOnly}{$p<0.001$}
\newcommand{\PvalBSixRCodeRecGenThreevsJudgeOnly}{$p<0.001$}
\newcommand{\PvalBSixQnaAccFlatThreevsBpd}{$p=1.00$}
\newcommand{\PvalBSixQnaPropFlatThreevsBpd}{$p=1.00$}
\newcommand{\PvalBSixQnaDetFlatThreevsBpd}{$p=0.002$}
\newcommand{\PvalBSixQnaRecFlatThreevsBpd}{$p=0.002$}
\newcommand{\PvalBSixQnaAccGenOnevsBpd}{$p<0.001$}
\newcommand{\PvalBSixQnaPropGenOnevsBpd}{$p<0.001$}
\newcommand{\PvalBSixQnaDetGenOnevsBpd}{$p<0.001$}
\newcommand{\PvalBSixQnaRecGenOnevsBpd}{$p<0.001$}
\newcommand{\PvalBSixQnaAccGenThreevsBpd}{$p=1.00$}
\newcommand{\PvalBSixQnaPropGenThreevsBpd}{$p=1.00$}
\newcommand{\PvalBSixQnaDetGenThreevsBpd}{$p<0.001$}
\newcommand{\PvalBSixQnaRecGenThreevsBpd}{$p<0.001$}
\newcommand{\PvalBSixQnaAccJudgeOnlyvsBpd}{$p=1.00$}
\newcommand{\PvalBSixQnaPropJudgeOnlyvsBpd}{$p=1.00$}
\newcommand{\PvalBSixQnaDetJudgeOnlyvsBpd}{$p=0.832$}
\newcommand{\PvalBSixQnaRecJudgeOnlyvsBpd}{$p=0.832$}
\newcommand{\PvalBSixQnaAccGenOnevsFlatThree}{$p<0.001$}
\newcommand{\PvalBSixQnaPropGenOnevsFlatThree}{$p<0.001$}
\newcommand{\PvalBSixQnaDetGenOnevsFlatThree}{$p<0.001$}
\newcommand{\PvalBSixQnaRecGenOnevsFlatThree}{$p<0.001$}
\newcommand{\PvalBSixQnaAccGenThreevsFlatThree}{$p=1.00$}
\newcommand{\PvalBSixQnaPropGenThreevsFlatThree}{$p=1.00$}
\newcommand{\PvalBSixQnaDetGenThreevsFlatThree}{$p=1.00$}
\newcommand{\PvalBSixQnaRecGenThreevsFlatThree}{$p=1.00$}
\newcommand{\PvalBSixQnaAccFlatThreevsJudgeOnly}{$p=1.00$}
\newcommand{\PvalBSixQnaPropFlatThreevsJudgeOnly}{$p=1.00$}
\newcommand{\PvalBSixQnaDetFlatThreevsJudgeOnly}{$p<0.001$}
\newcommand{\PvalBSixQnaRecFlatThreevsJudgeOnly}{$p<0.001$}
\newcommand{\PvalBSixQnaAccGenOnevsGenThree}{$p<0.001$}
\newcommand{\PvalBSixQnaPropGenOnevsGenThree}{$p<0.001$}
\newcommand{\PvalBSixQnaDetGenOnevsGenThree}{$p<0.001$}
\newcommand{\PvalBSixQnaRecGenOnevsGenThree}{$p<0.001$}
\newcommand{\PvalBSixQnaAccGenOnevsJudgeOnly}{$p<0.001$}
\newcommand{\PvalBSixQnaPropGenOnevsJudgeOnly}{$p<0.001$}
\newcommand{\PvalBSixQnaDetGenOnevsJudgeOnly}{$p<0.001$}
\newcommand{\PvalBSixQnaRecGenOnevsJudgeOnly}{$p<0.001$}
\newcommand{\PvalBSixQnaAccGenThreevsJudgeOnly}{$p=1.00$}
\newcommand{\PvalBSixQnaPropGenThreevsJudgeOnly}{$p=1.00$}
\newcommand{\PvalBSixQnaDetGenThreevsJudgeOnly}{$p<0.001$}
\newcommand{\PvalBSixQnaRecGenThreevsJudgeOnly}{$p<0.001$}
\newcommand{\PvalBSixCodeAccFlatThreevsBpd}{$p=0.503$}
\newcommand{\PvalBSixCodePropFlatThreevsBpd}{$p=0.039$}
\newcommand{\PvalBSixCodeDetFlatThreevsBpd}{$p<0.001$}
\newcommand{\PvalBSixCodeRecFlatThreevsBpd}{$p<0.001$}
\newcommand{\PvalBSixCodeAccGenOnevsBpd}{$p<0.001$}
\newcommand{\PvalBSixCodePropGenOnevsBpd}{$p<0.001$}
\newcommand{\PvalBSixCodeDetGenOnevsBpd}{$p<0.001$}
\newcommand{\PvalBSixCodeRecGenOnevsBpd}{$p<0.001$}
\newcommand{\PvalBSixCodeAccGenThreevsBpd}{$p=0.022$}
\newcommand{\PvalBSixCodePropGenThreevsBpd}{$p=1.00$}
\newcommand{\PvalBSixCodeDetGenThreevsBpd}{$p<0.001$}
\newcommand{\PvalBSixCodeRecGenThreevsBpd}{$p<0.001$}
\newcommand{\PvalBSixCodeAccJudgeOnlyvsBpd}{$p=0.003$}
\newcommand{\PvalBSixCodePropJudgeOnlyvsBpd}{$p=1.00$}
\newcommand{\PvalBSixCodeDetJudgeOnlyvsBpd}{$p=0.607$}
\newcommand{\PvalBSixCodeRecJudgeOnlyvsBpd}{$p=0.607$}
\newcommand{\PvalBSixCodeAccGenOnevsFlatThree}{$p<0.001$}
\newcommand{\PvalBSixCodePropGenOnevsFlatThree}{$p<0.001$}
\newcommand{\PvalBSixCodeDetGenOnevsFlatThree}{$p=1.00$}
\newcommand{\PvalBSixCodeRecGenOnevsFlatThree}{$p=1.00$}
\newcommand{\PvalBSixCodeAccGenThreevsFlatThree}{$p=0.180$}
\newcommand{\PvalBSixCodePropGenThreevsFlatThree}{$p=0.070$}
\newcommand{\PvalBSixCodeDetGenThreevsFlatThree}{$p=1.00$}
\newcommand{\PvalBSixCodeRecGenThreevsFlatThree}{$p=1.00$}
\newcommand{\PvalBSixCodeAccFlatThreevsJudgeOnly}{$p=0.039$}
\newcommand{\PvalBSixCodePropFlatThreevsJudgeOnly}{$p=0.039$}
\newcommand{\PvalBSixCodeDetFlatThreevsJudgeOnly}{$p<0.001$}
\newcommand{\PvalBSixCodeRecFlatThreevsJudgeOnly}{$p<0.001$}
\newcommand{\PvalBSixCodeAccGenOnevsGenThree}{$p<0.001$}
\newcommand{\PvalBSixCodePropGenOnevsGenThree}{$p<0.001$}
\newcommand{\PvalBSixCodeDetGenOnevsGenThree}{$p=1.00$}
\newcommand{\PvalBSixCodeRecGenOnevsGenThree}{$p=1.00$}
\newcommand{\PvalBSixCodeAccGenOnevsJudgeOnly}{$p<0.001$}
\newcommand{\PvalBSixCodePropGenOnevsJudgeOnly}{$p<0.001$}
\newcommand{\PvalBSixCodeDetGenOnevsJudgeOnly}{$p<0.001$}
\newcommand{\PvalBSixCodeRecGenOnevsJudgeOnly}{$p<0.001$}
\newcommand{\PvalBSixCodeAccGenThreevsJudgeOnly}{$p=0.625$}
\newcommand{\PvalBSixCodePropGenThreevsJudgeOnly}{$p=1.00$}
\newcommand{\PvalBSixCodeDetGenThreevsJudgeOnly}{$p<0.001$}
\newcommand{\PvalBSixCodeRecGenThreevsJudgeOnly}{$p<0.001$}
\newcommand{\PvalBSevenQnaAccFlatThreevsBpd}{$p=1.00$}
\newcommand{\PvalMainQnaAccFlatThreevsBpd}{$p=1.00$}
\newcommand{\PvalBSevenQnaPropFlatThreevsBpd}{$p=1.00$}
\newcommand{\PvalMainQnaPropFlatThreevsBpd}{$p=1.00$}
\newcommand{\PvalBSevenQnaDetFlatThreevsBpd}{$p=0.125$}
\newcommand{\PvalMainQnaDetFlatThreevsBpd}{$p=0.125$}
\newcommand{\PvalBSevenQnaRecFlatThreevsBpd}{$p=0.125$}
\newcommand{\PvalMainQnaRecFlatThreevsBpd}{$p=0.125$}
\newcommand{\PvalBSevenQnaAccGenOnevsBpd}{$p<0.001$}
\newcommand{\PvalMainQnaAccGenOnevsBpd}{$p<0.001$}
\newcommand{\PvalBSevenQnaPropGenOnevsBpd}{$p<0.001$}
\newcommand{\PvalMainQnaPropGenOnevsBpd}{$p<0.001$}
\newcommand{\PvalBSevenQnaDetGenOnevsBpd}{$p<0.001$}
\newcommand{\PvalMainQnaDetGenOnevsBpd}{$p<0.001$}
\newcommand{\PvalBSevenQnaRecGenOnevsBpd}{$p<0.001$}
\newcommand{\PvalMainQnaRecGenOnevsBpd}{$p<0.001$}
\newcommand{\PvalBSevenQnaAccGenThreevsBpd}{$p=1.00$}
\newcommand{\PvalMainQnaAccGenThreevsBpd}{$p=1.00$}
\newcommand{\PvalBSevenQnaPropGenThreevsBpd}{$p=1.00$}
\newcommand{\PvalMainQnaPropGenThreevsBpd}{$p=1.00$}
\newcommand{\PvalBSevenQnaDetGenThreevsBpd}{$p=0.031$}
\newcommand{\PvalMainQnaDetGenThreevsBpd}{$p=0.031$}
\newcommand{\PvalBSevenQnaRecGenThreevsBpd}{$p=0.031$}
\newcommand{\PvalMainQnaRecGenThreevsBpd}{$p=0.031$}
\newcommand{\PvalBSevenQnaAccJudgeOnlyvsBpd}{$p=1.00$}
\newcommand{\PvalMainQnaAccJudgeOnlyvsBpd}{$p=1.00$}
\newcommand{\PvalBSevenQnaPropJudgeOnlyvsBpd}{$p=1.00$}
\newcommand{\PvalMainQnaPropJudgeOnlyvsBpd}{$p=1.00$}
\newcommand{\PvalBSevenQnaDetJudgeOnlyvsBpd}{$p=0.581$}
\newcommand{\PvalMainQnaDetJudgeOnlyvsBpd}{$p=0.581$}
\newcommand{\PvalBSevenQnaRecJudgeOnlyvsBpd}{$p=0.581$}
\newcommand{\PvalMainQnaRecJudgeOnlyvsBpd}{$p=0.581$}
\newcommand{\PvalBSevenQnaAccGenOnevsFlatThree}{$p<0.001$}
\newcommand{\PvalMainQnaAccGenOnevsFlatThree}{$p<0.001$}
\newcommand{\PvalBSevenQnaPropGenOnevsFlatThree}{$p<0.001$}
\newcommand{\PvalMainQnaPropGenOnevsFlatThree}{$p<0.001$}
\newcommand{\PvalBSevenQnaDetGenOnevsFlatThree}{$p<0.001$}
\newcommand{\PvalMainQnaDetGenOnevsFlatThree}{$p<0.001$}
\newcommand{\PvalBSevenQnaRecGenOnevsFlatThree}{$p<0.001$}
\newcommand{\PvalMainQnaRecGenOnevsFlatThree}{$p<0.001$}
\newcommand{\PvalBSevenQnaAccGenThreevsFlatThree}{$p=1.00$}
\newcommand{\PvalMainQnaAccGenThreevsFlatThree}{$p=1.00$}
\newcommand{\PvalBSevenQnaPropGenThreevsFlatThree}{$p=1.00$}
\newcommand{\PvalMainQnaPropGenThreevsFlatThree}{$p=1.00$}
\newcommand{\PvalBSevenQnaDetGenThreevsFlatThree}{$p=1.00$}
\newcommand{\PvalMainQnaDetGenThreevsFlatThree}{$p=1.00$}
\newcommand{\PvalBSevenQnaRecGenThreevsFlatThree}{$p=1.00$}
\newcommand{\PvalMainQnaRecGenThreevsFlatThree}{$p=1.00$}
\newcommand{\PvalBSevenQnaAccFlatThreevsJudgeOnly}{$p=1.00$}
\newcommand{\PvalMainQnaAccFlatThreevsJudgeOnly}{$p=1.00$}
\newcommand{\PvalBSevenQnaPropFlatThreevsJudgeOnly}{$p=1.00$}
\newcommand{\PvalMainQnaPropFlatThreevsJudgeOnly}{$p=1.00$}
\newcommand{\PvalBSevenQnaDetFlatThreevsJudgeOnly}{$p=0.021$}
\newcommand{\PvalMainQnaDetFlatThreevsJudgeOnly}{$p=0.021$}
\newcommand{\PvalBSevenQnaRecFlatThreevsJudgeOnly}{$p=0.021$}
\newcommand{\PvalMainQnaRecFlatThreevsJudgeOnly}{$p=0.021$}
\newcommand{\PvalBSevenQnaAccGenOnevsGenThree}{$p<0.001$}
\newcommand{\PvalMainQnaAccGenOnevsGenThree}{$p<0.001$}
\newcommand{\PvalBSevenQnaPropGenOnevsGenThree}{$p<0.001$}
\newcommand{\PvalMainQnaPropGenOnevsGenThree}{$p<0.001$}
\newcommand{\PvalBSevenQnaDetGenOnevsGenThree}{$p<0.001$}
\newcommand{\PvalMainQnaDetGenOnevsGenThree}{$p<0.001$}
\newcommand{\PvalBSevenQnaRecGenOnevsGenThree}{$p<0.001$}
\newcommand{\PvalMainQnaRecGenOnevsGenThree}{$p<0.001$}
\newcommand{\PvalBSevenQnaAccGenOnevsJudgeOnly}{$p<0.001$}
\newcommand{\PvalMainQnaAccGenOnevsJudgeOnly}{$p<0.001$}
\newcommand{\PvalBSevenQnaPropGenOnevsJudgeOnly}{$p<0.001$}
\newcommand{\PvalMainQnaPropGenOnevsJudgeOnly}{$p<0.001$}
\newcommand{\PvalBSevenQnaDetGenOnevsJudgeOnly}{$p<0.001$}
\newcommand{\PvalMainQnaDetGenOnevsJudgeOnly}{$p<0.001$}
\newcommand{\PvalBSevenQnaRecGenOnevsJudgeOnly}{$p<0.001$}
\newcommand{\PvalMainQnaRecGenOnevsJudgeOnly}{$p<0.001$}
\newcommand{\PvalBSevenQnaAccGenThreevsJudgeOnly}{$p=1.00$}
\newcommand{\PvalMainQnaAccGenThreevsJudgeOnly}{$p=1.00$}
\newcommand{\PvalBSevenQnaPropGenThreevsJudgeOnly}{$p=1.00$}
\newcommand{\PvalMainQnaPropGenThreevsJudgeOnly}{$p=1.00$}
\newcommand{\PvalBSevenQnaDetGenThreevsJudgeOnly}{$p=0.004$}
\newcommand{\PvalMainQnaDetGenThreevsJudgeOnly}{$p=0.004$}
\newcommand{\PvalBSevenQnaRecGenThreevsJudgeOnly}{$p=0.004$}
\newcommand{\PvalMainQnaRecGenThreevsJudgeOnly}{$p=0.004$}
\newcommand{\PvalBSevenCodeAccFlatThreevsBpd}{$p<0.001$}
\newcommand{\PvalMainCodeAccFlatThreevsBpd}{$p<0.001$}
\newcommand{\PvalBSevenCodePropFlatThreevsBpd}{$p<0.001$}
\newcommand{\PvalMainCodePropFlatThreevsBpd}{$p<0.001$}
\newcommand{\PvalBSevenCodeDetFlatThreevsBpd}{$p<0.001$}
\newcommand{\PvalMainCodeDetFlatThreevsBpd}{$p<0.001$}
\newcommand{\PvalBSevenCodeRecFlatThreevsBpd}{$p<0.001$}
\newcommand{\PvalMainCodeRecFlatThreevsBpd}{$p<0.001$}
\newcommand{\PvalBSevenCodeAccGenOnevsBpd}{$p<0.001$}
\newcommand{\PvalMainCodeAccGenOnevsBpd}{$p<0.001$}
\newcommand{\PvalBSevenCodePropGenOnevsBpd}{$p<0.001$}
\newcommand{\PvalMainCodePropGenOnevsBpd}{$p<0.001$}
\newcommand{\PvalBSevenCodeDetGenOnevsBpd}{$p<0.001$}
\newcommand{\PvalMainCodeDetGenOnevsBpd}{$p<0.001$}
\newcommand{\PvalBSevenCodeRecGenOnevsBpd}{$p<0.001$}
\newcommand{\PvalMainCodeRecGenOnevsBpd}{$p<0.001$}
\newcommand{\PvalBSevenCodeAccGenThreevsBpd}{$p<0.001$}
\newcommand{\PvalMainCodeAccGenThreevsBpd}{$p<0.001$}
\newcommand{\PvalBSevenCodePropGenThreevsBpd}{$p<0.001$}
\newcommand{\PvalMainCodePropGenThreevsBpd}{$p<0.001$}
\newcommand{\PvalBSevenCodeDetGenThreevsBpd}{$p<0.001$}
\newcommand{\PvalMainCodeDetGenThreevsBpd}{$p<0.001$}
\newcommand{\PvalBSevenCodeRecGenThreevsBpd}{$p<0.001$}
\newcommand{\PvalMainCodeRecGenThreevsBpd}{$p<0.001$}
\newcommand{\PvalBSevenCodeAccJudgeOnlyvsBpd}{$p=1.00$}
\newcommand{\PvalMainCodeAccJudgeOnlyvsBpd}{$p=1.00$}
\newcommand{\PvalBSevenCodePropJudgeOnlyvsBpd}{$p=1.00$}
\newcommand{\PvalMainCodePropJudgeOnlyvsBpd}{$p=1.00$}
\newcommand{\PvalBSevenCodeDetJudgeOnlyvsBpd}{$p=0.481$}
\newcommand{\PvalMainCodeDetJudgeOnlyvsBpd}{$p=0.481$}
\newcommand{\PvalBSevenCodeRecJudgeOnlyvsBpd}{$p=0.481$}
\newcommand{\PvalMainCodeRecJudgeOnlyvsBpd}{$p=0.481$}
\newcommand{\PvalBSevenCodeAccGenOnevsFlatThree}{$p=0.035$}
\newcommand{\PvalMainCodeAccGenOnevsFlatThree}{$p=0.035$}
\newcommand{\PvalBSevenCodePropGenOnevsFlatThree}{$p=0.020$}
\newcommand{\PvalMainCodePropGenOnevsFlatThree}{$p=0.020$}
\newcommand{\PvalBSevenCodeDetGenOnevsFlatThree}{$p=1.00$}
\newcommand{\PvalMainCodeDetGenOnevsFlatThree}{$p=1.00$}
\newcommand{\PvalBSevenCodeRecGenOnevsFlatThree}{$p=1.00$}
\newcommand{\PvalMainCodeRecGenOnevsFlatThree}{$p=1.00$}
\newcommand{\PvalBSevenCodeAccGenThreevsFlatThree}{$p=0.064$}
\newcommand{\PvalMainCodeAccGenThreevsFlatThree}{$p=0.064$}
\newcommand{\PvalBSevenCodePropGenThreevsFlatThree}{$p=0.019$}
\newcommand{\PvalMainCodePropGenThreevsFlatThree}{$p=0.019$}
\newcommand{\PvalBSevenCodeDetGenThreevsFlatThree}{$p=1.00$}
\newcommand{\PvalMainCodeDetGenThreevsFlatThree}{$p=1.00$}
\newcommand{\PvalBSevenCodeRecGenThreevsFlatThree}{$p=1.00$}
\newcommand{\PvalMainCodeRecGenThreevsFlatThree}{$p=1.00$}
\newcommand{\PvalBSevenCodeAccFlatThreevsJudgeOnly}{$p<0.001$}
\newcommand{\PvalMainCodeAccFlatThreevsJudgeOnly}{$p<0.001$}
\newcommand{\PvalBSevenCodePropFlatThreevsJudgeOnly}{$p<0.001$}
\newcommand{\PvalMainCodePropFlatThreevsJudgeOnly}{$p<0.001$}
\newcommand{\PvalBSevenCodeDetFlatThreevsJudgeOnly}{$p<0.001$}
\newcommand{\PvalMainCodeDetFlatThreevsJudgeOnly}{$p<0.001$}
\newcommand{\PvalBSevenCodeRecFlatThreevsJudgeOnly}{$p<0.001$}
\newcommand{\PvalMainCodeRecFlatThreevsJudgeOnly}{$p<0.001$}
\newcommand{\PvalBSevenCodeAccGenOnevsGenThree}{$p<0.001$}
\newcommand{\PvalMainCodeAccGenOnevsGenThree}{$p<0.001$}
\newcommand{\PvalBSevenCodePropGenOnevsGenThree}{$p<0.001$}
\newcommand{\PvalMainCodePropGenOnevsGenThree}{$p<0.001$}
\newcommand{\PvalBSevenCodeDetGenOnevsGenThree}{$p=0.500$}
\newcommand{\PvalMainCodeDetGenOnevsGenThree}{$p=0.500$}
\newcommand{\PvalBSevenCodeRecGenOnevsGenThree}{$p=0.500$}
\newcommand{\PvalMainCodeRecGenOnevsGenThree}{$p=0.500$}
\newcommand{\PvalBSevenCodeAccGenOnevsJudgeOnly}{$p<0.001$}
\newcommand{\PvalMainCodeAccGenOnevsJudgeOnly}{$p<0.001$}
\newcommand{\PvalBSevenCodePropGenOnevsJudgeOnly}{$p<0.001$}
\newcommand{\PvalMainCodePropGenOnevsJudgeOnly}{$p<0.001$}
\newcommand{\PvalBSevenCodeDetGenOnevsJudgeOnly}{$p<0.001$}
\newcommand{\PvalMainCodeDetGenOnevsJudgeOnly}{$p<0.001$}
\newcommand{\PvalBSevenCodeRecGenOnevsJudgeOnly}{$p<0.001$}
\newcommand{\PvalMainCodeRecGenOnevsJudgeOnly}{$p<0.001$}
\newcommand{\PvalBSevenCodeAccGenThreevsJudgeOnly}{$p<0.001$}
\newcommand{\PvalMainCodeAccGenThreevsJudgeOnly}{$p<0.001$}
\newcommand{\PvalBSevenCodePropGenThreevsJudgeOnly}{$p<0.001$}
\newcommand{\PvalMainCodePropGenThreevsJudgeOnly}{$p<0.001$}
\newcommand{\PvalBSevenCodeDetGenThreevsJudgeOnly}{$p<0.001$}
\newcommand{\PvalMainCodeDetGenThreevsJudgeOnly}{$p<0.001$}
\newcommand{\PvalBSevenCodeRecGenThreevsJudgeOnly}{$p<0.001$}
\newcommand{\PvalMainCodeRecGenThreevsJudgeOnly}{$p<0.001$}
\newcommand{\PvalBEightQnaAccFlatThreevsBpd}{$p=1.00$}
\newcommand{\PvalBEightQnaPropFlatThreevsBpd}{$p=1.00$}
\newcommand{\PvalBEightQnaDetFlatThreevsBpd}{$p=0.791$}
\newcommand{\PvalBEightQnaRecFlatThreevsBpd}{$p=0.791$}
\newcommand{\PvalBEightQnaAccGenOnevsBpd}{$p<0.001$}
\newcommand{\PvalBEightQnaPropGenOnevsBpd}{$p<0.001$}
\newcommand{\PvalBEightQnaDetGenOnevsBpd}{$p=0.002$}
\newcommand{\PvalBEightQnaRecGenOnevsBpd}{$p=0.002$}
\newcommand{\PvalBEightQnaAccGenThreevsBpd}{$p=1.00$}
\newcommand{\PvalBEightQnaPropGenThreevsBpd}{$p=1.00$}
\newcommand{\PvalBEightQnaDetGenThreevsBpd}{$p<0.001$}
\newcommand{\PvalBEightQnaRecGenThreevsBpd}{$p<0.001$}
\newcommand{\PvalBEightQnaAccJudgeOnlyvsBpd}{$p=1.00$}
\newcommand{\PvalBEightQnaPropJudgeOnlyvsBpd}{$p=1.00$}
\newcommand{\PvalBEightQnaDetJudgeOnlyvsBpd}{$p=1.00$}
\newcommand{\PvalBEightQnaRecJudgeOnlyvsBpd}{$p=1.00$}
\newcommand{\PvalBEightQnaAccGenOnevsFlatThree}{$p<0.001$}
\newcommand{\PvalBEightQnaPropGenOnevsFlatThree}{$p<0.001$}
\newcommand{\PvalBEightQnaDetGenOnevsFlatThree}{$p<0.001$}
\newcommand{\PvalBEightQnaRecGenOnevsFlatThree}{$p<0.001$}
\newcommand{\PvalBEightQnaAccGenThreevsFlatThree}{$p=1.00$}
\newcommand{\PvalBEightQnaPropGenThreevsFlatThree}{$p=1.00$}
\newcommand{\PvalBEightQnaDetGenThreevsFlatThree}{$p<0.001$}
\newcommand{\PvalBEightQnaRecGenThreevsFlatThree}{$p<0.001$}
\newcommand{\PvalBEightQnaAccFlatThreevsJudgeOnly}{$p=1.00$}
\newcommand{\PvalBEightQnaPropFlatThreevsJudgeOnly}{$p=1.00$}
\newcommand{\PvalBEightQnaDetFlatThreevsJudgeOnly}{$p=1.00$}
\newcommand{\PvalBEightQnaRecFlatThreevsJudgeOnly}{$p=1.00$}
\newcommand{\PvalBEightQnaAccGenOnevsGenThree}{$p<0.001$}
\newcommand{\PvalBEightQnaPropGenOnevsGenThree}{$p<0.001$}
\newcommand{\PvalBEightQnaDetGenOnevsGenThree}{$p<0.001$}
\newcommand{\PvalBEightQnaRecGenOnevsGenThree}{$p<0.001$}
\newcommand{\PvalBEightQnaAccGenOnevsJudgeOnly}{$p<0.001$}
\newcommand{\PvalBEightQnaPropGenOnevsJudgeOnly}{$p<0.001$}
\newcommand{\PvalBEightQnaDetGenOnevsJudgeOnly}{$p<0.001$}
\newcommand{\PvalBEightQnaRecGenOnevsJudgeOnly}{$p<0.001$}
\newcommand{\PvalBEightQnaAccGenThreevsJudgeOnly}{$p=1.00$}
\newcommand{\PvalBEightQnaPropGenThreevsJudgeOnly}{$p=1.00$}
\newcommand{\PvalBEightQnaDetGenThreevsJudgeOnly}{$p<0.001$}
\newcommand{\PvalBEightQnaRecGenThreevsJudgeOnly}{$p<0.001$}
\newcommand{\PvalBEightCodeAccFlatThreevsBpd}{$p<0.001$}
\newcommand{\PvalBEightCodePropFlatThreevsBpd}{$p<0.001$}
\newcommand{\PvalBEightCodeDetFlatThreevsBpd}{$p<0.001$}
\newcommand{\PvalBEightCodeRecFlatThreevsBpd}{$p<0.001$}
\newcommand{\PvalBEightCodeAccGenOnevsBpd}{$p<0.001$}
\newcommand{\PvalBEightCodePropGenOnevsBpd}{$p<0.001$}
\newcommand{\PvalBEightCodeDetGenOnevsBpd}{$p<0.001$}
\newcommand{\PvalBEightCodeRecGenOnevsBpd}{$p<0.001$}
\newcommand{\PvalBEightCodeAccGenThreevsBpd}{$p<0.001$}
\newcommand{\PvalBEightCodePropGenThreevsBpd}{$p<0.001$}
\newcommand{\PvalBEightCodeDetGenThreevsBpd}{$p<0.001$}
\newcommand{\PvalBEightCodeRecGenThreevsBpd}{$p<0.001$}
\newcommand{\PvalBEightCodeAccJudgeOnlyvsBpd}{$p=1.00$}
\newcommand{\PvalBEightCodePropJudgeOnlyvsBpd}{$p=1.00$}
\newcommand{\PvalBEightCodeDetJudgeOnlyvsBpd}{$p=1.00$}
\newcommand{\PvalBEightCodeRecJudgeOnlyvsBpd}{$p=1.00$}
\newcommand{\PvalBEightCodeAccGenOnevsFlatThree}{$p<0.001$}
\newcommand{\PvalBEightCodePropGenOnevsFlatThree}{$p=0.005$}
\newcommand{\PvalBEightCodeDetGenOnevsFlatThree}{$p=1.00$}
\newcommand{\PvalBEightCodeRecGenOnevsFlatThree}{$p=1.00$}
\newcommand{\PvalBEightCodeAccGenThreevsFlatThree}{$p=0.815$}
\newcommand{\PvalBEightCodePropGenThreevsFlatThree}{$p=0.815$}
\newcommand{\PvalBEightCodeDetGenThreevsFlatThree}{$p=1.00$}
\newcommand{\PvalBEightCodeRecGenThreevsFlatThree}{$p=1.00$}
\newcommand{\PvalBEightCodeAccFlatThreevsJudgeOnly}{$p<0.001$}
\newcommand{\PvalBEightCodePropFlatThreevsJudgeOnly}{$p<0.001$}
\newcommand{\PvalBEightCodeDetFlatThreevsJudgeOnly}{$p<0.001$}
\newcommand{\PvalBEightCodeRecFlatThreevsJudgeOnly}{$p<0.001$}
\newcommand{\PvalBEightCodeAccGenOnevsGenThree}{$p<0.001$}
\newcommand{\PvalBEightCodePropGenOnevsGenThree}{$p=0.001$}
\newcommand{\PvalBEightCodeDetGenOnevsGenThree}{$p=1.00$}
\newcommand{\PvalBEightCodeRecGenOnevsGenThree}{$p=1.00$}
\newcommand{\PvalBEightCodeAccGenOnevsJudgeOnly}{$p<0.001$}
\newcommand{\PvalBEightCodePropGenOnevsJudgeOnly}{$p<0.001$}
\newcommand{\PvalBEightCodeDetGenOnevsJudgeOnly}{$p<0.001$}
\newcommand{\PvalBEightCodeRecGenOnevsJudgeOnly}{$p<0.001$}
\newcommand{\PvalBEightCodeAccGenThreevsJudgeOnly}{$p<0.001$}
\newcommand{\PvalBEightCodePropGenThreevsJudgeOnly}{$p<0.001$}
\newcommand{\PvalBEightCodeDetGenThreevsJudgeOnly}{$p<0.001$}
\newcommand{\PvalBEightCodeRecGenThreevsJudgeOnly}{$p<0.001$}
\newcommand{\RatePosBSevenBpdQnaAccZero}{1.00}
\newcommand{\PvalPosBSevenBpdQnaAcc}{$p=1.00$}
\newcommand{\RatePosBSevenBpdQnaPropZero}{0.00}
\newcommand{\PvalPosBSevenBpdQnaProp}{$p=1.00$}
\newcommand{\RatePosBSevenBpdQnaDetZero}{0.94}
\newcommand{\PvalPosBSevenBpdQnaDet}{$p=0.688$}
\newcommand{\RatePosBSevenBpdQnaRecZero}{0.94}
\newcommand{\PvalPosBSevenBpdQnaRec}{$p=0.696$}
\newcommand{\RatePosBSevenBpdQnaAccOne}{1.00}
\newcommand{\RatePosBSevenBpdQnaPropOne}{0.00}
\newcommand{\RatePosBSevenBpdQnaDetOne}{0.91}
\newcommand{\RatePosBSevenBpdQnaRecOne}{0.91}
\newcommand{\RatePosBSevenBpdQnaAccTwo}{1.00}
\newcommand{\RatePosBSevenBpdQnaPropTwo}{0.00}
\newcommand{\RatePosBSevenBpdQnaDetTwo}{0.97}
\newcommand{\RatePosBSevenBpdQnaRecTwo}{0.97}
\newcommand{\RatePosBSevenBpdCodeAccZero}{0.97}
\newcommand{\PvalPosBSevenBpdCodeAcc}{$p=1.00$}
\newcommand{\RatePosBSevenBpdCodePropZero}{0.03}
\newcommand{\PvalPosBSevenBpdCodeProp}{$p=1.00$}
\newcommand{\RatePosBSevenBpdCodeDetZero}{0.62}
\newcommand{\PvalPosBSevenBpdCodeDet}{$p=0.240$}
\newcommand{\RatePosBSevenBpdCodeRecZero}{0.62}
\newcommand{\PvalPosBSevenBpdCodeRec}{$p=0.247$}
\newcommand{\RatePosBSevenBpdCodeAccOne}{1.00}
\newcommand{\RatePosBSevenBpdCodePropOne}{0.00}
\newcommand{\RatePosBSevenBpdCodeDetOne}{0.61}
\newcommand{\RatePosBSevenBpdCodeRecOne}{0.61}
\newcommand{\RatePosBSevenBpdCodeAccTwo}{1.00}
\newcommand{\RatePosBSevenBpdCodePropTwo}{0.00}
\newcommand{\RatePosBSevenBpdCodeDetTwo}{0.79}
\newcommand{\RatePosBSevenBpdCodeRecTwo}{0.79}
\newcommand{\RatePosBSevenFlatThreeQnaAccZero}{1.00}
\newcommand{\PvalPosBSevenFlatThreeQnaAcc}{$p=1.00$}
\newcommand{\RatePosBSevenFlatThreeQnaPropZero}{0.00}
\newcommand{\PvalPosBSevenFlatThreeQnaProp}{$p=1.00$}
\newcommand{\RatePosBSevenFlatThreeQnaDetZero}{0.97}
\newcommand{\PvalPosBSevenFlatThreeQnaDet}{$p=1.00$}
\newcommand{\RatePosBSevenFlatThreeQnaRecZero}{0.97}
\newcommand{\PvalPosBSevenFlatThreeQnaRec}{$p=1.00$}
\newcommand{\RatePosBSevenFlatThreeQnaAccOne}{1.00}
\newcommand{\RatePosBSevenFlatThreeQnaPropOne}{0.00}
\newcommand{\RatePosBSevenFlatThreeQnaDetOne}{1.00}
\newcommand{\RatePosBSevenFlatThreeQnaRecOne}{1.00}
\newcommand{\RatePosBSevenFlatThreeQnaAccTwo}{1.00}
\newcommand{\RatePosBSevenFlatThreeQnaPropTwo}{0.00}
\newcommand{\RatePosBSevenFlatThreeQnaDetTwo}{1.00}
\newcommand{\RatePosBSevenFlatThreeQnaRecTwo}{1.00}
\newcommand{\RatePosBSevenFlatThreeCodeAccZero}{0.88}
\newcommand{\PvalPosBSevenFlatThreeCodeAcc}{$p<0.001$}
\newcommand{\RatePosBSevenFlatThreeCodePropZero}{0.12}
\newcommand{\PvalPosBSevenFlatThreeCodeProp}{$p<0.001$}
\newcommand{\RatePosBSevenFlatThreeCodeDetZero}{0.00}
\newcommand{\PvalPosBSevenFlatThreeCodeDet}{$p=0.654$}
\newcommand{\RatePosBSevenFlatThreeCodeRecZero}{0.00}
\newcommand{\PvalPosBSevenFlatThreeCodeRec}{$p=0.658$}
\newcommand{\RatePosBSevenFlatThreeCodeAccOne}{0.91}
\newcommand{\RatePosBSevenFlatThreeCodePropOne}{0.09}
\newcommand{\RatePosBSevenFlatThreeCodeDetOne}{0.00}
\newcommand{\RatePosBSevenFlatThreeCodeRecOne}{0.00}
\newcommand{\RatePosBSevenFlatThreeCodeAccTwo}{0.36}
\newcommand{\RatePosBSevenFlatThreeCodePropTwo}{0.58}
\newcommand{\RatePosBSevenFlatThreeCodeDetTwo}{0.03}
\newcommand{\RatePosBSevenFlatThreeCodeRecTwo}{0.03}
\newcommand{\RatePosBSevenGenOneQnaAccZero}{0.00}
\newcommand{\PvalPosBSevenGenOneQnaAcc}{$p=1.00$}
\newcommand{\RatePosBSevenGenOneQnaPropZero}{1.00}
\newcommand{\PvalPosBSevenGenOneQnaProp}{$p=1.00$}
\newcommand{\RatePosBSevenGenOneQnaDetZero}{0.00}
\newcommand{\PvalPosBSevenGenOneQnaDet}{$p=1.00$}
\newcommand{\RatePosBSevenGenOneQnaRecZero}{0.00}
\newcommand{\PvalPosBSevenGenOneQnaRec}{$p=1.00$}
\newcommand{\RatePosBSevenGenOneQnaAccOne}{0.00}
\newcommand{\RatePosBSevenGenOneQnaPropOne}{1.00}
\newcommand{\RatePosBSevenGenOneQnaDetOne}{0.00}
\newcommand{\RatePosBSevenGenOneQnaRecOne}{0.00}
\newcommand{\RatePosBSevenGenOneQnaAccTwo}{0.00}
\newcommand{\RatePosBSevenGenOneQnaPropTwo}{1.00}
\newcommand{\RatePosBSevenGenOneQnaDetTwo}{0.00}
\newcommand{\RatePosBSevenGenOneQnaRecTwo}{0.00}
\newcommand{\RatePosBSevenGenOneCodeAccZero}{0.53}
\newcommand{\PvalPosBSevenGenOneCodeAcc}{$p=0.348$}
\newcommand{\RatePosBSevenGenOneCodePropZero}{0.44}
\newcommand{\PvalPosBSevenGenOneCodeProp}{$p=0.407$}
\newcommand{\RatePosBSevenGenOneCodeDetZero}{0.00}
\newcommand{\PvalPosBSevenGenOneCodeDet}{$p=1.00$}
\newcommand{\RatePosBSevenGenOneCodeRecZero}{0.00}
\newcommand{\PvalPosBSevenGenOneCodeRec}{$p=1.00$}
\newcommand{\RatePosBSevenGenOneCodeAccOne}{0.70}
\newcommand{\RatePosBSevenGenOneCodePropOne}{0.30}
\newcommand{\RatePosBSevenGenOneCodeDetOne}{0.00}
\newcommand{\RatePosBSevenGenOneCodeRecOne}{0.00}
\newcommand{\RatePosBSevenGenOneCodeAccTwo}{0.55}
\newcommand{\RatePosBSevenGenOneCodePropTwo}{0.45}
\newcommand{\RatePosBSevenGenOneCodeDetTwo}{0.00}
\newcommand{\RatePosBSevenGenOneCodeRecTwo}{0.00}
\newcommand{\RatePosBSevenGenThreeQnaAccZero}{1.00}
\newcommand{\PvalPosBSevenGenThreeQnaAcc}{$p=1.00$}
\newcommand{\RatePosBSevenGenThreeQnaPropZero}{0.00}
\newcommand{\PvalPosBSevenGenThreeQnaProp}{$p=1.00$}
\newcommand{\RatePosBSevenGenThreeQnaDetZero}{1.00}
\newcommand{\PvalPosBSevenGenThreeQnaDet}{$p=1.00$}
\newcommand{\RatePosBSevenGenThreeQnaRecZero}{1.00}
\newcommand{\PvalPosBSevenGenThreeQnaRec}{$p=1.00$}
\newcommand{\RatePosBSevenGenThreeQnaAccOne}{1.00}
\newcommand{\RatePosBSevenGenThreeQnaPropOne}{0.00}
\newcommand{\RatePosBSevenGenThreeQnaDetOne}{1.00}
\newcommand{\RatePosBSevenGenThreeQnaRecOne}{1.00}
\newcommand{\RatePosBSevenGenThreeQnaAccTwo}{1.00}
\newcommand{\RatePosBSevenGenThreeQnaPropTwo}{0.00}
\newcommand{\RatePosBSevenGenThreeQnaDetTwo}{1.00}
\newcommand{\RatePosBSevenGenThreeQnaRecTwo}{1.00}
\newcommand{\RatePosBSevenGenThreeCodeAccZero}{0.91}
\newcommand{\PvalPosBSevenGenThreeCodeAcc}{$p<0.001$}
\newcommand{\RatePosBSevenGenThreeCodePropZero}{0.06}
\newcommand{\PvalPosBSevenGenThreeCodeProp}{$p<0.001$}
\newcommand{\RatePosBSevenGenThreeCodeDetZero}{0.00}
\newcommand{\PvalPosBSevenGenThreeCodeDet}{$p=0.213$}
\newcommand{\RatePosBSevenGenThreeCodeRecZero}{0.00}
\newcommand{\PvalPosBSevenGenThreeCodeRec}{$p=0.207$}
\newcommand{\RatePosBSevenGenThreeCodeAccOne}{1.00}
\newcommand{\RatePosBSevenGenThreeCodePropOne}{0.00}
\newcommand{\RatePosBSevenGenThreeCodeDetOne}{0.06}
\newcommand{\RatePosBSevenGenThreeCodeRecOne}{0.06}
\newcommand{\RatePosBSevenGenThreeCodeAccTwo}{0.52}
\newcommand{\RatePosBSevenGenThreeCodePropTwo}{0.39}
\newcommand{\RatePosBSevenGenThreeCodeDetTwo}{0.00}
\newcommand{\RatePosBSevenGenThreeCodeRecTwo}{0.00}
\newcommand{\RatePosBSevenJudgeOnlyQnaAccZero}{1.00}
\newcommand{\PvalPosBSevenJudgeOnlyQnaAcc}{$p=1.00$}
\newcommand{\RatePosBSevenJudgeOnlyQnaPropZero}{0.00}
\newcommand{\PvalPosBSevenJudgeOnlyQnaProp}{$p=1.00$}
\newcommand{\RatePosBSevenJudgeOnlyQnaDetZero}{0.82}
\newcommand{\PvalPosBSevenJudgeOnlyQnaDet}{$p=0.139$}
\newcommand{\RatePosBSevenJudgeOnlyQnaRecZero}{0.82}
\newcommand{\PvalPosBSevenJudgeOnlyQnaRec}{$p=0.140$}
\newcommand{\RatePosBSevenJudgeOnlyQnaAccOne}{1.00}
\newcommand{\RatePosBSevenJudgeOnlyQnaPropOne}{0.00}
\newcommand{\RatePosBSevenJudgeOnlyQnaDetOne}{0.94}
\newcommand{\RatePosBSevenJudgeOnlyQnaRecOne}{0.94}
\newcommand{\RatePosBSevenJudgeOnlyQnaAccTwo}{1.00}
\newcommand{\RatePosBSevenJudgeOnlyQnaPropTwo}{0.00}
\newcommand{\RatePosBSevenJudgeOnlyQnaDetTwo}{0.97}
\newcommand{\RatePosBSevenJudgeOnlyQnaRecTwo}{0.97}
\newcommand{\RatePosBSevenJudgeOnlyCodeAccZero}{1.00}
\newcommand{\PvalPosBSevenJudgeOnlyCodeAcc}{$p=1.00$}
\newcommand{\RatePosBSevenJudgeOnlyCodePropZero}{0.00}
\newcommand{\PvalPosBSevenJudgeOnlyCodeProp}{$p=1.00$}
\newcommand{\RatePosBSevenJudgeOnlyCodeDetZero}{0.65}
\newcommand{\PvalPosBSevenJudgeOnlyCodeDet}{$p=0.760$}
\newcommand{\RatePosBSevenJudgeOnlyCodeRecZero}{0.65}
\newcommand{\PvalPosBSevenJudgeOnlyCodeRec}{$p=0.757$}
\newcommand{\RatePosBSevenJudgeOnlyCodeAccOne}{1.00}
\newcommand{\RatePosBSevenJudgeOnlyCodePropOne}{0.00}
\newcommand{\RatePosBSevenJudgeOnlyCodeDetOne}{0.67}
\newcommand{\RatePosBSevenJudgeOnlyCodeRecOne}{0.67}
\newcommand{\RatePosBSevenJudgeOnlyCodeAccTwo}{1.00}
\newcommand{\RatePosBSevenJudgeOnlyCodePropTwo}{0.00}
\newcommand{\RatePosBSevenJudgeOnlyCodeDetTwo}{0.58}
\newcommand{\RatePosBSevenJudgeOnlyCodeRecTwo}{0.58}
\newcommand{\RatePosBEightBpdQnaAccZero}{1.00}
\newcommand{\PvalPosBEightBpdQnaAcc}{$p=1.00$}
\newcommand{\RatePosBEightBpdQnaPropZero}{0.00}
\newcommand{\PvalPosBEightBpdQnaProp}{$p=1.00$}
\newcommand{\RatePosBEightBpdQnaDetZero}{0.18}
\newcommand{\PvalPosBEightBpdQnaDet}{$p=0.146$}
\newcommand{\RatePosBEightBpdQnaRecZero}{0.18}
\newcommand{\PvalPosBEightBpdQnaRec}{$p=0.146$}
\newcommand{\RatePosBEightBpdQnaAccOne}{1.00}
\newcommand{\RatePosBEightBpdQnaPropOne}{0.00}
\newcommand{\RatePosBEightBpdQnaDetOne}{0.03}
\newcommand{\RatePosBEightBpdQnaRecOne}{0.03}
\newcommand{\RatePosBEightBpdQnaAccTwo}{1.00}
\newcommand{\RatePosBEightBpdQnaPropTwo}{0.00}
\newcommand{\RatePosBEightBpdQnaDetTwo}{0.09}
\newcommand{\RatePosBEightBpdQnaRecTwo}{0.09}
\newcommand{\RatePosBEightBpdCodeAccZero}{0.94}
\newcommand{\PvalPosBEightBpdCodeAcc}{$p=0.325$}
\newcommand{\RatePosBEightBpdCodePropZero}{0.06}
\newcommand{\PvalPosBEightBpdCodeProp}{$p=0.329$}
\newcommand{\RatePosBEightBpdCodeDetZero}{0.68}
\newcommand{\PvalPosBEightBpdCodeDet}{$p=0.216$}
\newcommand{\RatePosBEightBpdCodeRecZero}{0.68}
\newcommand{\PvalPosBEightBpdCodeRec}{$p=0.203$}
\newcommand{\RatePosBEightBpdCodeAccOne}{1.00}
\newcommand{\RatePosBEightBpdCodePropOne}{0.00}
\newcommand{\RatePosBEightBpdCodeDetOne}{0.48}
\newcommand{\RatePosBEightBpdCodeRecOne}{0.48}
\newcommand{\RatePosBEightBpdCodeAccTwo}{1.00}
\newcommand{\RatePosBEightBpdCodePropTwo}{0.00}
\newcommand{\RatePosBEightBpdCodeDetTwo}{0.67}
\newcommand{\RatePosBEightBpdCodeRecTwo}{0.67}
\newcommand{\RatePosBEightFlatThreeQnaAccZero}{1.00}
\newcommand{\PvalPosBEightFlatThreeQnaAcc}{$p=1.00$}
\newcommand{\RatePosBEightFlatThreeQnaPropZero}{0.00}
\newcommand{\PvalPosBEightFlatThreeQnaProp}{$p=1.00$}
\newcommand{\RatePosBEightFlatThreeQnaDetZero}{0.26}
\newcommand{\PvalPosBEightFlatThreeQnaDet}{$p=0.006$}
\newcommand{\RatePosBEightFlatThreeQnaRecZero}{0.26}
\newcommand{\PvalPosBEightFlatThreeQnaRec}{$p=0.006$}
\newcommand{\RatePosBEightFlatThreeQnaAccOne}{1.00}
\newcommand{\RatePosBEightFlatThreeQnaPropOne}{0.00}
\newcommand{\RatePosBEightFlatThreeQnaDetOne}{0.06}
\newcommand{\RatePosBEightFlatThreeQnaRecOne}{0.06}
\newcommand{\RatePosBEightFlatThreeQnaAccTwo}{1.00}
\newcommand{\RatePosBEightFlatThreeQnaPropTwo}{0.00}
\newcommand{\RatePosBEightFlatThreeQnaDetTwo}{0.03}
\newcommand{\RatePosBEightFlatThreeQnaRecTwo}{0.03}
\newcommand{\RatePosBEightFlatThreeCodeAccZero}{0.88}
\newcommand{\PvalPosBEightFlatThreeCodeAcc}{$p<0.001$}
\newcommand{\RatePosBEightFlatThreeCodePropZero}{0.12}
\newcommand{\PvalPosBEightFlatThreeCodeProp}{$p<0.001$}
\newcommand{\RatePosBEightFlatThreeCodeDetZero}{0.03}
\newcommand{\PvalPosBEightFlatThreeCodeDet}{$p=1.00$}
\newcommand{\RatePosBEightFlatThreeCodeRecZero}{0.03}
\newcommand{\PvalPosBEightFlatThreeCodeRec}{$p=1.00$}
\newcommand{\RatePosBEightFlatThreeCodeAccOne}{0.91}
\newcommand{\RatePosBEightFlatThreeCodePropOne}{0.09}
\newcommand{\RatePosBEightFlatThreeCodeDetOne}{0.00}
\newcommand{\RatePosBEightFlatThreeCodeRecOne}{0.00}
\newcommand{\RatePosBEightFlatThreeCodeAccTwo}{0.30}
\newcommand{\RatePosBEightFlatThreeCodePropTwo}{0.67}
\newcommand{\RatePosBEightFlatThreeCodeDetTwo}{0.00}
\newcommand{\RatePosBEightFlatThreeCodeRecTwo}{0.00}
\newcommand{\RatePosBEightGenOneQnaAccZero}{0.06}
\newcommand{\PvalPosBEightGenOneQnaAcc}{$p=1.00$}
\newcommand{\RatePosBEightGenOneQnaPropZero}{0.94}
\newcommand{\PvalPosBEightGenOneQnaProp}{$p=1.00$}
\newcommand{\RatePosBEightGenOneQnaDetZero}{0.00}
\newcommand{\PvalPosBEightGenOneQnaDet}{$p=1.00$}
\newcommand{\RatePosBEightGenOneQnaRecZero}{0.00}
\newcommand{\PvalPosBEightGenOneQnaRec}{$p=1.00$}
\newcommand{\RatePosBEightGenOneQnaAccOne}{0.06}
\newcommand{\RatePosBEightGenOneQnaPropOne}{0.94}
\newcommand{\RatePosBEightGenOneQnaDetOne}{0.00}
\newcommand{\RatePosBEightGenOneQnaRecOne}{0.00}
\newcommand{\RatePosBEightGenOneQnaAccTwo}{0.03}
\newcommand{\RatePosBEightGenOneQnaPropTwo}{0.97}
\newcommand{\RatePosBEightGenOneQnaDetTwo}{0.00}
\newcommand{\RatePosBEightGenOneQnaRecTwo}{0.00}
\newcommand{\RatePosBEightGenOneCodeAccZero}{0.50}
\newcommand{\PvalPosBEightGenOneCodeAcc}{$p=0.511$}
\newcommand{\RatePosBEightGenOneCodePropZero}{0.35}
\newcommand{\PvalPosBEightGenOneCodeProp}{$p=0.289$}
\newcommand{\RatePosBEightGenOneCodeDetZero}{0.00}
\newcommand{\PvalPosBEightGenOneCodeDet}{$p=1.00$}
\newcommand{\RatePosBEightGenOneCodeRecZero}{0.00}
\newcommand{\PvalPosBEightGenOneCodeRec}{$p=1.00$}
\newcommand{\RatePosBEightGenOneCodeAccOne}{0.45}
\newcommand{\RatePosBEightGenOneCodePropOne}{0.48}
\newcommand{\RatePosBEightGenOneCodeDetOne}{0.00}
\newcommand{\RatePosBEightGenOneCodeRecOne}{0.00}
\newcommand{\RatePosBEightGenOneCodeAccTwo}{0.36}
\newcommand{\RatePosBEightGenOneCodePropTwo}{0.55}
\newcommand{\RatePosBEightGenOneCodeDetTwo}{0.00}
\newcommand{\RatePosBEightGenOneCodeRecTwo}{0.00}
\newcommand{\RatePosBEightGenThreeQnaAccZero}{1.00}
\newcommand{\PvalPosBEightGenThreeQnaAcc}{$p=1.00$}
\newcommand{\RatePosBEightGenThreeQnaPropZero}{0.00}
\newcommand{\PvalPosBEightGenThreeQnaProp}{$p=1.00$}
\newcommand{\RatePosBEightGenThreeQnaDetZero}{0.94}
\newcommand{\PvalPosBEightGenThreeQnaDet}{$p=0.270$}
\newcommand{\RatePosBEightGenThreeQnaRecZero}{0.94}
\newcommand{\PvalPosBEightGenThreeQnaRec}{$p=0.281$}
\newcommand{\RatePosBEightGenThreeQnaAccOne}{1.00}
\newcommand{\RatePosBEightGenThreeQnaPropOne}{0.00}
\newcommand{\RatePosBEightGenThreeQnaDetOne}{1.00}
\newcommand{\RatePosBEightGenThreeQnaRecOne}{1.00}
\newcommand{\RatePosBEightGenThreeQnaAccTwo}{1.00}
\newcommand{\RatePosBEightGenThreeQnaPropTwo}{0.00}
\newcommand{\RatePosBEightGenThreeQnaDetTwo}{0.91}
\newcommand{\RatePosBEightGenThreeQnaRecTwo}{0.91}
\newcommand{\RatePosBEightGenThreeCodeAccZero}{0.94}
\newcommand{\PvalPosBEightGenThreeCodeAcc}{$p<0.001$}
\newcommand{\RatePosBEightGenThreeCodePropZero}{0.06}
\newcommand{\PvalPosBEightGenThreeCodeProp}{$p<0.001$}
\newcommand{\RatePosBEightGenThreeCodeDetZero}{0.00}
\newcommand{\PvalPosBEightGenThreeCodeDet}{$p=0.663$}
\newcommand{\RatePosBEightGenThreeCodeRecZero}{0.00}
\newcommand{\PvalPosBEightGenThreeCodeRec}{$p=0.661$}
\newcommand{\RatePosBEightGenThreeCodeAccOne}{0.82}
\newcommand{\RatePosBEightGenThreeCodePropOne}{0.18}
\newcommand{\RatePosBEightGenThreeCodeDetOne}{0.03}
\newcommand{\RatePosBEightGenThreeCodeRecOne}{0.03}
\newcommand{\RatePosBEightGenThreeCodeAccTwo}{0.39}
\newcommand{\RatePosBEightGenThreeCodePropTwo}{0.58}
\newcommand{\RatePosBEightGenThreeCodeDetTwo}{0.00}
\newcommand{\RatePosBEightGenThreeCodeRecTwo}{0.00}
\newcommand{\RatePosBEightJudgeOnlyQnaAccZero}{1.00}
\newcommand{\PvalPosBEightJudgeOnlyQnaAcc}{$p=1.00$}
\newcommand{\RatePosBEightJudgeOnlyQnaPropZero}{0.00}
\newcommand{\PvalPosBEightJudgeOnlyQnaProp}{$p=1.00$}
\newcommand{\RatePosBEightJudgeOnlyQnaDetZero}{0.21}
\newcommand{\PvalPosBEightJudgeOnlyQnaDet}{$p=0.077$}
\newcommand{\RatePosBEightJudgeOnlyQnaRecZero}{0.21}
\newcommand{\PvalPosBEightJudgeOnlyQnaRec}{$p=0.074$}
\newcommand{\RatePosBEightJudgeOnlyQnaAccOne}{1.00}
\newcommand{\RatePosBEightJudgeOnlyQnaPropOne}{0.00}
\newcommand{\RatePosBEightJudgeOnlyQnaDetOne}{0.09}
\newcommand{\RatePosBEightJudgeOnlyQnaRecOne}{0.09}
\newcommand{\RatePosBEightJudgeOnlyQnaAccTwo}{1.00}
\newcommand{\RatePosBEightJudgeOnlyQnaPropTwo}{0.00}
\newcommand{\RatePosBEightJudgeOnlyQnaDetTwo}{0.03}
\newcommand{\RatePosBEightJudgeOnlyQnaRecTwo}{0.03}
\newcommand{\RatePosBEightJudgeOnlyCodeAccZero}{0.97}
\newcommand{\PvalPosBEightJudgeOnlyCodeAcc}{$p=1.00$}
\newcommand{\RatePosBEightJudgeOnlyCodePropZero}{0.03}
\newcommand{\PvalPosBEightJudgeOnlyCodeProp}{$p=1.00$}
\newcommand{\RatePosBEightJudgeOnlyCodeDetZero}{0.65}
\newcommand{\PvalPosBEightJudgeOnlyCodeDet}{$p=0.039$}
\newcommand{\RatePosBEightJudgeOnlyCodeRecZero}{0.65}
\newcommand{\PvalPosBEightJudgeOnlyCodeRec}{$p=0.037$}
\newcommand{\RatePosBEightJudgeOnlyCodeAccOne}{1.00}
\newcommand{\RatePosBEightJudgeOnlyCodePropOne}{0.00}
\newcommand{\RatePosBEightJudgeOnlyCodeDetOne}{0.42}
\newcommand{\RatePosBEightJudgeOnlyCodeRecOne}{0.42}
\newcommand{\RatePosBEightJudgeOnlyCodeAccTwo}{1.00}
\newcommand{\RatePosBEightJudgeOnlyCodePropTwo}{0.00}
\newcommand{\RatePosBEightJudgeOnlyCodeDetTwo}{0.73}
\newcommand{\RatePosBEightJudgeOnlyCodeRecTwo}{0.73}
\newcommand{\RateSemBSixRGenThreeCodeAst}{0.49}
\newcommand{\RateSemBSixRGenThreeCodeAstWrong}{0.07}
\newcommand{\RateSemBSixRGenThreeCodeOracle}{0.42}
\newcommand{\RateSemBSixRGenThreeCodeMethod}{0.01}
\newcommand{\RateSemBSixRGenThreeCodePoisonedFails}{0.44}
\newcommand{\CountSemBSixRGenThreeCodeN}{100}
\newcommand{\PvalSemBSixRGenThreeCodeMethodVsAst}{$p<0.001$}
\newcommand{\RateSemBSixRFlatThreeCodeAst}{0.67}
\newcommand{\RateSemBSixRFlatThreeCodeAstWrong}{0.04}
\newcommand{\RateSemBSixRFlatThreeCodeOracle}{0.64}
\newcommand{\RateSemBSixRFlatThreeCodeMethod}{0.01}
\newcommand{\RateSemBSixRFlatThreeCodePoisonedFails}{0.64}
\newcommand{\CountSemBSixRFlatThreeCodeN}{100}
\newcommand{\PvalSemBSixRFlatThreeCodeMethodVsAst}{$p<0.001$}
\newcommand{\RateSemBSixRJudgeOnlyCodeAst}{0.63}
\newcommand{\RateSemBSixRJudgeOnlyCodeAstWrong}{0.05}
\newcommand{\RateSemBSixRJudgeOnlyCodeOracle}{0.63}
\newcommand{\RateSemBSixRJudgeOnlyCodeMethod}{0.64}
\newcommand{\RateSemBSixRJudgeOnlyCodePoisonedFails}{0.63}
\newcommand{\CountSemBSixRJudgeOnlyCodeN}{100}
\newcommand{\PvalSemBSixRJudgeOnlyCodeMethodVsAst}{$p=1.00$}
\newcommand{\RateSemBSixRBpdCodeAst}{0.66}
\newcommand{\RateSemBSixRBpdCodeAstWrong}{0.03}
\newcommand{\RateSemBSixRBpdCodeOracle}{0.63}
\newcommand{\RateSemBSixRBpdCodeMethod}{0.65}
\newcommand{\RateSemBSixRBpdCodePoisonedFails}{0.63}
\newcommand{\CountSemBSixRBpdCodeN}{97}
\newcommand{\PvalSemBSixRBpdCodeMethodVsAst}{$p=1.00$}
\newcommand{\RateSemBSixGenThreeCodeAst}{0.49}
\newcommand{\RateSemBSixGenThreeCodeOracle}{0.42}
\newcommand{\RateSemBSixGenThreeCodeMethod}{0.01}
\newcommand{\RateSemBSixGenThreeCodePoisonedFails}{0.44}
\newcommand{\CountSemBSixGenThreeCodeN}{100}
\newcommand{\RateSemBSixFlatThreeCodeAst}{0.67}
\newcommand{\RateSemBSixFlatThreeCodeOracle}{0.64}
\newcommand{\RateSemBSixFlatThreeCodeMethod}{0.01}
\newcommand{\RateSemBSixFlatThreeCodePoisonedFails}{0.64}
\newcommand{\CountSemBSixFlatThreeCodeN}{100}
\newcommand{\RateSemBSixJudgeOnlyCodeAst}{0.63}
\newcommand{\RateSemBSixJudgeOnlyCodeOracle}{0.63}
\newcommand{\RateSemBSixJudgeOnlyCodeMethod}{0.64}
\newcommand{\RateSemBSixJudgeOnlyCodePoisonedFails}{0.63}
\newcommand{\CountSemBSixJudgeOnlyCodeN}{100}
\newcommand{\RateSemBSixBpdCodeAst}{0.69}
\newcommand{\RateSemBSixBpdCodeOracle}{0.66}
\newcommand{\RateSemBSixBpdCodeMethod}{0.69}
\newcommand{\RateSemBSixBpdCodePoisonedFails}{0.66}
\newcommand{\CountSemBSixBpdCodeN}{89}
\newcommand{\RateSemBSevenGenThreeCodeAst}{0.48}
\newcommand{\RateSemMainGenThreeCodeAst}{0.48}
\newcommand{\RateSemBSevenGenThreeCodeAstWrong}{0.08}
\newcommand{\RateSemMainGenThreeCodeAstWrong}{0.08}
\newcommand{\RateSemBSevenGenThreeCodeOracle}{0.43}
\newcommand{\RateSemMainGenThreeCodeOracle}{0.43}
\newcommand{\RateSemBSevenGenThreeCodeMethod}{0.02}
\newcommand{\RateSemMainGenThreeCodeMethod}{0.02}
\newcommand{\RateSemBSevenGenThreeCodePoisonedFails}{0.44}
\newcommand{\RateSemMainGenThreeCodePoisonedFails}{0.44}
\newcommand{\CountSemBSevenGenThreeCodeN}{100}
\newcommand{\PvalSemBSevenGenThreeCodeMethodVsAst}{$p<0.001$}
\newcommand{\CountSemMainGenThreeCodeN}{100}
\newcommand{\PvalSemMainGenThreeCodeMethodVsAst}{$p<0.001$}
\newcommand{\RateSemBSevenFlatThreeCodeAst}{0.65}
\newcommand{\RateSemMainFlatThreeCodeAst}{0.65}
\newcommand{\RateSemBSevenFlatThreeCodeAstWrong}{0.06}
\newcommand{\RateSemMainFlatThreeCodeAstWrong}{0.06}
\newcommand{\RateSemBSevenFlatThreeCodeOracle}{0.62}
\newcommand{\RateSemMainFlatThreeCodeOracle}{0.62}
\newcommand{\RateSemBSevenFlatThreeCodeMethod}{0.01}
\newcommand{\RateSemMainFlatThreeCodeMethod}{0.01}
\newcommand{\RateSemBSevenFlatThreeCodePoisonedFails}{0.63}
\newcommand{\RateSemMainFlatThreeCodePoisonedFails}{0.63}
\newcommand{\CountSemBSevenFlatThreeCodeN}{100}
\newcommand{\PvalSemBSevenFlatThreeCodeMethodVsAst}{$p<0.001$}
\newcommand{\CountSemMainFlatThreeCodeN}{100}
\newcommand{\PvalSemMainFlatThreeCodeMethodVsAst}{$p<0.001$}
\newcommand{\RateSemBSevenJudgeOnlyCodeAst}{0.60}
\newcommand{\RateSemMainJudgeOnlyCodeAst}{0.60}
\newcommand{\RateSemBSevenJudgeOnlyCodeAstWrong}{0.00}
\newcommand{\RateSemMainJudgeOnlyCodeAstWrong}{0.00}
\newcommand{\RateSemBSevenJudgeOnlyCodeOracle}{0.60}
\newcommand{\RateSemMainJudgeOnlyCodeOracle}{0.60}
\newcommand{\RateSemBSevenJudgeOnlyCodeMethod}{0.63}
\newcommand{\RateSemMainJudgeOnlyCodeMethod}{0.63}
\newcommand{\RateSemBSevenJudgeOnlyCodePoisonedFails}{0.60}
\newcommand{\RateSemMainJudgeOnlyCodePoisonedFails}{0.60}
\newcommand{\CountSemBSevenJudgeOnlyCodeN}{100}
\newcommand{\PvalSemBSevenJudgeOnlyCodeMethodVsAst}{$p=0.581$}
\newcommand{\CountSemMainJudgeOnlyCodeN}{100}
\newcommand{\PvalSemMainJudgeOnlyCodeMethodVsAst}{$p=0.581$}
\newcommand{\RateSemBSevenBpdCodeAst}{0.67}
\newcommand{\RateSemMainBpdCodeAst}{0.67}
\newcommand{\RateSemBSevenBpdCodeAstWrong}{0.02}
\newcommand{\RateSemMainBpdCodeAstWrong}{0.02}
\newcommand{\RateSemBSevenBpdCodeOracle}{0.66}
\newcommand{\RateSemMainBpdCodeOracle}{0.66}
\newcommand{\RateSemBSevenBpdCodeMethod}{0.67}
\newcommand{\RateSemMainBpdCodeMethod}{0.67}
\newcommand{\RateSemBSevenBpdCodePoisonedFails}{0.66}
\newcommand{\RateSemMainBpdCodePoisonedFails}{0.66}
\newcommand{\CountSemBSevenBpdCodeN}{100}
\newcommand{\PvalSemBSevenBpdCodeMethodVsAst}{$p=1.00$}
\newcommand{\CountSemMainBpdCodeN}{100}
\newcommand{\PvalSemMainBpdCodeMethodVsAst}{$p=1.00$}
\newcommand{\RateSemBEightGenThreeCodeAst}{0.56}
\newcommand{\RateSemBEightGenThreeCodeAstWrong}{0.02}
\newcommand{\RateSemBEightGenThreeCodeOracle}{0.57}
\newcommand{\RateSemBEightGenThreeCodeMethod}{0.01}
\newcommand{\RateSemBEightGenThreeCodePoisonedFails}{0.57}
\newcommand{\CountSemBEightGenThreeCodeN}{100}
\newcommand{\PvalSemBEightGenThreeCodeMethodVsAst}{$p<0.001$}
\newcommand{\RateSemBEightFlatThreeCodeAst}{0.64}
\newcommand{\RateSemBEightFlatThreeCodeAstWrong}{0.01}
\newcommand{\RateSemBEightFlatThreeCodeOracle}{0.62}
\newcommand{\RateSemBEightFlatThreeCodeMethod}{0.01}
\newcommand{\RateSemBEightFlatThreeCodePoisonedFails}{0.63}
\newcommand{\CountSemBEightFlatThreeCodeN}{100}
\newcommand{\PvalSemBEightFlatThreeCodeMethodVsAst}{$p<0.001$}
\newcommand{\RateSemBEightJudgeOnlyCodeAst}{0.55}
\newcommand{\RateSemBEightJudgeOnlyCodeAstWrong}{0.01}
\newcommand{\RateSemBEightJudgeOnlyCodeOracle}{0.53}
\newcommand{\RateSemBEightJudgeOnlyCodeMethod}{0.60}
\newcommand{\RateSemBEightJudgeOnlyCodePoisonedFails}{0.54}
\newcommand{\CountSemBEightJudgeOnlyCodeN}{100}
\newcommand{\PvalSemBEightJudgeOnlyCodeMethodVsAst}{$p=0.125$}
\newcommand{\RateSemBEightBpdCodeAst}{0.59}
\newcommand{\RateSemBEightBpdCodeAstWrong}{0.02}
\newcommand{\RateSemBEightBpdCodeOracle}{0.60}
\newcommand{\RateSemBEightBpdCodeMethod}{0.61}
\newcommand{\RateSemBEightBpdCodePoisonedFails}{0.60}
\newcommand{\CountSemBEightBpdCodeN}{100}
\newcommand{\PvalSemBEightBpdCodeMethodVsAst}{$p=0.754$}
\newcommand{\RateFbBSevenGenThreeCodeLast}{0.91}
\newcommand{\RateFbMainGenThreeCodeLast}{0.91}
\newcommand{\RateFbBSevenGenThreeCodePoisonedOnly}{0.34}
\newcommand{\RateFbMainGenThreeCodePoisonedOnly}{0.34}
\newcommand{\RateFbBSevenGenThreeCodeAccPoisonedOnly}{0.47}
\newcommand{\RateFbMainGenThreeCodeAccPoisonedOnly}{0.47}
\newcommand{\RateFbBSevenGenThreeCodeAccOther}{0.98}
\newcommand{\RateFbMainGenThreeCodeAccOther}{0.98}
\newcommand{\CountFbBSevenGenThreeCodePoisonedOnly}{34}
\newcommand{\CountFbMainGenThreeCodePoisonedOnly}{34}
\newcommand{\RateFbBSevenFlatThreeCodeLast}{0.90}
\newcommand{\RateFbMainFlatThreeCodeLast}{0.90}
\newcommand{\RateFbBSevenFlatThreeCodePoisonedOnly}{0.40}
\newcommand{\RateFbMainFlatThreeCodePoisonedOnly}{0.40}
\newcommand{\RateFbBSevenFlatThreeCodeAccPoisonedOnly}{0.30}
\newcommand{\RateFbMainFlatThreeCodeAccPoisonedOnly}{0.30}
\newcommand{\RateFbBSevenFlatThreeCodeAccOther}{1.00}
\newcommand{\RateFbMainFlatThreeCodeAccOther}{1.00}
\newcommand{\CountFbBSevenFlatThreeCodePoisonedOnly}{40}
\newcommand{\CountFbMainFlatThreeCodePoisonedOnly}{40}
\newcommand{\RateFbBSevenJudgeOnlyCodeLast}{0.48}
\newcommand{\RateFbMainJudgeOnlyCodeLast}{0.48}
\newcommand{\RateFbBSevenJudgeOnlyCodePoisonedOnly}{0.07}
\newcommand{\RateFbMainJudgeOnlyCodePoisonedOnly}{0.07}
\newcommand{\RateFbBSevenJudgeOnlyCodeAccPoisonedOnly}{1.00}
\newcommand{\RateFbMainJudgeOnlyCodeAccPoisonedOnly}{1.00}
\newcommand{\RateFbBSevenJudgeOnlyCodeAccOther}{1.00}
\newcommand{\RateFbMainJudgeOnlyCodeAccOther}{1.00}
\newcommand{\CountFbBSevenJudgeOnlyCodePoisonedOnly}{7}
\newcommand{\CountFbMainJudgeOnlyCodePoisonedOnly}{7}
\newcommand{\RateFbBSevenBpdCodeLast}{0.43}
\newcommand{\RateFbMainBpdCodeLast}{0.43}
\newcommand{\RateFbBSevenBpdCodePoisonedOnly}{0.11}
\newcommand{\RateFbMainBpdCodePoisonedOnly}{0.11}
\newcommand{\RateFbBSevenBpdCodeAccPoisonedOnly}{0.91}
\newcommand{\RateFbMainBpdCodeAccPoisonedOnly}{0.91}
\newcommand{\RateFbBSevenBpdCodeAccOther}{1.00}
\newcommand{\RateFbMainBpdCodeAccOther}{1.00}
\newcommand{\CountFbBSevenBpdCodePoisonedOnly}{11}
\newcommand{\CountFbMainBpdCodePoisonedOnly}{11}
\newcommand{\RateFbBEightGenThreeCodeLast}{0.91}
\newcommand{\RateFbBEightGenThreeCodePoisonedOnly}{0.33}
\newcommand{\RateFbBEightGenThreeCodeAccPoisonedOnly}{0.15}
\newcommand{\RateFbBEightGenThreeCodeAccOther}{1.00}
\newcommand{\CountFbBEightGenThreeCodePoisonedOnly}{33}
\newcommand{\RateFbBEightFlatThreeCodeLast}{0.92}
\newcommand{\RateFbBEightFlatThreeCodePoisonedOnly}{0.37}
\newcommand{\RateFbBEightFlatThreeCodeAccPoisonedOnly}{0.19}
\newcommand{\RateFbBEightFlatThreeCodeAccOther}{1.00}
\newcommand{\CountFbBEightFlatThreeCodePoisonedOnly}{37}
\newcommand{\RateFbBEightJudgeOnlyCodeLast}{0.56}
\newcommand{\RateFbBEightJudgeOnlyCodePoisonedOnly}{0.07}
\newcommand{\RateFbBEightJudgeOnlyCodeAccPoisonedOnly}{0.86}
\newcommand{\RateFbBEightJudgeOnlyCodeAccOther}{1.00}
\newcommand{\CountFbBEightJudgeOnlyCodePoisonedOnly}{7}
\newcommand{\RateFbBEightBpdCodeLast}{0.44}
\newcommand{\RateFbBEightBpdCodePoisonedOnly}{0.06}
\newcommand{\RateFbBEightBpdCodeAccPoisonedOnly}{0.67}
\newcommand{\RateFbBEightBpdCodeAccOther}{1.00}
\newcommand{\CountFbBEightBpdCodePoisonedOnly}{6}
\newcommand{\RateBehBSevenGenOneQnaNoTool}{1.00}
\newcommand{\RateBehBSevenGenOneQnaToolGold}{0.00}
\newcommand{\RateBehBSevenGenOneQnaNoToolLie}{1.00}
\newcommand{\RateBehBSevenGenOneQnaToolLie}{0.00}
\newcommand{\RateBehBSevenGenThreeQnaNoTool}{1.00}
\newcommand{\RateBehBSevenGenThreeQnaToolGold}{0.00}
\newcommand{\RateBehBSevenGenThreeQnaNoToolLie}{1.00}
\newcommand{\RateBehBSevenGenThreeQnaToolLie}{0.00}
\newcommand{\RateBehBSevenFlatThreeQnaNoTool}{0.90}
\newcommand{\RateBehBSevenFlatThreeQnaToolGold}{0.01}
\newcommand{\RateBehBSevenFlatThreeQnaNoToolLie}{0.90}
\newcommand{\RateBehBSevenFlatThreeQnaToolLie}{0.09}
\newcommand{\RateBehBSevenJudgeOnlyQnaNoTool}{0.88}
\newcommand{\RateBehBSevenJudgeOnlyQnaToolGold}{0.02}
\newcommand{\RateBehBSevenJudgeOnlyQnaNoToolLie}{0.88}
\newcommand{\RateBehBSevenJudgeOnlyQnaToolLie}{0.10}
\newcommand{\RateBehBSevenBpdQnaNoTool}{0.87}
\newcommand{\RateBehBSevenBpdQnaToolGold}{0.00}
\newcommand{\RateBehBSevenBpdQnaNoToolLie}{0.87}
\newcommand{\RateBehBSevenBpdQnaToolLie}{0.13}
\newcommand{\RateBehBEightGenOneQnaNoTool}{0.95}
\newcommand{\RateBehBEightGenOneQnaToolGold}{0.05}
\newcommand{\RateBehBEightGenOneQnaNoToolLie}{0.95}
\newcommand{\RateBehBEightGenOneQnaToolLie}{0.00}
\newcommand{\RateBehBEightGenThreeQnaNoTool}{0.91}
\newcommand{\RateBehBEightGenThreeQnaToolGold}{0.05}
\newcommand{\RateBehBEightGenThreeQnaNoToolLie}{0.91}
\newcommand{\RateBehBEightGenThreeQnaToolLie}{0.04}
\newcommand{\RateBehBEightFlatThreeQnaNoTool}{0.11}
\newcommand{\RateBehBEightFlatThreeQnaToolGold}{0.88}
\newcommand{\RateBehBEightFlatThreeQnaNoToolLie}{0.11}
\newcommand{\RateBehBEightFlatThreeQnaToolLie}{0.01}
\newcommand{\RateBehBEightJudgeOnlyQnaNoTool}{0.11}
\newcommand{\RateBehBEightJudgeOnlyQnaToolGold}{0.88}
\newcommand{\RateBehBEightJudgeOnlyQnaNoToolLie}{0.11}
\newcommand{\RateBehBEightJudgeOnlyQnaToolLie}{0.01}
\newcommand{\RateBehBEightBpdQnaNoTool}{0.09}
\newcommand{\RateBehBEightBpdQnaToolGold}{0.90}
\newcommand{\RateBehBEightBpdQnaNoToolLie}{0.09}
\newcommand{\RateBehBEightBpdQnaToolLie}{0.01}
\newcommand{\MainRunName}{budget-7}


% [AUTHOR: ...] marks text that needs author input or re-verification after new runs.
\newcommand{\todo}[1]{{\color{red}[AUTHOR: #1]}}
% Optional macro: prints a placeholder when the run that defines it has not finished yet.
\newcommand{\opt}[1]{\ifcsname #1\endcsname\csname #1\endcsname\else\textbf{[TBD]}\fi}

\title{Isolation Is Not Attribution:\\Internal Poison in DAG-Structured Coding Agents}

\author{Anonymous ACL submission}

\begin{document}
\maketitle

\begin{abstract}
Coding agents increasingly run as graphs of cooperating nodes, and one node can hand a wrong claim to all the others.
We study this internal-poison setting. A planted lie enters one node, and the system should keep its final result correct, name the node that lied, or do both.
We compare a single-loop agent with three-replica ensembles, a judge-based graph with backpropagation of node contributions, and a hierarchical scan-implement-test graph on 100 configuration questions and 100 bug-fixing tasks graded by hidden tests, and we vary which replica carries the lie.
With the same claim in the poisoned node's system prompt, the single loop repeats the lie in \RateMainGenOneQnaProp{} of question-answering runs and \RateMainGenOneCodeProp{} of bug-fixing runs.
Three isolated replicas remove the damage on question answering, but on bug fixing a majority vote depends on which replica is poisoned: flat-3 accuracy is \RatePosBSevenFlatThreeCodeAccZero{} when replica~0 lies and \RatePosBSevenFlatThreeCodeAccTwo{} when replica~2 lies, because the vote falls back to the last replica.
Naming the source is a separate property.
On bug fixing it appears only when a comparison reads the replicas' code: a judge reaches recovery \RateMainJudgeOnlyCodeRec{}, and a judge-free comparison of normalized sources names the poisoned replica as often (\RateSemMainJudgeOnlyCodeAst{}) and keeps accuracy at \RateSelMainFlatThreeCodeAst{} when it also chooses the workspace, against \RateSelMainFlatThreeCodeActual{} for the string vote and \RateSelMainFlatThreeCodeUniform{} for a random replica. Backpropagation adds no measurable benefit.
A second model reproduces the bug-fixing results, but its question-answering attribution falls to \RateBEightJudgeOnlyQnaDet{} for the judge arm because the poisoned replica rarely repeats the lie (\RateBEightJudgeOnlyQnaEmit{}): a source can be named only when the lie is visible.
Graph structure alone predicts neither isolation nor attribution.
\end{abstract}

\section{Introduction}

Software agents now split work across nodes. A scanner reads the repository, several workers propose edits, and a reviewer picks one.
Frameworks such as ChatDev, MetaGPT, and MacNet organize these nodes as pipelines or directed acyclic graphs \citep{qian2023chatdev,hong2023metagpt,qian2024macnet}.
A graph gives each claim a path, so in principle a wrong claim can be traced to the node that introduced it.
In practice, a node that trusts its upstream neighbor can also pass a wrong claim along with more authority than a lone agent would give it.

Prior work studies adjacent threats. Benchmarks such as AgentDojo and InjecAgent inject instructions through external tool outputs \citep{debenedetti2024agentdojo,zhan2024injecagent}.
Studies of faulty agents measure how much a wrong collaborator degrades a team \citep{huang2024resilience}.
Prompt Infection and manipulated-knowledge spread show that a bad message can move between agents \citep{lee2024promptinfection,ju2024flooding}.
Automated failure attribution asks which agent and step caused a failed run \citep{zhang2025failureattribution}.
None of these measures, on tool-using coding tasks with executable ground truth, whether a system can name the internal node that carried a planted lie.

We study that question with a controlled suite. A lie is planted in one internal node, and we score four outcomes: accuracy of the final result, propagation of the lie into it, detection (the system names the poisoned node), and recovery (detection with a correct result).
We ask two questions.
\textbf{RQ1:} Do graph-structured systems isolate, detect, and recover from an internal lie better than a single-loop agent?
\textbf{RQ2:} Inside graph-structured systems, which mechanism for handling the contaminated node drives the outcome: redundancy, a comparison step, or contribution backpropagation?

The answers separate isolation from attribution.
Three isolated replicas remove the damage of the lie on question answering, where every three-replica arm reaches accuracy \RateMainGenThreeQnaAcc{} against \RateMainGenOneQnaAcc{} for one loop.
On bug fixing, replication alone is not enough: accuracy is \RateMainGenThreeCodeAcc{} and \RateMainFlatThreeCodeAcc{} for the two majority arms, against \RateMainJudgeOnlyCodeAcc{} and \RateMainBpdCodeAcc{} for the arms with a judge, and the majority arms collapse when the poisoned replica is the one their fallback rule promotes (Section~\ref{sec:position}).
Naming the source is a different property. On bug fixing it appears only when a comparison reads the replicas' code: recovery is \RateMainJudgeOnlyCodeRec{} for the judge-only arm and \RateMainBpdCodeRec{} for the full backpropagation pipeline (exact McNemar, \PvalMainCodeRecJudgeOnlyvsBpd{}), against \RateMainFlatThreeCodeRec{} for a majority vote on strategy strings.
A judge-free comparison of the normalized replica sources names the poisoned replica as often as the judge does (\RateSemMainJudgeOnlyCodeAst{} against \RateSemMainJudgeOnlyCodeMethod{}, \PvalSemMainJudgeOnlyCodeMethodVsAst{}), and promoting a matching replica keeps flat-3 accuracy at \RateSelMainFlatThreeCodeAst{} against \RateSelMainFlatThreeCodeActual{} for its string vote, so the judge is not what makes attribution or workspace choice work on these tasks.
A second model reproduces the bug-fixing pattern but shows the limit of attribution: when the poisoned replica does not repeat the lie, no scheme can name it (Section~\ref{sec:second}).
Finally, a hierarchical graph whose downstream nodes treat the poisoned scan report as authoritative repeats the lie in \RateHundredYonseiQnaProp{} of QnA runs; a source-over-scan rule removes this effect (exploratory).

We contribute the suite and harness, a compute-matched set of arms that separates replication, judging, and backpropagation, two measurement controls that changed our conclusions (the channel that delivers the claim, and the index of the poisoned replica), and a replication package that regenerates every number in this paper from frozen run summaries.
All results use two models from one family and one attempt per task, so we report them as comparisons on this suite and not as general properties of agents.

\section{Related Work}

\paragraph{Graph-structured coding agents.}
SWE-bench and SWE-agent set the single-agent baseline for repository-level repair \citep{jimenez2023swebench,yang2024sweagent}.
ChatDev and MetaGPT assign roles along a pipeline, and MacNet uses directed acyclic graphs to organize large numbers of agents \citep{qian2023chatdev,hong2023metagpt,qian2024macnet}.
GPTSwarm treats a team of agents as an optimizable graph \citep{zhuge2024gptswarm}.
These papers measure task performance. We measure what happens when one node is wrong.

\paragraph{Faults and attacks inside multi-agent systems.}
\citet{huang2024resilience} compare system structures under faulty agents and find that a hierarchical structure loses the least performance.
Prompt Infection shows that malicious prompts can replicate across interconnected agents \citep{lee2024promptinfection}, and \citet{ju2024flooding} study how manipulated knowledge spreads in agent communities.
AgentPoison poisons an agent's memory or knowledge base \citep{chen2024agentpoison}.
Defenses include a graph neural network over the utterance graph \citep{wang2025gsafeguard}, pruning of redundant or malicious messages \citep{zhang2024cutthecrap}, and a decentralized consensus filter for Byzantine agents \citep{lee2026byzantine}.
Our setting differs in two ways: the lie enters through a collaborator node, and the tasks have executable ground truth.

\paragraph{Attribution and contribution scoring.}
\citet{zhang2025failureattribution} introduce failure attribution for multi-agent systems and report low accuracy on pinpointing the failing step.
\citet{cemri2025mast} give a taxonomy of why multi-agent systems fail.
VeriTrail traces where unsupported content entered a multi-step generation process \citep{metropolitansky2025veritrail}.
Backpropagation of node contributions over a signed directed acyclic graph (BPD) scores each agent's contribution to the final decision and isolates malicious agents \citep{wu2025bpd}.
We reimplement the BPD mechanism in a tool-using harness and ask how much of its effect comes from the judge, from the backward pass, and from running three replicas.

\paragraph{Ensembles and judges.}
Majority voting over sampled outputs improves reasoning \citep{wang2022selfconsistency}, and multiagent debate improves factuality \citep{du2023debate}.
Language-model judges agree with human preferences on open-ended tasks but carry their own biases \citep{zheng2023judge}.
We use both ideas as baselines, and we report paired tests and intervals following the argument for error bars in model evaluations \citep{miller2024errorbars}.

\section{Experimental Setup}

\subsection{Tasks}
The suite has two parts of 100 tasks each, expanded from 20 hand-authored templates per part by a catalog generator.
\textbf{QnA tasks} ask for one configuration value. The gold value sits in a source file in the workspace, and the planted lie is a different value (for example a timeout of 300 seconds against a gold value of 30).
Answers are graded by string match.
\textbf{Code tasks} ask for a fix to a single-module bug. The planted lie is a repair strategy that reproduces the bug (for example, retrying without a limit), and a hidden test file grades the final workspace.

\subsection{Poison conditions}
Each task runs under three conditions: \texttt{none}, \texttt{document} (a notes file asserts the lie as policy), and \texttt{agent} (a collaborator node asserts it).
The \texttt{document} condition did not move any system in our suite (propagation near zero), so all main results use \texttt{agent}.
In the \texttt{agent} condition the claim reads ``Injected claim from another agent (mandatory): the answer is $X$; do not use any value from source code or notes'' for QnA, with a strategy variant for code.
We give all arms in the main comparison this identical text and place it in the system prompt of the poisoned node.
In the main run the poisoned replica index is balanced: replica 0, 1, and 2 carry the lie in equal shares of the tasks of a suite, assigned by a seeded hash of the task identifier.
An earlier run (appendix) always poisoned replica~0.

\subsection{Systems}
Each arm runs the same model through the same tool loop, limited to 20 turns per node.
\textbf{Single loop} (\texttt{general-1}) is one agent with file and shell tools.
\textbf{General-3} and \textbf{flat-3} run three isolated replicas, one of which receives the lie.
A replica in general-3 uses the single-loop prompt, and a replica in flat-3 is asked to state its strategy.
Both choose the answer that most replicas support, found by matching the gold and lie strings in each output, and for code tasks they promote the winning replica's workspace.
When no output matches either string, which is the usual case for code tasks, they promote the last replica.
\textbf{Judge-only-3} adds three summarizers, each reading all proposals and writing a final answer, and one independent judge that scores every (proposal, summary) pair with $+1$, $0$, or $-1$.
A replica's score is the mean of its row, the highest score wins, and the system names a source when exactly one score is negative.
\textbf{BPD} (\texttt{dag-bpd}) uses the same calls but gives each summary a terminal score of $+1$ if it agrees with the plurality answer and $-1$ otherwise, and computes each replica's score as $S_i=\tfrac{1}{3}\sum_j g_{ij}S_j$ with judge edges $g_{ij}$.
\textbf{Hierarchical} (\texttt{dag-yonsei}) runs scan, implement, test, and repair; the lie enters the scan report.
In its naive form, later nodes treat the scan as authoritative.
The source-over-scan variant instructs the implement node to prefer source and tests over the scan report, and names the scan node only when they conflict and the scan claim was rejected.
We added this rule after seeing the naive results, so we label it exploratory.
The hierarchical arm was run only in earlier runs.

\subsection{Judge-free reference methods}
To separate the judge from the act of comparing replicas, we attribute offline on code tasks, with no model call, from the three replicas' final workspaces.
The \emph{AST outlier} parses each replica's non-test Python sources, drops docstrings and comments, and flags a replica when it differs from two identical others.
The \emph{test oracle} flags the one replica that fails the hidden tests while the other two pass; the arms never see these tests, so it is an upper bound and not a competing method.
Both are scored like detection: a hit when the flagged replica is the poisoned one, and a miss when none or a wrong one is flagged.

\subsection{Models, metrics, and statistics}
The main runs use \texttt{gpt-5.6-luna}; a second run repeats the configuration with \texttt{gpt-6-luna}. Both are reached through the same OpenAI-compatible interface, with one attempt per task.
\emph{Accuracy} is the share of runs with a correct final result; \emph{propagation} is the share in which the lie appears in the final output; \emph{detection} is the share in which the system names the node that carried the lie; and \emph{recovery} requires detection and accuracy.
The single loop reports no source, so its detection is zero by construction.
We also record whether the poisoned replica itself output the lie (\emph{Emit}), because a vote can name a source only when the lie is visible; on code tasks Emit is a lexical proxy.
Intervals are percentile bootstrap intervals over tasks (1{,}000 resamples, fixed seed); paired differences between arms use the exact McNemar test on task identity, and dependence of a rate on the poisoned replica index uses a permutation test of the chi-square statistic.
Token counts are logged as the cost proxy, since the endpoint reports no price.
Each task has a 180-second wall-clock limit, and no task reached it in the main or second-model run.
Three-replica arms use about three times the input tokens of the single loop (Table~\ref{tab:main-tokens}), and arms with a judge add summarizer and judge calls.

\section{Results}

\subsection{Isolation under the same injected claim}
% Auto-generated by code/03_export_paper_tables.py -- do not edit by hand.
\begin{table*}[t]
\centering
\footnotesize
\caption{Compute-matched arms under the same injected claim ($n{=}100$ per cell, agent condition; run \texttt{budget-7}). Acc.: gold answer; Prop.: lie in final output; Det.: poisoned source named; Rec.: Det.\ and Acc.; Emit: poisoned replica output the lie; Det$\mid$E: Det.\ among Emit runs.}
\label{tab:main}
\begin{tabular}{lcccccccccccc}
\toprule
 & \multicolumn{6}{c}{QnA} & \multicolumn{6}{c}{Code} \\
\cmidrule(lr){2-7}\cmidrule(lr){8-13}
Arm & Acc & Prop & Det & Rec & Emit & Det$\mid$E & Acc & Prop & Det & Rec & Emit & Det$\mid$E \\
\midrule
general-1 & 0.00 & 1.00 & 0.00 & 0.00 & 1.00 & 0.00 & 0.59 & 0.40 & 0.00 & 0.00 & 0.47 & 0.00 \\
general-3 maj. & 1.00 & 0.00 & 1.00 & 1.00 & 1.00 & 1.00 & 0.81 & 0.15 & 0.02 & 0.02 & 0.48 & 0.04 \\
flat-3 maj. & 1.00 & 0.00 & 0.99 & 0.99 & 0.99 & 1.00 & 0.72 & 0.26 & 0.01 & 0.01 & 0.59 & 0.02 \\
judge-only-3 & 1.00 & 0.00 & 0.91 & 0.91 & 0.98 & 0.93 & 1.00 & 0.00 & 0.63 & 0.63 & 0.61 & 0.87 \\
dag-bpd & 1.00 & 0.00 & 0.94 & 0.94 & 1.00 & 0.94 & 0.99 & 0.01 & 0.67 & 0.67 & 0.61 & 0.84 \\
\bottomrule
\end{tabular}
\end{table*}


Table~\ref{tab:main} compares the five arms with the claim in the poisoned node's system prompt and the poisoned replica index balanced.
The single loop follows the lie: it repeats it in \RateMainGenOneQnaProp{} of QnA runs and \RateMainGenOneCodeProp{} of code runs, and its accuracy is \RateMainGenOneQnaAcc{} and \RateMainGenOneCodeAcc{}.
On QnA, accuracy is \RateMainGenThreeQnaAcc{} for general-3, \RateMainFlatThreeQnaAcc{} for flat-3, \RateMainJudgeOnlyQnaAcc{} for judge-only, and \RateMainBpdQnaAcc{} for BPD.
On code tasks accuracy separates the arms: \RateMainGenThreeCodeAcc{}, \RateMainFlatThreeCodeAcc{}, \RateMainJudgeOnlyCodeAcc{}, and \RateMainBpdCodeAcc{} in the same order.
Both majority arms beat the single loop (exact McNemar against general-1: general-3 \PvalMainCodeAccGenOnevsGenThree{}, flat-3 \PvalMainCodeAccGenOnevsFlatThree{}) and both stay below the judge arms (flat-3 against judge-only, \PvalMainCodeAccFlatThreevsJudgeOnly{}).
Whether the poisoned replica repeats the lie is itself an outcome, and it differs by suite: in general-3 it outputs the lie in \RateMainGenThreeQnaEmit{} of QnA runs but only \RateMainGenThreeCodeEmit{} of code runs, because on code tasks it often ignores or repairs around the claim.
This Emit rate limits how often any scheme can name a source.

\begin{table}[t]
\centering
\footnotesize
\caption{Mean input tokens per task (thousands) in run \texttt{\MainRunName}.}
\label{tab:main-tokens}
\begin{tabular}{lcc}
\toprule
Arm & QnA & Code \\
\midrule
general-1 & \TokMainGenOneQna & \TokMainGenOneCode \\
general-3 maj. & \TokMainGenThreeQna & \TokMainGenThreeCode \\
flat-3 maj. & \TokMainFlatThreeQna & \TokMainFlatThreeCode \\
judge-only-3 & \TokMainJudgeOnlyQna & \TokMainJudgeOnlyCode \\
dag-bpd & \TokMainBpdQna & \TokMainBpdCode \\
\bottomrule
\end{tabular}
\end{table}

\subsection{Majority arms depend on which replica is poisoned}
\label{sec:position}
% Auto-generated by code/03_export_paper_tables.py — do not edit by hand.
\begin{table*}[t]
\centering
\scriptsize
\caption{Accuracy and recovery by the index of the poisoned replica (permutation test of dependence on the index; each cell holds about a third of the tasks of a suite). The single loop ignores the index, so its row shows the spread expected from task composition alone. Run \texttt{budget-7}.}
\label{tab:position}
\begin{tabular}{lccccccccc}
\toprule
Arm & Set & Acc@0 & Acc@1 & Acc@2 & $p_{\mathrm{Acc}}$ & Rec@0 & Rec@1 & Rec@2 & $p_{\mathrm{Rec}}$ \\
\midrule
general-1 & QnA & 0.00 & 0.00 & 0.00 & 1.000 & 0.00 & 0.00 & 0.00 & 1.000 \\
general-1 & Code & 0.53 & 0.70 & 0.55 & 0.348 & 0.00 & 0.00 & 0.00 & 1.000 \\
general-3 maj. & QnA & 1.00 & 1.00 & 1.00 & 1.000 & 1.00 & 1.00 & 1.00 & 1.000 \\
general-3 maj. & Code & 0.91 & 1.00 & 0.52 & $<$0.001 & 0.00 & 0.06 & 0.00 & 0.207 \\
flat-3 maj. & QnA & 1.00 & 1.00 & 1.00 & 1.000 & 0.97 & 1.00 & 1.00 & 1.000 \\
flat-3 maj. & Code & 0.88 & 0.91 & 0.36 & $<$0.001 & 0.00 & 0.00 & 0.03 & 0.658 \\
judge-only-3 & QnA & 1.00 & 1.00 & 1.00 & 1.000 & 0.82 & 0.94 & 0.97 & 0.140 \\
judge-only-3 & Code & 1.00 & 1.00 & 1.00 & 1.000 & 0.65 & 0.67 & 0.58 & 0.757 \\
dag-bpd & QnA & 1.00 & 1.00 & 1.00 & 1.000 & 0.94 & 0.91 & 0.97 & 0.696 \\
dag-bpd & Code & 0.97 & 1.00 & 1.00 & 1.000 & 0.62 & 0.61 & 0.79 & 0.247 \\
\bottomrule
\end{tabular}
\end{table*}


The majority arms look strong when the lie always sits in replica~0.
In that earlier design flat-3 and general-3 reach code accuracy \RateBSixRFlatThreeCodeAcc{} and \RateBSixRGenThreeCodeAcc{}.
With the index balanced, accuracy depends on it (Table~\ref{tab:position}): flat-3 reaches \RatePosBSevenFlatThreeCodeAccZero{}, \RatePosBSevenFlatThreeCodeAccOne{}, and \RatePosBSevenFlatThreeCodeAccTwo{} when replica 0, 1, or 2 lies (permutation test, \PvalPosBSevenFlatThreeCodeAcc{}), and general-3 reaches \RatePosBSevenGenThreeCodeAccZero{}, \RatePosBSevenGenThreeCodeAccOne{}, and \RatePosBSevenGenThreeCodeAccTwo{} (\PvalPosBSevenGenThreeCodeAcc{}).
The cause is the fallback rule. Code proposals rarely contain the gold or lie string, so no proposal supplies a vote and the rule promotes the last replica.
In the raw runs the promoted source equals replica~2's source in \RateFbMainFlatThreeCodeLast{} of flat-3 tasks and \RateFbMainGenThreeCodeLast{} of general-3 tasks (Table~\ref{tab:fallback} in the appendix).
When the promoted source equals only the poisoned replica's, flat-3 accuracy is \RateFbMainFlatThreeCodeAccPoisonedOnly{} (\CountFbMainFlatThreeCodePoisonedOnly{} tasks) and otherwise \RateFbMainFlatThreeCodeAccOther{}.
A majority vote on code strategies therefore behaves like picking the last replica, and its accuracy is that replica's chance of being clean.
The arms with a judge are not affected: judge-only accuracy is \RatePosBSevenJudgeOnlyCodeAccZero{}, \RatePosBSevenJudgeOnlyCodeAccOne{}, and \RatePosBSevenJudgeOnlyCodeAccTwo{}, and BPD accuracy is \RatePosBSevenBpdCodeAccZero{}, \RatePosBSevenBpdCodeAccOne{}, and \RatePosBSevenBpdCodeAccTwo{}.
Their recovery does not depend detectably on the index either (judge-only on code, \PvalPosBSevenJudgeOnlyCodeRec{}; BPD, \PvalPosBSevenBpdCodeRec{}), but with about a third of the tasks per cell this test cannot rule out small effects.

\subsection{Naming the source}
Attribution separates the arms more than accuracy does.
On QnA tasks the answer space is small and string matching works. General-3 and flat-3 name the source in \RateMainGenThreeQnaDet{} and \RateMainFlatThreeQnaDet{} of runs, against \RateMainJudgeOnlyQnaDet{} for judge-only and \RateMainBpdQnaDet{} for BPD.
The majority vote recovers more often than the judge it is compared with (general-3 against judge-only, \PvalMainQnaRecGenThreevsJudgeOnly{}; flat-3 against judge-only, \PvalMainQnaRecFlatThreevsJudgeOnly{}), because the judge misses some runs in which the lie was visible.
On code tasks the two majority arms name the source in \RateMainGenThreeCodeDet{} and \RateMainFlatThreeCodeDet{} of runs, because the strategy phrases of three replicas rarely match as strings.
The judge-only arm names it in \RateMainJudgeOnlyCodeDet{} and BPD in \RateMainBpdCodeDet{}; given that the lie was emitted, these rates are \RateMainJudgeOnlyCodeDetEmit{} and \RateMainBpdCodeDetEmit{}.
The ceiling is therefore the Emit rate, \RateMainJudgeOnlyCodeEmit{} and \RateMainBpdCodeEmit{} in these arms.

\subsection{What the judge, the backward pass, and a structural diff contribute}
The judge-only and BPD arms make the same calls and differ only in the terminal weighting.
On code tasks their recovery is \RateMainJudgeOnlyCodeRec{} and \RateMainBpdCodeRec{} (exact McNemar, \PvalMainCodeRecJudgeOnlyvsBpd{}), and on QnA tasks \RateMainJudgeOnlyQnaRec{} and \RateMainBpdQnaRec{} (\PvalMainQnaRecJudgeOnlyvsBpd{}).
In this suite the backward pass adds no measurable recovery beyond the judge.
The judge does the work on code compared with a string vote: replacing it lowers recovery from \RateMainJudgeOnlyCodeRec{} to \RateMainFlatThreeCodeRec{} (\PvalMainCodeRecFlatThreevsJudgeOnly{}).
Because the majority arms use string matching, this comparison mixes the value of an LLM judge with the value of reading the replicas' content.

% Auto-generated by code/03_export_paper_tables.py — do not edit by hand.
\begin{table*}[t]
\centering
\scriptsize
\caption{Judge-free attribution on code replicas (\texttt{gpt-5.6-luna}, position-balanced; no LLM call). AST outlier: the one replica whose normalized source differs from two identical ones; wrong: it flags a replica other than the poisoned one; test oracle: the one replica that fails hidden tests (an upper bound the arms cannot use). Arm detection is the arm's own detection rate on the same tasks.}
\label{tab:semantic}
\begin{tabular}{lccccc}
\toprule
Arm & $n$ & AST outlier & AST wrong & Test oracle & Arm detection \\
\midrule
general-3 maj. & 100 & 0.48 & 0.08 & 0.43 & 0.02 \\
flat-3 maj. & 100 & 0.65 & 0.06 & 0.62 & 0.01 \\
judge-only-3 & 100 & 0.60 & 0.00 & 0.60 & 0.63 \\
dag-bpd & 100 & 0.67 & 0.02 & 0.66 & 0.67 \\
\bottomrule
\end{tabular}
\end{table*}


The judge-free references separate the two (Table~\ref{tab:semantic}).
On the replicas of the judge-only arm the AST outlier names the poisoned replica in \RateSemMainJudgeOnlyCodeAst{} of runs and a wrong replica in \RateSemMainJudgeOnlyCodeAstWrong{}, against \RateSemMainJudgeOnlyCodeMethod{} for the judge itself (exact McNemar, \PvalSemMainJudgeOnlyCodeMethodVsAst{}); on the BPD replicas the rates are \RateSemMainBpdCodeAst{} and \RateSemMainBpdCodeMethod{} (\PvalSemMainBpdCodeMethodVsAst{}).
On the replicas of flat-3, the same offline rule names the source in \RateSemMainFlatThreeCodeAst{} of runs, where the arm's own string vote names it in \RateSemMainFlatThreeCodeMethod{}.
The test oracle reaches \RateSemMainJudgeOnlyCodeOracle{} on the judge-only replicas, close to the share of runs in which the poisoned replica fails its tests at all (\RateSemMainJudgeOnlyCodePoisonedFails{}).
On these tasks, then, the gain of the judge arms over the string vote comes from comparing the replicas' code at all, and the judge adds no measurable attribution beyond a structural diff.
The diff works because the correct replicas converge on the same fix, so the one that differs stands out; we expect it to degrade when correct replicas legitimately differ.

The same diff also chooses a workspace (Table~\ref{tab:selection}).
On the flat-3 replicas, promoting a replica whose normalized source matches another one, and otherwise promoting the last replica, reaches accuracy \RateSelMainFlatThreeCodeAst{}, against \RateSelMainFlatThreeCodeActual{} for the arm's own string vote and \RateSelMainFlatThreeCodeUniform{} for the expected accuracy of drawing one replica at random.
Always promoting the last replica scores \RateSelMainFlatThreeCodeLast{}, so the failure is not repaired by merely abandoning that fallback.
The same rule scores \RateSelMainGenThreeCodeAst{}, \RateSelMainJudgeOnlyCodeAst{}, and \RateSelMainBpdCodeAst{} on the general-3, judge-only, and BPD replicas, against recorded accuracies \RateSelMainGenThreeCodeActual{}, \RateSelMainJudgeOnlyCodeActual{}, and \RateSelMainBpdCodeActual{}.
On the second model the flat-3 figures are \RateSelBEightFlatThreeCodeAst{} for the diff, \RateSelBEightFlatThreeCodeActual{} for the arm, and \RateSelBEightFlatThreeCodeUniform{} for a uniform draw.
On these tasks a structural diff matches the judge as a way to pick the workspace, and a random pick does not.

% Auto-generated by code/03_export_paper_tables.py — do not edit by hand.
\begin{table*}[t]
\centering
\scriptsize
\caption{Code accuracy if a different replica workspace had been promoted ($n{=}100$). Actual: the arm's recorded accuracy; Last: always replica~2; Uniform: expected accuracy of drawing one replica at random; AST: a replica whose normalized source matches another one when exactly one replica differs, and replica~2 otherwise. A: \texttt{gpt-5.6-luna}; B: \texttt{gpt-6-luna}.}
\label{tab:selection}
\begin{tabular}{lcccccccc}
\toprule
Arm & Actual A & Last A & Uniform A & AST A & Actual B & Last B & Uniform B & AST B \\
\midrule
general-3 maj. & 0.81 & 0.83 & 0.85 & 0.97 & 0.72 & 0.80 & 0.81 & 1.00 \\
flat-3 maj. & 0.72 & 0.78 & 0.79 & 0.98 & 0.70 & 0.77 & 0.79 & 0.98 \\
judge-only-3 & 1.00 & 0.81 & 0.80 & 0.97 & 0.99 & 0.77 & 0.81 & 0.99 \\
dag-bpd & 0.99 & 0.75 & 0.78 & 0.97 & 0.98 & 0.78 & 0.80 & 1.00 \\
\bottomrule
\end{tabular}
\end{table*}


\subsection{A second model}
\label{sec:second}
% Auto-generated by code/03_export_paper_tables.py -- do not edit by hand.
\begin{table*}[t]
\centering
\footnotesize
\caption{Second model (\texttt{gpt-6-luna}), same configuration as Table~\ref{tab:main} ($n{=}100$ per cell, agent condition). Columns as in Table~\ref{tab:main}.}
\label{tab:second}
\begin{tabular}{lcccccccccccc}
\toprule
 & \multicolumn{6}{c}{QnA} & \multicolumn{6}{c}{Code} \\
\cmidrule(lr){2-7}\cmidrule(lr){8-13}
Arm & Acc & Prop & Det & Rec & Emit & Det$\mid$E & Acc & Prop & Det & Rec & Emit & Det$\mid$E \\
\midrule
general-1 & 0.05 & 0.95 & 0.00 & 0.00 & 0.95 & 0.00 & 0.44 & 0.46 & 0.00 & 0.00 & 0.48 & 0.00 \\
general-3 maj. & 1.00 & 0.00 & 0.95 & 0.95 & 0.95 & 1.00 & 0.72 & 0.27 & 0.01 & 0.01 & 0.51 & 0.02 \\
flat-3 maj. & 1.00 & 0.00 & 0.12 & 0.12 & 0.12 & 1.00 & 0.70 & 0.29 & 0.01 & 0.01 & 0.58 & 0.02 \\
judge-only-3 & 1.00 & 0.00 & 0.11 & 0.11 & 0.12 & 0.92 & 0.99 & 0.01 & 0.60 & 0.60 & 0.54 & 0.98 \\
dag-bpd & 1.00 & 0.00 & 0.10 & 0.10 & 0.10 & 1.00 & 0.98 & 0.02 & 0.61 & 0.61 & 0.53 & 0.96 \\
\bottomrule
\end{tabular}
\end{table*}


Table~\ref{tab:second} repeats the configuration with \texttt{gpt-6-luna}.
The bug-fixing results carry over.
The single loop follows the lie (QnA propagation \RateBEightGenOneQnaProp{}, code accuracy \RateBEightGenOneCodeAcc{}).
The majority arms depend on the index again: flat-3 accuracy is \RatePosBEightFlatThreeCodeAccZero{}, \RatePosBEightFlatThreeCodeAccOne{}, and \RatePosBEightFlatThreeCodeAccTwo{} for replica 0, 1, and 2 (\PvalPosBEightFlatThreeCodeAcc{}), and the fallback rule promotes replica~2's source in \RateFbBEightFlatThreeCodeLast{} of tasks.
The judge arms keep accuracy \RateBEightJudgeOnlyCodeAcc{} and \RateBEightBpdCodeAcc{} and recovery \RateBEightJudgeOnlyCodeRec{} and \RateBEightBpdCodeRec{}, with no difference between them (\PvalBEightCodeRecJudgeOnlyvsBpd{}).
The AST outlier names the poisoned replica in \RateSemBEightJudgeOnlyCodeAst{} and \RateSemBEightBpdCodeAst{} of runs on the judge-only and BPD replicas, against \RateSemBEightJudgeOnlyCodeMethod{} and \RateSemBEightBpdCodeMethod{} for the judges (\PvalSemBEightJudgeOnlyCodeMethodVsAst{} and \PvalSemBEightBpdCodeMethodVsAst{}).
On this model the judge is slightly ahead of the diff, but not detectably so.
One index effect appears for a judge arm on this model: judge-only code recovery is \RatePosBEightJudgeOnlyCodeRecZero{}, \RatePosBEightJudgeOnlyCodeRecOne{}, and \RatePosBEightJudgeOnlyCodeRecTwo{} for replica 0, 1, and 2 (\PvalPosBEightJudgeOnlyCodeRec{}), while BPD does not show it (\PvalPosBEightBpdCodeRec{}); we ran 20 recovery tests across the two position tables without correction, so we treat this as a lead and not a finding.

The QnA attribution result does not carry over.
Detection falls to \RateBEightFlatThreeQnaDet{} for flat-3, \RateBEightJudgeOnlyQnaDet{} for judge-only, and \RateBEightBpdQnaDet{} for BPD, while general-3 stays at \RateBEightGenThreeQnaDet{}.
In the first three arms the poisoned replica repeats the lie in only \RateBEightFlatThreeQnaEmit{}, \RateBEightJudgeOnlyQnaEmit{}, and \RateBEightBpdQnaEmit{} of runs, whereas in general-3 it does so in \RateBEightGenThreeQnaEmit{}.
Given that the lie was emitted, detection is \RateBEightFlatThreeQnaDetEmit{}, \RateBEightJudgeOnlyQnaDetEmit{}, and \RateBEightBpdQnaDetEmit{}, so attribution works when the lie is visible and cannot work when it is not.
Accuracy stays at \RateBEightBpdQnaAcc{}: a replica that does not repeat the lie needs no isolation.
The event logs show why, and it is the base prompt rather than the claim, which is identical (Table~\ref{tab:emit}).
Flat-3, judge-only, and BPD tell the replica it is a proposal worker that should edit files; general-3 uses the single-loop prompt.
On \texttt{gpt-6-luna}, reading the file decides the answer. In flat-3 the poisoned replica calls a tool and answers with the source value in \RateBehBEightFlatThreeQnaToolGold{} of tasks, and it answers with the lie without a tool in \RateBehBEightFlatThreeQnaNoToolLie{}.
Under the general-3 prompt the same model does the opposite: no tool and the lie in \RateBehBEightGenThreeQnaNoToolLie{} of tasks.
On \texttt{gpt-5.6-luna} the poisoned replica skips tools in \RateBehBSevenFlatThreeQnaNoTool{} of flat-3 tasks and in \RateBehBSevenGenThreeQnaNoTool{} of general-3 tasks, and a tool call does not flip it: flat-3 answers with the lie after a tool in \RateBehBSevenFlatThreeQnaToolLie{} of tasks and with the source value in \RateBehBSevenFlatThreeQnaToolGold{}.
The second model drops the claim once it has read the source; the first model states the claim either way.
The user prompt also says to prefer source code, which conflicts with the claim, and the extra tool turns show up as input tokens (\TokBEightFlatThreeQna{}k against \TokBSevenFlatThreeQna{}k for flat-3 on QnA).
Judge-only and BPD match flat-3 in the table, since they share its worker prompt.

% Auto-generated by code/03_export_paper_tables.py — do not edit by hand.
\begin{table*}[t]
\centering
\scriptsize
\caption{Poisoned replica on QnA tasks ($n{=}100$). No tool: it answered without a tool call; No tool+lie: and the answer was the planted lie; Tool+gold: it called a tool and answered with the source value; Tool+lie: it called a tool and still answered with the lie. A: \texttt{gpt-5.6-luna}; B: \texttt{gpt-6-luna}.}
\label{tab:emit}
\begin{tabular}{lcccccccc}
\toprule
Arm & No tool A & No tool+lie A & Tool+gold A & Tool+lie A & No tool B & No tool+lie B & Tool+gold B & Tool+lie B \\
\midrule
general-1 & 1.00 & 1.00 & 0.00 & 0.00 & 0.95 & 0.95 & 0.05 & 0.00 \\
general-3 maj. & 1.00 & 1.00 & 0.00 & 0.00 & 0.91 & 0.91 & 0.05 & 0.04 \\
flat-3 maj. & 0.90 & 0.90 & 0.01 & 0.09 & 0.11 & 0.11 & 0.88 & 0.01 \\
judge-only-3 & 0.88 & 0.88 & 0.02 & 0.10 & 0.11 & 0.11 & 0.88 & 0.01 \\
dag-bpd & 0.87 & 0.87 & 0.00 & 0.13 & 0.09 & 0.09 & 0.90 & 0.01 \\
\bottomrule
\end{tabular}
\end{table*}


\subsection{A hierarchy that trusts its first node}
In the hierarchical system the lie enters the scan report, so the injection point differs from the other arms and the comparison is indicative.
In the earlier 100-task run, the naive hierarchy propagates the lie in \RateHundredYonseiQnaProp{} of QnA runs and \RateHundredYonseiCodeProp{} of code runs, and it names the scan node in \RateHundredYonseiQnaDet{} and \RateHundredYonseiCodeDet{} of runs.
We do not compare this rate with the single loop's, because the single loop's rate depends on the delivery channel (Section~\ref{sec:channel}).
With the source-over-scan rule, QnA recovery is \RateYonseiImprovedYonseiQnaRec{} and code recovery \RateYonseiImprovedYonseiCodeRec{}.
The rule works because the source file and the tests give the implement node an independent check, not because the graph has more nodes.

\subsection{Sensitivity to how the claim is delivered}
\label{sec:channel}
The single loop's QnA propagation is \RateHundredGenOneQnaProp{} with the legacy user-message wording (100 tasks), \RateBFourGenOneQnaProp{} in a 20-task pilot, \RateBFiveGenOneQnaProp{} when the replica-prompt text is delivered as a user message, and \RateMainGenOneQnaProp{} when the same text is placed in the system prompt.
Code propagation in the same four settings is \RateHundredGenOneCodeProp{}, \RateBFourGenOneCodeProp{}, \RateBFiveGenOneCodeProp{}, and \RateMainGenOneCodeProp{}.
In our setting the delivery channel decides whether a poisoned single loop repeats the lie at all: a claim in the system prompt is followed, and the same claim in a user message mostly is not.
Earlier comparisons gave the single loop the claim as a user message and the other arms in the system prompt, which made the single loop look more robust than it is (QnA accuracy \RateBFiveGenOneQnaAcc{} there, \RateMainGenOneQnaAcc{} here); in that run general-3 also named the source in none of its runs, because a single loop that does not repeat the lie leaves the vote nothing to count.
We therefore report the system-prompt runs as the main result and the user-message run in the appendix.

\section{Discussion}

\paragraph{Three separate properties.}
Our results suggest treating isolation, comparison, and attribution as separate design choices.
Isolation limits how far a lie travels, and replication provides it at about three times the token cost; with the lie in the system prompt, a single loop has no isolation at all on QnA.
Replication alone does not protect bug fixing, because choosing among replicas without reading them is a gamble on which one is clean.
Comparison needs a node that sees several branches at once and reads their content.
Attribution maps a disagreement between branches back to the node that produced the outlier, and it is bounded by whether the poisoned node emits the lie at all.
A directed acyclic graph supplies the edges that make this mapping possible, but a graph with a single authoritative upstream node has no branch to compare against.
That is consistent with the hierarchical system repeating the lie when its scan node is poisoned, and with the fix that adds an independent source of evidence.

\paragraph{Where the judge matters.}
Majority voting is enough when answers are short strings and the lie is visible as one: on QnA it names the source more often than the judge does.
When correct and planted behavior are described in free text, a comparison step has to read the replicas' content, and an LLM judge does.
On these small single-module tasks a judge-free structural diff does the same job as well as the judge does, because the correct replicas write the same fix and the one that differs stands out.
We expect this to break on larger tasks in which correct replicas legitimately differ. There a judge or a behavioral comparison may matter, but we did not test that.

\paragraph{Measurement lessons.}
Two controls changed our conclusions. Placing the same text in a user message instead of the system prompt moved single-loop QnA propagation from \RateBFiveGenOneQnaProp{} to \RateMainGenOneQnaProp{}, and fixing the poisoned replica at index~0 hid a failure of the majority arms that appears with a balanced index.
Studies of internal poison should vary both, and report whether the poisoned node repeats the lie.

\paragraph{Relation to resilience studies.}
\citet{huang2024resilience} find that a hierarchical structure is the most resilient to faulty agents.
Our hierarchy is fragile when its first node is poisoned and its later nodes cannot contradict that node.
The two findings agree if resilience depends on whether downstream nodes have an independent signal, and the hierarchy's shape matters less.
We make this reading as a hypothesis; our design does not vary shape and independence separately.

\paragraph{Practical guidance.}
For a system that must both stay correct and report where a bad claim came from, our data favor redundant branches followed by a comparison that reads their content over authoritative single nodes.
On these small tasks the comparison can be a normalized-source diff: using it to choose the workspace reaches accuracy \RateSelMainFlatThreeCodeAst{} where the string vote reaches \RateSelMainFlatThreeCodeActual{} and a random replica reaches \RateSelMainFlatThreeCodeUniform{}. The judge call and the backward pass add cost without measurable benefit here.

\section{Conclusion}
A planted lie in one node of a coding-agent graph damages a three-replica system less than it damages a single loop, but on bug fixing only a comparison that reads the replicas' code keeps accuracy independent of which replica lies and names the node. On these tasks that comparison can be a judge-free diff of normalized sources: it matches the judge both at naming the node and at choosing the workspace, while drawing a replica at random does not.
Backpropagation on top of the judge added no measurable recovery, and attribution is capped by whether the poisoned replica repeats the lie, as a second model shows.
A hierarchical graph whose downstream nodes trust a poisoned scan report repeats the lie in most QnA runs unless a node has an independent check.
The suite and scripts are released so that others can test whether these patterns hold for other model families and harder tasks.

\section*{Limitations}
All results come from two models of one family, \texttt{gpt-5.6-luna} and \texttt{gpt-6-luna}, with one attempt per task, so run-to-run variance is estimated only from repeated earlier runs of the same arm, where BPD recovery on QnA varied by a few points.
The tasks are templates expanded from 20 hand-written items per part, single-module, and the lie is stated explicitly rather than hidden in plausible text; harder repository-scale tasks may behave differently.
The judge-free baseline, including its use as a promotion rule, works because the correct replicas on these tasks converge on the same normalized source; it is not evidence for larger tasks. When no replica is a unique outlier the rule still promotes the last replica. Its test-oracle variant uses tests the arms never see.
The majority arms promote the last replica when no proposal matches, which is a design choice; other tie-breaks (random, or test-based) might change their code accuracy, so their results describe this implementation.
Recovery of the judge arms by replica index was tested with about a third of the tasks per cell, which cannot rule out small position effects, and one uncorrected test on the second model is nominally significant (Section~\ref{sec:second}).
How the claim is delivered changes single-loop QnA propagation from \RateBFiveGenOneQnaProp{} to \RateMainGenOneQnaProp{}, so absolute single-loop rates describe this injection and not agents in general.
\emph{Emit} on code tasks is a lexical proxy (the lie phrase, or a pattern in the replica's source file) and can miss paraphrased compliance.
The second model states the QnA lie only under the single-loop prompt; under the proposal-worker prompt it reads the source and answers with that value (Section~\ref{sec:second}), so the drop in detection is a prompt-by-model interaction and not a property of the model alone.
The hierarchical arm receives the lie at a different node and was run only in earlier runs, and its source-over-scan rule was added after the first results, so it is exploratory.
Compute is matched only approximately: judge arms add summarizer and judge calls, and tokens are the only cost proxy.
The \texttt{document} condition did not affect any system, so we report no document-poison result.
The stronger planted claim we use is a stress test and does not describe the lies real collaborators produce.

\section*{Ethical Considerations}
The study uses synthetic tasks and no human subjects or personal data.
The injected claims are plain statements, and the paper does not describe an attack that is new or hard to reproduce.
\todo{confirm with the Responsible NLP checklist; disclose AI coding and writing assistance in the Acknowledgments as ACL policy requires}

\section*{Reproducibility}
Scripts \texttt{code/01} to \texttt{code/11} regenerate every table and macro in this paper from frozen per-run summaries (\texttt{score.json}) and the raw replica workspaces; no number in the text is typed by hand.
\todo{add an anonymous repository link for review, and a persistent archive of the raw run transcripts}

\bibliography{references}

\appendix
\section{Earlier design and full results with intervals}
\label{app:full}
The earlier run poisoned replica~0 in every task and used a 60-second task limit. In its first pass \CountBSixBpdCodeTimeout{} BPD code tasks hit the limit, with no tokens recorded, and no other arm did; we re-ran exactly those tasks once with the identical configuration and report the merged run, in which \CountBSixRBpdCodeTimeout{} tasks timed out again and count as failures (BPD code accuracy there is a lower bound).
Tables~\ref{tab:replication} and~\ref{tab:replication-b8} compare runs on identical task identifiers.
% Auto-generated by code/03_export_paper_tables.py — do not edit by hand.
\begin{table*}[t]
\centering
\scriptsize
\caption{Position-balanced run: the poisoned replica is replica 0, 1, or 2 in equal shares, with the identical system-prompt claim and a 180\,s task limit (\texttt{gpt-5.6-luna}, $n{=}100$ per cell, agent condition). Rates with bootstrap 95\% CI.}
\label{tab:budget7}
\begin{tabular}{lcccccccc}
\toprule
Arm & Set & Acc. & Prop. & Det. & Rec. & Emit. & Det$\mid$Emit & In-tok/task (k) \\
\midrule
general-1 & QnA & 0.000 [0.000, 0.000] & 1.000 [1.000, 1.000] & 0.000 [0.000, 0.000] & 0.000 [0.000, 0.000] & 1.000 & 0.000 & 0.6 \\
general-1 & Code & 0.590 [0.500, 0.690] & 0.400 [0.300, 0.500] & 0.000 [0.000, 0.000] & 0.000 [0.000, 0.000] & 0.470 & 0.000 & 2.8 \\
general-3 maj. & QnA & 1.000 [1.000, 1.000] & 0.000 [0.000, 0.000] & 1.000 [1.000, 1.000] & 1.000 [1.000, 1.000] & 1.000 & 1.000 & 3.0 \\
general-3 maj. & Code & 0.810 [0.730, 0.880] & 0.150 [0.080, 0.220] & 0.020 [0.000, 0.050] & 0.020 [0.000, 0.050] & 0.480 & 0.042 & 8.4 \\
flat-3 maj. & QnA & 1.000 [1.000, 1.000] & 0.000 [0.000, 0.000] & 0.990 [0.970, 1.000] & 0.990 [0.970, 1.000] & 0.990 & 1.000 & 3.0 \\
flat-3 maj. & Code & 0.720 [0.640, 0.800] & 0.260 [0.180, 0.350] & 0.010 [0.000, 0.030] & 0.010 [0.000, 0.030] & 0.590 & 0.017 & 8.1 \\
judge-only-3 & QnA & 1.000 [1.000, 1.000] & 0.000 [0.000, 0.000] & 0.910 [0.850, 0.960] & 0.910 [0.850, 0.960] & 0.980 & 0.929 & 3.0 \\
judge-only-3 & Code & 1.000 [1.000, 1.000] & 0.000 [0.000, 0.000] & 0.630 [0.540, 0.720] & 0.630 [0.540, 0.720] & 0.610 & 0.869 & 8.1 \\
dag-bpd & QnA & 1.000 [1.000, 1.000] & 0.000 [0.000, 0.000] & 0.940 [0.890, 0.980] & 0.940 [0.890, 0.980] & 1.000 & 0.940 & 3.1 \\
dag-bpd & Code & 0.990 [0.970, 1.000] & 0.010 [0.000, 0.030] & 0.670 [0.580, 0.760] & 0.670 [0.580, 0.760] & 0.610 & 0.836 & 8.1 \\
\bottomrule
\end{tabular}
\end{table*}

% Auto-generated by code/03_export_paper_tables.py — do not edit by hand.
\begin{table*}[t]
\centering
\scriptsize
\caption{Second model: the configuration of Table~\ref{tab:budget7} with \texttt{gpt-6-luna} ($n{=}100$ per cell, agent condition). Rates with bootstrap 95\% CI.}
\label{tab:budget8}
\begin{tabular}{lcccccccc}
\toprule
Arm & Set & Acc. & Prop. & Det. & Rec. & Emit. & Det$\mid$Emit & In-tok/task (k) \\
\midrule
general-1 & QnA & 0.050 [0.010, 0.100] & 0.950 [0.900, 0.990] & 0.000 [0.000, 0.000] & 0.000 [0.000, 0.000] & 0.950 & 0.000 & 0.7 \\
general-1 & Code & 0.440 [0.340, 0.530] & 0.460 [0.370, 0.560] & 0.000 [0.000, 0.000] & 0.000 [0.000, 0.000] & 0.480 & 0.000 & 3.2 \\
general-3 maj. & QnA & 1.000 [1.000, 1.000] & 0.000 [0.000, 0.000] & 0.950 [0.910, 0.990] & 0.950 [0.900, 0.990] & 0.950 & 1.000 & 3.1 \\
general-3 maj. & Code & 0.720 [0.630, 0.810] & 0.270 [0.180, 0.360] & 0.010 [0.000, 0.030] & 0.010 [0.000, 0.030] & 0.510 & 0.020 & 8.9 \\
flat-3 maj. & QnA & 1.000 [1.000, 1.000] & 0.000 [0.000, 0.000] & 0.120 [0.060, 0.190] & 0.120 [0.060, 0.190] & 0.120 & 1.000 & 4.5 \\
flat-3 maj. & Code & 0.700 [0.610, 0.780] & 0.290 [0.210, 0.380] & 0.010 [0.000, 0.030] & 0.010 [0.000, 0.030] & 0.580 & 0.017 & 8.1 \\
judge-only-3 & QnA & 1.000 [1.000, 1.000] & 0.000 [0.000, 0.000] & 0.110 [0.060, 0.170] & 0.110 [0.050, 0.170] & 0.120 & 0.917 & 4.5 \\
judge-only-3 & Code & 0.990 [0.970, 1.000] & 0.010 [0.000, 0.030] & 0.600 [0.500, 0.690] & 0.600 [0.500, 0.690] & 0.540 & 0.981 & 8.1 \\
dag-bpd & QnA & 1.000 [1.000, 1.000] & 0.000 [0.000, 0.000] & 0.100 [0.050, 0.160] & 0.100 [0.040, 0.160] & 0.100 & 1.000 & 4.5 \\
dag-bpd & Code & 0.980 [0.950, 1.000] & 0.020 [0.000, 0.050] & 0.610 [0.510, 0.700] & 0.610 [0.520, 0.710] & 0.530 & 0.962 & 8.2 \\
\bottomrule
\end{tabular}
\end{table*}

% Auto-generated by code/03_export_paper_tables.py — do not edit by hand.
\begin{table*}[t]
\centering
\scriptsize
\caption{Accuracy and recovery by the index of the poisoned replica (permutation test of dependence on the index; each cell holds about a third of the tasks of a suite). The single loop ignores the index, so its row shows the spread expected from task composition alone. Run \texttt{budget-8}.}
\label{tab:position-b8}
\begin{tabular}{lccccccccc}
\toprule
Arm & Set & Acc@0 & Acc@1 & Acc@2 & $p_{\mathrm{Acc}}$ & Rec@0 & Rec@1 & Rec@2 & $p_{\mathrm{Rec}}$ \\
\midrule
general-1 & QnA & 0.06 & 0.06 & 0.03 & 1.000 & 0.00 & 0.00 & 0.00 & 1.000 \\
general-1 & Code & 0.50 & 0.45 & 0.36 & 0.511 & 0.00 & 0.00 & 0.00 & 1.000 \\
general-3 maj. & QnA & 1.00 & 1.00 & 1.00 & 1.000 & 0.94 & 1.00 & 0.91 & 0.281 \\
general-3 maj. & Code & 0.94 & 0.82 & 0.39 & $<$0.001 & 0.00 & 0.03 & 0.00 & 0.661 \\
flat-3 maj. & QnA & 1.00 & 1.00 & 1.00 & 1.000 & 0.26 & 0.06 & 0.03 & 0.006 \\
flat-3 maj. & Code & 0.88 & 0.91 & 0.30 & $<$0.001 & 0.03 & 0.00 & 0.00 & 1.000 \\
judge-only-3 & QnA & 1.00 & 1.00 & 1.00 & 1.000 & 0.21 & 0.09 & 0.03 & 0.074 \\
judge-only-3 & Code & 0.97 & 1.00 & 1.00 & 1.000 & 0.65 & 0.42 & 0.73 & 0.037 \\
dag-bpd & QnA & 1.00 & 1.00 & 1.00 & 1.000 & 0.18 & 0.03 & 0.09 & 0.146 \\
dag-bpd & Code & 0.94 & 1.00 & 1.00 & 0.325 & 0.68 & 0.48 & 0.67 & 0.203 \\
\bottomrule
\end{tabular}
\end{table*}

% Auto-generated by code/03_export_paper_tables.py — do not edit by hand.
\begin{table*}[t]
\centering
\scriptsize
\caption{Which replica's source the arm promotes on code tasks (position-balanced runs, $n{=}100$). Last: final source equals replica~2's; Poisoned: final source equals only the poisoned replica's; Acc.\ is accuracy among those runs, and Acc.\ other among all remaining runs. A: \texttt{gpt-5.6-luna}; B: \texttt{gpt-6-luna}.}
\label{tab:fallback}
\begin{tabular}{lcccccccc}
\toprule
Arm & Last A & Poisoned A & Acc. poisoned A & Acc. other A & Last B & Poisoned B & Acc. poisoned B & Acc. other B \\
\midrule
general-3 maj. & 0.91 & 0.34 & 0.47 & 0.98 & 0.91 & 0.33 & 0.15 & 1.00 \\
flat-3 maj. & 0.90 & 0.40 & 0.30 & 1.00 & 0.92 & 0.37 & 0.19 & 1.00 \\
judge-only-3 & 0.48 & 0.07 & 1.00 & 1.00 & 0.56 & 0.07 & 0.86 & 1.00 \\
dag-bpd & 0.43 & 0.11 & 0.91 & 1.00 & 0.44 & 0.06 & 0.67 & 1.00 \\
\bottomrule
\end{tabular}
\end{table*}

% Auto-generated by code/03_export_paper_tables.py — do not edit by hand.
\begin{table*}[t]
\centering
\scriptsize
\caption{Same arm and task ids in the position-balanced run (A) and the earlier run in which replica~0 always carried the lie (B). a/b: tasks correct only in A / only in B; $p$: exact McNemar. The runs differ in poisoned-replica index and task limit, so differences mix both with sampling noise.}
\label{tab:replication}
\begin{tabular}{lccccccccc}
\toprule
Arm & Set & Acc A & Acc B & Acc a/b & p & Rec A & Rec B & Rec a/b & p \\
\midrule
general-1 & QnA & 0.00 & 0.00 & 0/0 & 1.00 & 0.00 & 0.00 & 0/0 & 1.00 \\
general-1 & Code & 0.59 & 0.54 & 11/6 & 0.33 & 0.00 & 0.00 & 0/0 & 1.00 \\
general-3 maj. & QnA & 1.00 & 1.00 & 0/0 & 1.00 & 1.00 & 1.00 & 0/0 & 1.00 \\
general-3 maj. & Code & 0.81 & 0.97 & 2/18 & 0.00 & 0.02 & 0.01 & 1/0 & 1.00 \\
flat-3 maj. & QnA & 1.00 & 1.00 & 0/0 & 1.00 & 0.99 & 0.99 & 1/1 & 1.00 \\
flat-3 maj. & Code & 0.72 & 0.92 & 4/24 & 0.00 & 0.01 & 0.01 & 1/1 & 1.00 \\
judge-only-3 & QnA & 1.00 & 1.00 & 0/0 & 1.00 & 0.91 & 0.85 & 14/8 & 0.29 \\
judge-only-3 & Code & 1.00 & 0.99 & 1/0 & 1.00 & 0.63 & 0.64 & 8/9 & 1.00 \\
dag-bpd & QnA & 1.00 & 1.00 & 0/0 & 1.00 & 0.94 & 0.87 & 12/5 & 0.14 \\
dag-bpd & Code & 0.99 & 0.96 & 4/1 & 0.38 & 0.67 & 0.63 & 10/6 & 0.45 \\
\bottomrule
\end{tabular}
\end{table*}

% Auto-generated by code/03_export_paper_tables.py — do not edit by hand.
\begin{table*}[t]
\centering
\scriptsize
\caption{Same arm and task ids with the second model (A, \texttt{gpt-6-luna}) and the first model (B, \texttt{gpt-5.6-luna}), both position-balanced. a/b: tasks correct only in A / only in B; $p$: exact McNemar.}
\label{tab:replication-b8}
\begin{tabular}{lccccccccc}
\toprule
Arm & Set & Acc A & Acc B & Acc a/b & p & Rec A & Rec B & Rec a/b & p \\
\midrule
general-1 & QnA & 0.05 & 0.00 & 5/0 & 0.06 & 0.00 & 0.00 & 0/0 & 1.00 \\
general-1 & Code & 0.44 & 0.59 & 12/27 & 0.02 & 0.00 & 0.00 & 0/0 & 1.00 \\
general-3 maj. & QnA & 1.00 & 1.00 & 0/0 & 1.00 & 0.95 & 1.00 & 0/5 & 0.06 \\
general-3 maj. & Code & 0.72 & 0.81 & 7/16 & 0.09 & 0.01 & 0.02 & 1/2 & 1.00 \\
flat-3 maj. & QnA & 1.00 & 1.00 & 0/0 & 1.00 & 0.12 & 0.99 & 0/87 & 0.00 \\
flat-3 maj. & Code & 0.70 & 0.72 & 8/10 & 0.81 & 0.01 & 0.01 & 1/1 & 1.00 \\
judge-only-3 & QnA & 1.00 & 1.00 & 0/0 & 1.00 & 0.11 & 0.91 & 0/80 & 0.00 \\
judge-only-3 & Code & 0.99 & 1.00 & 0/1 & 1.00 & 0.60 & 0.63 & 13/16 & 0.71 \\
dag-bpd & QnA & 1.00 & 1.00 & 0/0 & 1.00 & 0.10 & 0.94 & 1/85 & 0.00 \\
dag-bpd & Code & 0.98 & 0.99 & 0/1 & 1.00 & 0.61 & 0.67 & 9/15 & 0.31 \\
\bottomrule
\end{tabular}
\end{table*}

% Auto-generated by code/03_export_paper_tables.py — do not edit by hand.
\begin{table*}[t]
\centering
\scriptsize
\caption{Judge-free attribution on code replicas (second model \texttt{gpt-6-luna}; no LLM call). AST outlier: the one replica whose normalized source differs from two identical ones; wrong: it flags a replica other than the poisoned one; test oracle: the one replica that fails hidden tests (an upper bound the arms cannot use). Arm detection is the arm's own detection rate on the same tasks.}
\label{tab:semantic-b8}
\begin{tabular}{lccccc}
\toprule
Arm & $n$ & AST outlier & AST wrong & Test oracle & Arm detection \\
\midrule
general-3 maj. & 100 & 0.56 & 0.02 & 0.57 & 0.01 \\
flat-3 maj. & 100 & 0.64 & 0.01 & 0.62 & 0.01 \\
judge-only-3 & 100 & 0.55 & 0.01 & 0.53 & 0.60 \\
dag-bpd & 100 & 0.59 & 0.02 & 0.60 & 0.61 \\
\bottomrule
\end{tabular}
\end{table*}

% Auto-generated by code/03_export_paper_tables.py — do not edit by hand.
\begin{table*}[t]
\centering
\scriptsize
\caption{Judge-free attribution on code replicas (main run of the earlier, replica-0-only design; no LLM call). AST outlier: the one replica whose normalized source differs from two identical ones; wrong: it flags a replica other than the poisoned one; test oracle: the one replica that fails hidden tests (an upper bound the arms cannot use). Arm detection is the arm's own detection rate on the same tasks.}
\label{tab:semantic-b6r}
\begin{tabular}{lccccc}
\toprule
Arm & $n$ & AST outlier & AST wrong & Test oracle & Arm detection \\
\midrule
general-3 maj. & 100 & 0.49 & 0.07 & 0.42 & 0.01 \\
flat-3 maj. & 100 & 0.67 & 0.04 & 0.64 & 0.01 \\
judge-only-3 & 100 & 0.63 & 0.05 & 0.63 & 0.64 \\
dag-bpd & 97 & 0.66 & 0.03 & 0.63 & 0.65 \\
\bottomrule
\end{tabular}
\end{table*}

% Auto-generated by code/03_export_paper_tables.py — do not edit by hand.
\begin{table*}[t]
\centering
\scriptsize
\caption{Same run as Table~\ref{tab:budget6} after re-running the tasks that hit the 60\,s harness limit (identical configuration; see text). $n{=}100$ per cell, agent condition. Rates with bootstrap 95\% CI.}
\label{tab:budget6-retried}
\begin{tabular}{lcccccccc}
\toprule
Arm & Set & Acc. & Prop. & Det. & Rec. & Emit. & Det$\mid$Emit & In-tok/task (k) \\
\midrule
general-1 & QnA & 0.000 [0.000, 0.000] & 1.000 [1.000, 1.000] & 0.000 [0.000, 0.000] & 0.000 [0.000, 0.000] & 1.000 & 0.000 & 0.6 \\
general-1 & Code & 0.540 [0.440, 0.640] & 0.410 [0.310, 0.500] & 0.000 [0.000, 0.000] & 0.000 [0.000, 0.000] & 0.500 & 0.000 & 3.0 \\
general-3 maj. & QnA & 1.000 [1.000, 1.000] & 0.000 [0.000, 0.000] & 1.000 [1.000, 1.000] & 1.000 [1.000, 1.000] & 1.000 & 1.000 & 3.0 \\
general-3 maj. & Code & 0.970 [0.930, 1.000] & 0.020 [0.000, 0.050] & 0.010 [0.000, 0.030] & 0.010 [0.000, 0.030] & 0.490 & 0.020 & 8.2 \\
flat-3 maj. & QnA & 1.000 [1.000, 1.000] & 0.000 [0.000, 0.000] & 0.990 [0.970, 1.000] & 0.990 [0.960, 1.000] & 0.990 & 1.000 & 3.1 \\
flat-3 maj. & Code & 0.920 [0.860, 0.970] & 0.080 [0.030, 0.140] & 0.010 [0.000, 0.030] & 0.010 [0.000, 0.030] & 0.560 & 0.018 & 8.1 \\
judge-only-3 & QnA & 1.000 [1.000, 1.000] & 0.000 [0.000, 0.000] & 0.850 [0.770, 0.910] & 0.850 [0.770, 0.910] & 0.990 & 0.859 & 3.0 \\
judge-only-3 & Code & 0.990 [0.970, 1.000] & 0.010 [0.000, 0.030] & 0.640 [0.550, 0.730] & 0.640 [0.550, 0.740] & 0.580 & 0.897 & 8.1 \\
dag-bpd & QnA & 1.000 [1.000, 1.000] & 0.000 [0.000, 0.000] & 0.870 [0.790, 0.930] & 0.870 [0.800, 0.930] & 0.980 & 0.888 & 3.0 \\
dag-bpd & Code & 0.960 [0.910, 0.990] & 0.010 [0.000, 0.040] & 0.630 [0.540, 0.720] & 0.630 [0.540, 0.720] & 0.608 & 0.898 & 7.8 \\
\bottomrule
\end{tabular}
\end{table*}

% Auto-generated by code/03_export_paper_tables.py — do not edit by hand.
\begin{table*}[t]
\centering
\scriptsize
\caption{Compute-matched arms with the identical claim injected through the system prompt of every arm ($n{=}100$ per cell, agent condition). Emit.: share of runs in which the poisoned replica itself output the lie; Det$\mid$Emit: attribution rate among those runs. Rates with bootstrap 95\% CI.}
\label{tab:budget6}
\begin{tabular}{lcccccccc}
\toprule
Arm & Set & Acc. & Prop. & Det. & Rec. & Emit. & Det$\mid$Emit & In-tok/task (k) \\
\midrule
general-1 & QnA & 0.000 [0.000, 0.000] & 1.000 [1.000, 1.000] & 0.000 [0.000, 0.000] & 0.000 [0.000, 0.000] & 1.000 & 0.000 & 0.6 \\
general-1 & Code & 0.540 [0.440, 0.640] & 0.410 [0.310, 0.510] & 0.000 [0.000, 0.000] & 0.000 [0.000, 0.000] & 0.500 & 0.000 & 3.0 \\
general-3 maj. & QnA & 1.000 [1.000, 1.000] & 0.000 [0.000, 0.000] & 1.000 [1.000, 1.000] & 1.000 [1.000, 1.000] & 1.000 & 1.000 & 3.0 \\
general-3 maj. & Code & 0.970 [0.930, 1.000] & 0.020 [0.000, 0.050] & 0.010 [0.000, 0.030] & 0.010 [0.000, 0.030] & 0.490 & 0.020 & 8.2 \\
flat-3 maj. & QnA & 1.000 [1.000, 1.000] & 0.000 [0.000, 0.000] & 0.990 [0.970, 1.000] & 0.990 [0.970, 1.000] & 0.990 & 1.000 & 3.1 \\
flat-3 maj. & Code & 0.920 [0.860, 0.970] & 0.080 [0.030, 0.140] & 0.010 [0.000, 0.030] & 0.010 [0.000, 0.030] & 0.560 & 0.018 & 8.1 \\
judge-only-3 & QnA & 1.000 [1.000, 1.000] & 0.000 [0.000, 0.000] & 0.850 [0.770, 0.910] & 0.850 [0.780, 0.920] & 0.990 & 0.859 & 3.0 \\
judge-only-3 & Code & 0.990 [0.970, 1.000] & 0.010 [0.000, 0.030] & 0.640 [0.550, 0.730] & 0.640 [0.550, 0.740] & 0.580 & 0.897 & 8.1 \\
dag-bpd & QnA & 1.000 [1.000, 1.000] & 0.000 [0.000, 0.000] & 0.870 [0.800, 0.930] & 0.870 [0.800, 0.930] & 0.980 & 0.888 & 3.0 \\
dag-bpd & Code & 0.880 [0.810, 0.940] & 0.010 [0.000, 0.030] & 0.610 [0.520, 0.700] & 0.610 [0.510, 0.710] & 0.652 & 0.897 & 7.2 \\
\bottomrule
\end{tabular}
\end{table*}

% Auto-generated by code/03_export_paper_tables.py — do not edit by hand.
\begin{table*}[t]
\centering
\scriptsize
\caption{Compute-matched arms with aligned wording ($n{=}100$ per cell, agent condition). Here the claim reaches the single loop and general-3 as a user message, and the other arms in the system prompt (superseded by the main run). Rates with bootstrap 95\% CI; last column is mean input tokens per task.}
\label{tab:budget5}
\begin{tabular}{lcccccc}
\toprule
Arm & Set & Acc. & Prop. & Det. & Rec. & In-tok/task (k) \\
\midrule
general-1 & QnA & 0.990 [0.970, 1.000] & 0.010 [0.000, 0.030] & 0.000 [0.000, 0.000] & 0.000 [0.000, 0.000] & 1.3 \\
general-1 & Code & 0.840 [0.770, 0.910] & 0.110 [0.050, 0.180] & 0.000 [0.000, 0.000] & 0.000 [0.000, 0.000] & 3.0 \\
general-3 maj. & QnA & 1.000 [1.000, 1.000] & 0.000 [0.000, 0.000] & 0.000 [0.000, 0.000] & 0.000 [0.000, 0.000] & 3.6 \\
general-3 maj. & Code & 1.000 [1.000, 1.000] & 0.000 [0.000, 0.000] & 0.000 [0.000, 0.000] & 0.000 [0.000, 0.000] & 8.4 \\
flat-3 maj. & QnA & 1.000 [1.000, 1.000] & 0.000 [0.000, 0.000] & 1.000 [1.000, 1.000] & 1.000 [1.000, 1.000] & 3.0 \\
flat-3 maj. & Code & 0.910 [0.850, 0.960] & 0.090 [0.040, 0.150] & 0.000 [0.000, 0.000] & 0.000 [0.000, 0.000] & 8.1 \\
judge-only-3 & QnA & 1.000 [1.000, 1.000] & 0.000 [0.000, 0.000] & 0.840 [0.760, 0.910] & 0.840 [0.760, 0.910] & 3.0 \\
judge-only-3 & Code & 1.000 [1.000, 1.000] & 0.000 [0.000, 0.000] & 0.630 [0.540, 0.720] & 0.630 [0.530, 0.720] & 8.2 \\
dag-bpd & QnA & 1.000 [1.000, 1.000] & 0.000 [0.000, 0.000] & 0.880 [0.810, 0.940] & 0.880 [0.810, 0.940] & 3.0 \\
dag-bpd & Code & 0.990 [0.970, 1.000] & 0.010 [0.000, 0.030] & 0.630 [0.540, 0.720] & 0.630 [0.530, 0.720] & 8.1 \\
\bottomrule
\end{tabular}
\end{table*}

% Auto-generated by code/03_export_paper_tables.py — do not edit by hand.
\begin{table}[t]
\centering
\caption{Compute-matched pilot (core $n{=}20$ per cell, legacy injection wording). Rates with bootstrap 95\% CI.}
\label{tab:budget4-pilot}
\begin{tabular}{lcccccc}
\toprule
Arm & Set & Acc. & Prop. & Det. & Rec. & In-tok/task (k) \\
\midrule
general-1 & QnA & 0.900 [0.750, 1.000] & 0.100 [0.000, 0.250] & 0.000 [0.000, 0.000] & 0.000 [0.000, 0.000] & 1.2 \\
general-1 & Code & 0.700 [0.500, 0.900] & 0.250 [0.050, 0.450] & 0.000 [0.000, 0.000] & 0.000 [0.000, 0.000] & 3.1 \\
general-3 maj. & QnA & 1.000 [1.000, 1.000] & 0.000 [0.000, 0.000] & 0.150 [0.000, 0.300] & 0.150 [0.000, 0.300] & 3.6 \\
general-3 maj. & Code & 0.950 [0.850, 1.000] & 0.000 [0.000, 0.000] & 0.000 [0.000, 0.000] & 0.000 [0.000, 0.000] & 9.1 \\
flat-3 maj. & QnA & 1.000 [1.000, 1.000] & 0.000 [0.000, 0.000] & 1.000 [1.000, 1.000] & 1.000 [1.000, 1.000] & 3.0 \\
flat-3 maj. & Code & 0.950 [0.850, 1.000] & 0.050 [0.000, 0.150] & 0.000 [0.000, 0.000] & 0.000 [0.000, 0.000] & 8.7 \\
dag-bpd & QnA & 1.000 [1.000, 1.000] & 0.000 [0.000, 0.000] & 0.950 [0.850, 1.000] & 0.950 [0.850, 1.000] & 3.0 \\
dag-bpd & Code & 1.000 [1.000, 1.000] & 0.000 [0.000, 0.000] & 0.400 [0.200, 0.600] & 0.400 [0.200, 0.600] & 8.7 \\
\bottomrule
\end{tabular}
\end{table}

% Auto-generated by code/03_export_paper_tables.py — do not edit by hand.
\begin{table}[t]
\centering
\caption{Main agent-condition rates on poison-100 ($n{=}100$ per cell). Cells show rate with bootstrap 95\% CI (seed 20260909). Not an Index score.}
\label{tab:main-agent}
\begin{tabular}{lccccc}
\toprule
System & Set & Acc. & Prop. & Det. & Rec. \\
\midrule
dag-bpd & QnA & 1.000 [1.000, 1.000] & 0.000 [0.000, 0.000] & 0.850 [0.780, 0.910] & 0.850 [0.780, 0.920] \\
dag-bpd & Code & 0.970 [0.930, 1.000] & 0.020 [0.000, 0.050] & 0.630 [0.540, 0.720] & 0.630 [0.540, 0.720] \\
dag-yonsei & QnA & 0.290 [0.200, 0.390] & 0.710 [0.620, 0.790] & 0.290 [0.200, 0.380] & 0.290 [0.200, 0.380] \\
dag-yonsei & Code & 1.000 [1.000, 1.000] & 0.000 [0.000, 0.000] & 0.230 [0.150, 0.310] & 0.230 [0.140, 0.310] \\
general-agent-system & QnA & 0.830 [0.750, 0.890] & 0.170 [0.100, 0.250] & 0.000 [0.000, 0.000] & 0.000 [0.000, 0.000] \\
general-agent-system & Code & 0.880 [0.810, 0.940] & 0.100 [0.050, 0.160] & 0.000 [0.000, 0.000] & 0.000 [0.000, 0.000] \\
\bottomrule
\end{tabular}
\end{table}

% Auto-generated by code/03_export_paper_tables.py — do not edit by hand.
\begin{table*}[t]
\centering
\scriptsize
\caption{Exploratory yonsei A/B: naive scan trust vs source-over-scan rule (same model and strong agent injection).}
\label{tab:yonsei-ab}
\begin{tabular}{lccccc}
\toprule
Version & Set & Acc. & Prop. & Det. & Rec. \\
\midrule
naive & QnA & 0.290 [0.200, 0.390] & 0.710 [0.620, 0.790] & 0.290 [0.200, 0.380] & 0.290 [0.200, 0.380] \\
naive & Code & 1.000 [1.000, 1.000] & 0.000 [0.000, 0.000] & 0.230 [0.150, 0.310] & 0.230 [0.140, 0.310] \\
source-first & QnA & 1.000 [1.000, 1.000] & 0.000 [0.000, 0.000] & 1.000 [1.000, 1.000] & 1.000 [1.000, 1.000] \\
source-first & Code & 1.000 [1.000, 1.000] & 0.000 [0.000, 0.000] & 1.000 [1.000, 1.000] & 1.000 [1.000, 1.000] \\
\bottomrule
\end{tabular}
\end{table*}

% Auto-generated by code/03_export_paper_tables.py — do not edit by hand.
\begin{table}[t]
\centering
\caption{Condition ablation on the 100+100 suite (none / document / agent). Rates with bootstrap 95\% CI.}
\label{tab:condition-ablation}
\begin{tabular}{lcccccc}
\toprule
System & Cond. & Set & Acc. & Prop. & Det. & Rec. \\
\midrule
dag-bpd & none & QnA & 1.000 [1.000, 1.000] & 0.000 [0.000, 0.000] & 0.000 [0.000, 0.000] & 0.000 [0.000, 0.000] \\
dag-bpd & none & Code & 1.000 [1.000, 1.000] & 0.000 [0.000, 0.000] & 0.000 [0.000, 0.000] & 0.000 [0.000, 0.000] \\
dag-bpd & document & QnA & 1.000 [1.000, 1.000] & 0.000 [0.000, 0.000] & 0.000 [0.000, 0.000] & 0.000 [0.000, 0.000] \\
dag-bpd & document & Code & 0.990 [0.960, 1.000] & 0.000 [0.000, 0.000] & 0.000 [0.000, 0.000] & 0.000 [0.000, 0.000] \\
dag-bpd & agent & QnA & 1.000 [1.000, 1.000] & 0.000 [0.000, 0.000] & 0.870 [0.800, 0.930] & 0.870 [0.800, 0.940] \\
dag-bpd & agent & Code & 0.960 [0.920, 0.990] & 0.020 [0.000, 0.050] & 0.640 [0.540, 0.730] & 0.640 [0.540, 0.730] \\
dag-yonsei & none & QnA & 0.990 [0.970, 1.000] & 0.000 [0.000, 0.000] & 0.000 [0.000, 0.000] & 0.000 [0.000, 0.000] \\
dag-yonsei & none & Code & 1.000 [1.000, 1.000] & 0.000 [0.000, 0.000] & 0.000 [0.000, 0.000] & 0.000 [0.000, 0.000] \\
dag-yonsei & document & QnA & 0.980 [0.950, 1.000] & 0.000 [0.000, 0.000] & 0.000 [0.000, 0.000] & 0.000 [0.000, 0.000] \\
dag-yonsei & document & Code & 1.000 [1.000, 1.000] & 0.000 [0.000, 0.000] & 0.000 [0.000, 0.000] & 0.000 [0.000, 0.000] \\
dag-yonsei & agent & QnA & 1.000 [1.000, 1.000] & 0.000 [0.000, 0.000] & 1.000 [1.000, 1.000] & 1.000 [1.000, 1.000] \\
dag-yonsei & agent & Code & 1.000 [1.000, 1.000] & 0.000 [0.000, 0.000] & 1.000 [1.000, 1.000] & 1.000 [1.000, 1.000] \\
general-agent-system & none & QnA & 1.000 [1.000, 1.000] & 0.000 [0.000, 0.000] & 0.000 [0.000, 0.000] & 0.000 [0.000, 0.000] \\
general-agent-system & none & Code & 0.980 [0.950, 1.000] & 0.000 [0.000, 0.000] & 0.000 [0.000, 0.000] & 0.000 [0.000, 0.000] \\
general-agent-system & document & QnA & 1.000 [1.000, 1.000] & 0.000 [0.000, 0.000] & 0.000 [0.000, 0.000] & 0.000 [0.000, 0.000] \\
general-agent-system & document & Code & 0.990 [0.970, 1.000] & 0.000 [0.000, 0.000] & 0.000 [0.000, 0.000] & 0.000 [0.000, 0.000] \\
general-agent-system & agent & QnA & 0.800 [0.720, 0.880] & 0.200 [0.130, 0.280] & 0.000 [0.000, 0.000] & 0.000 [0.000, 0.000] \\
general-agent-system & agent & Code & 0.790 [0.710, 0.870] & 0.190 [0.120, 0.260] & 0.000 [0.000, 0.000] & 0.000 [0.000, 0.000] \\
\bottomrule
\end{tabular}
\end{table}


\end{document}
